% Lectura y comentario : la familia y sus problemas (viudez, herencia) ; juego de rol — Espagnol Tle A — Cours signé OnBuch+
% Programme officiel MINESEC. Compiler avec : tectonic E05-familia-problemas-viudez-herencia.tex
\newcommand{\DOCMATIERE}{Espagnol}
\newcommand{\DOCNIVEAU}{Tle A}
\newcommand{\DOCMODULE}{Módulo 1 : Vie familiale et intégration sociale}
\newcommand{\DOCLECON}{E05}
\newcommand{\DOCTITRE}{Lectura y comentario : la familia y sus problemas (viudez, herencia) ; juego de rol}
% OnBuch+ — préambule « cours signés » ENRICHI (compilé avec tectonic / XeLaTeX)
% Système visuel « Pop » d'OnBuch+ : fond crème (jamais blanc), bordures encre
% épaisses, ombres « dures » décalées, titres Archivo Black en capitales.
% Version MATHÉMATIQUES : graphes (pgfplots/tikz), boîtes pédagogiques
% (définition, propriété, méthode, attention, le savais-tu, exercice résolu,
% exercice, TP, bilan), page de garde et signature OnBuch+ ancrée en bas.
%
% Macros attendues dans main.tex AVANT \input{preamble} :
%   \DOCMATIERE  \DOCNIVEAU  \DOCMODULE  \DOCLECON (numéro)  \DOCTITRE
\documentclass[11pt]{article}

\usepackage{fontspec}
\usepackage[spanish,french]{babel} % français = langue principale ; l'espagnol via \esp{..} / espagnol
\usepackage[a4paper,margin=2cm,top=3.4cm,bottom=2.8cm,headheight=1.4cm,footskip=1.3cm]{geometry}
\usepackage{amsmath,amssymb,amsfonts}
\usepackage{mathtools} % \xmapsto, \coloneqq... souvent utilisés par les modèles
\usepackage{cancel} % simplifications barrées (bilans de forces)
% \abs{z} (module, valeur absolue) et \norm{u} (norme), utilisés par les modèles
\providecommand{\abs}{}\let\abs\relax\DeclarePairedDelimiter{\abs}{\lvert}{\rvert}
\providecommand{\norm}{}\let\norm\relax\DeclarePairedDelimiter{\norm}{\lVert}{\rVert}
\usepackage[table]{xcolor} % [table] : fournit aussi \rowcolors (alternance de lignes)
\usepackage{pagecolor}
\usepackage{tikz}
\usetikzlibrary{arrows.meta,positioning,calc,shapes.geometric,shapes.misc,shapes.arrows,decorations.pathmorphing,decorations.pathreplacing,decorations.markings,patterns,shadows,mindmap,trees,fit,backgrounds,matrix,intersections,angles,quotes}
\usepackage{pgfplots}
\usepackage[version=4]{mhchem} % équations nucléaires \ce{^{235}_{92}U}
\usepackage[european,siunitx]{circuitikz} % schémas électriques
\pgfplotsset{compat=1.18}
\usepgfplotslibrary{fillbetween,groupplots} % aires entre courbes (calcul intégral)
\usepackage{enumitem}
\usepackage{fancyhdr}
% [most] charge toutes les bibliothèques tcolorbox (listings, minted, raster,
% documentation...) et gonfle inutilement la mémoire TeX sur un document long
% (~300 boîtes) : on ne charge que ce qui est réellement utilisé.
\usepackage[skins,breakable]{tcolorbox}
\usepackage{graphicx}
\usepackage{colortbl}
\usepackage{booktabs}
\usepackage{tabularx}
\usepackage{diagbox} % case d'en-tête diagonale des tableaux à double entrée
% Colonnes extensibles centrée / gauche / droite (C, L, R), souvent utilisées par les modèles
\newcolumntype{C}{>{\centering\arraybackslash}X}
\newcolumntype{L}{>{\raggedright\arraybackslash}X}
\newcolumntype{R}{>{\raggedleft\arraybackslash}X}
\usepackage{array}
\usepackage{multicol}
\usepackage{multirow}
\usepackage{siunitx}
% parse-numbers=false : accepte aussi \SI{2,5 \times 10^{-3}}{...}
\sisetup{locale=FR,output-decimal-marker={,},group-minimum-digits=4,parse-numbers=false}
\DeclareSIUnit\ohm{\ensuremath{\symup{\Omega}}} % U+2126 absent de la police
\DeclareSIUnit\division{div} % réglages d'oscilloscope (ms/div, V/div)
\DeclareSIUnit\tr{tr} % tours (vitesse de rotation, ex. \SI{1500}{\tr\per\minute})
\DeclareSIUnit\molar{M} % concentration molaire (ex. \SI{3}{\molar}, \SI{1}{\micro\molar})
\DeclareSIUnit\an{an} % année (le modèle confond parfois avec \year, un compteur TeX) — ex. \SI{5}{\kilo\meter\per\an}
\DeclareSIUnit\FCFA{FCFA} % franc CFA (ex. \SI{500}{\FCFA\per\kilo\gram})
\DeclareSIUnit\degreeBrix{\textdegree Brix} % teneur en sucre d'un sirop/jus (réfractométrie)
\DeclareSIUnit\month{mois} % le modèle utilise parfois \month, un compteur TeX, comme unité de durée
\usepackage{qrcode}
% \degree : commande fréquemment utilisée par les modèles pour un angle ou une
% température en dehors de \SI{}{} (non fournie par les paquets chargés).
\providecommand{\degree}{^{\circ}}
% \qed (fin de démonstration), utilisé par les modèles sans amsthm
\providecommand{\qed}{\hfill$\square$}
% \barre : double barre des valeurs interdites dans un tableau de variations (tkz-tab)
\providecommand{\barre}{\ensuremath{\|}}
% environnement proof (amsthm non chargé) utilisé par les modèles
\ifdefined\proof\else\newenvironment{proof}[1][Démonstration]{\par\noindent\textit{#1.}\ }{\qed\par}\fi
\usepackage{lastpage}
\usepackage{hyperref}
\hypersetup{hidelinks}

% ============================================================
% Palette « Pop » OnBuch+ — mêmes hex que la classe P (plus_ui.dart)
% ============================================================
\definecolor{poporange}{HTML}{F59321}
\definecolor{popdark}{HTML}{E07A0C}
\definecolor{popcream}{HTML}{FFF6E8}
\definecolor{popink}{HTML}{1C1714}
\definecolor{poppurple}{HTML}{7A5AE0}
\definecolor{popgreen}{HTML}{2E9E6A}
\definecolor{poppink}{HTML}{E0457B}
\definecolor{popblue}{HTML}{2F7DE1}
\definecolor{popgold}{HTML}{C98A12}
\definecolor{popmuted}{HTML}{8A7F76}
\definecolor{popline}{HTML}{EADFD2}
\definecolor{popgray}{HTML}{8A7F76}
\colorlet{popgrey}{popgray}
% Alias tolérants (le modèle nomme parfois les couleurs différemment)
\colorlet{popwhite}{white}
\colorlet{popblack}{popink}
\colorlet{popred}{poppink}
\colorlet{popyellow}{popgold}
% Teintes claires (fonds de tableaux/encadrés) — le modèle nomme parfois
% les couleurs avec un suffixe L (ex. popblueL) sans les avoir définies.
\colorlet{poporangeL}{poporange!20}
\colorlet{popdarkL}{popdark!20}
\colorlet{poppurpleL}{poppurple!20}
\colorlet{popgreenL}{popgreen!20}
\colorlet{poppinkL}{poppink!20}
\colorlet{popblueL}{popblue!20}
\colorlet{popgoldL}{popgold!20}
\colorlet{popredL}{popred!20}
\colorlet{popgrayL}{popgray!20}
% Teintes claires pour fonds de tableaux / zones
\colorlet{popblueL}{popblue!10!white}
\colorlet{poporangeL}{poporange!14!white}
\colorlet{popgreenL}{popgreen!12!white}
\colorlet{poppurpleL}{poppurple!12!white}
\colorlet{poppinkL}{poppink!12!white}
\colorlet{popgoldL}{popgold!14!white}

\pagecolor{popcream}

% ============================================================
% Typographie
% ============================================================
\defaultfontfeatures{Ligatures=TeX}
\setmainfont{PlusJakartaSans-Medium.ttf}[Path=../../../fonts/,BoldFont=PlusJakartaSans-Bold.ttf]
% Symboles, API et emojis absents de la police principale : repli sur DejaVu Sans (caractères actifs)
\newfontface\symfont{DejaVuSans.ttf}[Path=../../../fonts/]
\makeatletter
\newcommand\symactive[1]{\catcode#1=13 \begingroup\lccode`\~=#1 \lowercase{\endgroup\def~}{{\symfont\char#1}}}
\newcommand\symblank[1]{\catcode#1=13 \begingroup\lccode`\~=#1 \lowercase{\endgroup\def~}{}}
\newcommand\symhyph[1]{\catcode#1=13 \begingroup\lccode`\~=#1 \lowercase{\endgroup\def~}{\mbox{-}}}
\newcount\symc
\newcommand\symrange[2]{\symc=#1 \@whilenum\symc<#2 \do{\expandafter\symactive\expandafter{\the\symc}\advance\symc by 1 }}
\newcommand\symblankrange[2]{\symc=#1 \@whilenum\symc<#2 \do{\expandafter\symblank\expandafter{\the\symc}\advance\symc by 1 }}
\makeatother
\symrange{"0250}{"02FF}\symactive{"2713}\symactive{"2714}\symactive{"2717}\symactive{"2718}\symactive{"2610}\symactive{"2611}\symactive{"2612}
\symhyph{"2011}\symhyph{"2E1A}\symblankrange{"1F300}{"1FAFF}\symblank{"FE0F}
% Maths en sans-serif (Fira Math) avec chiffres et lettres droites de Plus
% Jakarta, pour que les formules se fondent dans le texte.
\usepackage[math-style=ISO,bold-style=ISO]{unicode-math}
\setmathfont{FiraMath-Regular.otf}[Path=../../../fonts/]
\setmathfont{PlusJakartaSans-Medium.ttf}[Path=../../../fonts/,range={up/num,up/{Latin,latin},\mathup}]
\usepackage{icomma} % virgule décimale sans espace en mode math
% Caractères mathématiques Unicode absents des polices de texte (Plus Jakarta,
% Archivo Black) : rendus par leur équivalent mathématique, en texte comme en
% formule — sinon ils s'affichent en carré vide (ex. « Arithmétique dans ℤ »).
\usepackage{newunicodechar}
\newunicodechar{ℝ}{\ensuremath{\mathbb{R}}}
\newunicodechar{ℤ}{\ensuremath{\mathbb{Z}}}
\newunicodechar{ℂ}{\ensuremath{\mathbb{C}}}
\newunicodechar{ℕ}{\ensuremath{\mathbb{N}}}
\newunicodechar{ℚ}{\ensuremath{\mathbb{Q}}}
\newunicodechar{𝔻}{\ensuremath{\mathbb{D}}}
\newunicodechar{α}{\ensuremath{\alpha}}
\newunicodechar{β}{\ensuremath{\beta}}
\newunicodechar{γ}{\ensuremath{\gamma}}
\newunicodechar{δ}{\ensuremath{\delta}}
\newunicodechar{ε}{\ensuremath{\varepsilon}}
\newunicodechar{θ}{\ensuremath{\theta}}
\newunicodechar{λ}{\ensuremath{\lambda}}
\newunicodechar{μ}{\ensuremath{\mu}}
\newunicodechar{σ}{\ensuremath{\sigma}}
\newunicodechar{τ}{\ensuremath{\tau}}
\newunicodechar{φ}{\ensuremath{\varphi}}
\newunicodechar{ψ}{\ensuremath{\psi}}
\newunicodechar{ω}{\ensuremath{\omega}}
\newunicodechar{Σ}{\ensuremath{\Sigma}}
% majuscules produites par \MakeUppercase dans les titres (α-aminés -> Α)
\newunicodechar{Α}{\ensuremath{\alpha}}
\newunicodechar{Β}{\ensuremath{\beta}}
\newunicodechar{Γ}{\ensuremath{\Gamma}}
\newunicodechar{Δ}{\ensuremath{\Delta}}
\newunicodechar{Ε}{\ensuremath{\varepsilon}}
\newunicodechar{Θ}{\ensuremath{\Theta}}
\newunicodechar{Λ}{\ensuremath{\Lambda}}
\newunicodechar{Μ}{\ensuremath{\mu}}
\newunicodechar{Τ}{\ensuremath{\tau}}
\newunicodechar{Φ}{\ensuremath{\Phi}}
\newunicodechar{Ψ}{\ensuremath{\Psi}}
\newunicodechar{Ω}{\ensuremath{\Omega}}
\newunicodechar{↦}{\ensuremath{\mapsto}}
\newunicodechar{∈}{\ensuremath{\in}}
\newunicodechar{∉}{\ensuremath{\notin}}
\newunicodechar{∧}{\ensuremath{\wedge}}
\newunicodechar{⇔}{\ensuremath{\Leftrightarrow}}
\newunicodechar{⟺}{\ensuremath{\Longleftrightarrow}}
\newunicodechar{⇒}{\ensuremath{\Rightarrow}}
\newunicodechar{⟹}{\ensuremath{\Longrightarrow}}
\newunicodechar{∘}{\ensuremath{\circ}}
\newunicodechar{∪}{\ensuremath{\cup}}
\newunicodechar{∩}{\ensuremath{\cap}}
\newunicodechar{≡}{\ensuremath{\equiv}}
\newunicodechar{⊂}{\ensuremath{\subset}}
\newunicodechar{⊥}{\ensuremath{\perp}}
\newunicodechar{∃}{\ensuremath{\exists}}
\newunicodechar{∀}{\ensuremath{\forall}}
\newunicodechar{∛}{\ensuremath{\sqrt[3]{\,}}}
\newunicodechar{∜}{\ensuremath{\sqrt[4]{\,}}}
\newunicodechar{⃗}{\ensuremath{{}^{\to}}}
\newunicodechar{ⁿ}{\ifmmode^{n}\else\textsuperscript{\ensuremath{n}}\fi}
\newunicodechar{ˣ}{\ifmmode^{x}\else\textsuperscript{\ensuremath{x}}\fi}
\newunicodechar{ᵇ}{\ifmmode^{b}\else\textsuperscript{\ensuremath{b}}\fi}
\newunicodechar{ʳ}{\ifmmode^{r}\else\textsuperscript{\ensuremath{r}}\fi}
\newunicodechar{ᵘ}{\ifmmode^{u}\else\textsuperscript{\ensuremath{u}}\fi}
\newunicodechar{ʸ}{\ifmmode^{y}\else\textsuperscript{\ensuremath{y}}\fi}
\newunicodechar{ᵃ}{\ifmmode^{a}\else\textsuperscript{\ensuremath{a}}\fi}
\newunicodechar{ᵉ}{\ifmmode^{e}\else\textsuperscript{\ensuremath{e}}\fi}
\newunicodechar{ᵏ}{\ifmmode^{k}\else\textsuperscript{\ensuremath{k}}\fi}
\newunicodechar{ᵖ}{\ifmmode^{p}\else\textsuperscript{\ensuremath{p}}\fi}
\newunicodechar{ᴺ}{\ifmmode^{N}\else\textsuperscript{\ensuremath{N}}\fi}
\newunicodechar{ᶜ}{\ifmmode^{c}\else\textsuperscript{\ensuremath{c}}\fi}
\newunicodechar{ʰ}{\ifmmode^{h}\else\textsuperscript{\ensuremath{h}}\fi}
\newunicodechar{ᵠ}{\ifmmode^{q}\else\textsuperscript{\ensuremath{q}}\fi}
\newunicodechar{ₙ}{\ifmmode_{n}\else\textsubscript{\ensuremath{n}}\fi}
\newunicodechar{ₐ}{\ifmmode_{a}\else\textsubscript{\ensuremath{a}}\fi}
\newunicodechar{ᵢ}{\ifmmode_{i}\else\textsubscript{\ensuremath{i}}\fi}
\newunicodechar{ᵣ}{\ifmmode_{r}\else\textsubscript{\ensuremath{r}}\fi}
\newunicodechar{ₖ}{\ifmmode_{k}\else\textsubscript{\ensuremath{k}}\fi}
\newunicodechar{ᵦ}{\ifmmode_{\beta}\else\textsubscript{\ensuremath{\beta}}\fi}
\newunicodechar{ₓ}{\ifmmode_{x}\else\textsubscript{\ensuremath{x}}\fi}
\newfontfamily\popdisplay{ArchivoBlack-Regular.ttf}[Path=../../../fonts/]
\newfontfamily\poplogo{SpaceGrotesk-Bold.ttf}[Path=../../../fonts/]
\newfontfamily\popsemi{PlusJakartaSans-SemiBold.ttf}[Path=../../../fonts/]
\newfontfamily\popxbold{PlusJakartaSans-ExtraBold.ttf}[Path=../../../fonts/]
\newfontfamily\popxboldls{PlusJakartaSans-ExtraBold.ttf}[Path=../../../fonts/,LetterSpace=3.0]

\setlist[enumerate]{leftmargin=1.7em,itemsep=5pt,topsep=4pt}
\setlist[itemize]{leftmargin=1.7em,itemsep=4pt,topsep=4pt,label={\color{poporange}\textbullet}}
\setlength{\parindent}{0pt}
\setlength{\parskip}{5pt}
\renewcommand{\arraystretch}{1.25}

% pgfplots : style Pop par défaut
\pgfplotsset{
  every axis/.append style={axis line style={popink,line width=1pt},tick style={popink},
    grid style={popline,line width=0.5pt},label style={font=\small},tick label style={font=\footnotesize}},
  popcurve/.style={poporange,line width=1.8pt,no markers,smooth},
}
% Même style disponible dans un \draw TikZ simple (hors pgfplots)
\tikzset{popcurve/.style={poporange,line width=1.8pt}}
\tikzset{popgrid/.style={popline,very thin}}
\colorlet{popcurve}{poporange} % le nom du style est parfois utilisé comme couleur

% ============================================================
% Pill / squiggle
% ============================================================
\newcommand{\poppill}[3]{%
  \tikz[baseline=-0.55ex]\node[draw=popink,line width=1.2pt,rounded corners=2.6pt,%
    fill=#2,inner xsep=6.5pt,inner ysep=3pt]{%
    \popxboldls\fontsize{7}{7}\selectfont\color{#3}\MakeUppercase{#1}};%
}
\newcommand{\popsquiggle}[1]{%
  \tikz[baseline=-0.4ex]{
    \draw[#1,line width=2.6pt,line cap=round]
      plot [smooth] coordinates {(0,0.09) (0.34,0.01) (0.68,0.21) (1.02,0.01) (1.36,0.10)};
  }%
}
% Mot-clé surligné (marqueur orange sous le texte)
\newcommand{\cle}[1]{{\popxbold\color{popdark}#1}}

% ============================================================
% En-tête / pied
% ============================================================
\pagestyle{fancy}
\fancyhf{}
\renewcommand{\headrulewidth}{0pt}
\renewcommand{\footrulewidth}{0pt}
\lhead{\raisebox{-0.15ex}{\poplogo\fontsize{13}{13}\selectfont\color{popink}OnBuch\color{poporange}+}}
\rhead{\poppill{\DOCMATIERE\ \textbullet\ \DOCNIVEAU\ \textbullet\ Leçon \DOCLECON}{white}{popink}}
\lfoot{\popxbold\fontsize{7.6}{7.6}\selectfont\color{popmuted}\MakeUppercase{Cours signé OnBuch+}\ \textbullet\ {\popsemi\DOCTITRE}}
\rfoot{\tikz[baseline=-0.6ex]\node[rounded corners=8.5pt,fill=popink,minimum height=17pt,minimum width=17pt,inner xsep=5pt,inner sep=0]{\popxbold\fontsize{8}{8}\selectfont\color{popcream}\thepage\,/\,\pageref*{LastPage}};}

% ============================================================
% Titres
% ============================================================
\newcommand{\chaptitre}[2]{%
  \begin{tcolorbox}[enhanced,colback=poporange,colframe=popink,boxrule=2.6pt,arc=11pt,
      drop shadow=popink,left=15pt,right=15pt,top=12pt,bottom=11pt,before skip=3pt,after skip=18pt]
    \poppill{#2}{popink}{popcream}\par\vspace{9pt}
    {\raggedright\hyphenpenalty=10000\exhyphenpenalty=10000\popdisplay\fontsize{22}{24}\selectfont\color{popcream}\MakeUppercase{#1}\par}
  \end{tcolorbox}
}
% Section numérotée : pastille orange avec le numéro + titre Archivo
\newcounter{popsec}
\newcommand{\coursec}[1]{%
  \stepcounter{popsec}\setcounter{popsub}{0}%
  \par\vspace{16pt}\noindent
  \begin{minipage}{\linewidth}
  \tikz[baseline=-0.7ex]\node[draw=popink,line width=1.6pt,fill=poporange,rounded corners=4pt,minimum size=22pt,inner sep=0]{\popdisplay\fontsize{12}{12}\selectfont\color{popcream}\thepopsec};%
  \hspace{8pt}{\popdisplay\fontsize{15}{17}\selectfont\color{popink}\MakeUppercase{#1}}\par
  \vspace{4pt}\noindent{\color{poporange}\rule{36pt}{3.6pt}}
  \end{minipage}\par\nopagebreak\vspace{8pt}
}
\newcounter{popsub}
\newcommand{\courssub}[1]{%
  \stepcounter{popsub}%
  \par\vspace{9pt}\noindent{\popxbold\fontsize{11.5}{13}\selectfont\color{popink}{\color{poporange}\thepopsec.\thepopsub}\hspace{6pt}#1}\par\nopagebreak\vspace{4pt}
}

% ============================================================
% Boîtes pédagogiques — même recette : fond blanc, bordure épaisse colorée,
% ombre dure encre, badge plein. Titre optionnel [#1].
% ============================================================
\tcbset{popbase/.style={breakable,enhanced,colback=white,boxrule=2.3pt,arc=9pt,
  shadow={3pt}{-3pt}{0pt}{popink},left=14pt,right=14pt,top=11pt,bottom=11pt,before skip=11pt,after skip=15pt,
  fontupper=\popsemi\fontsize{10.6}{14.5}\selectfont\color{popink},
  fontlower=\popsemi\fontsize{10.6}{14.5}\selectfont\color{popink}}}
% Coche dessinée (le glyphe ✓ n'existe pas dans les polices du document)
\newcommand{\popcheck}[1]{\tikz[baseline=-0.5ex]\draw[#1,line width=1.6pt,line cap=round,line join=round](-0.13,0.0)--(-0.04,-0.09)--(0.13,0.1);}
\providecommand{\coche}{\popcheck{popgreen}}
\providecommand{\resultat}[1]{\par\smallskip\noindent\fcolorbox{popgreen}{popgreenL}{\parbox{\dimexpr\linewidth-2\fboxsep-2\fboxrule\relax}{\textbf{Résultat.} #1}}\par\smallskip}
\newcommand{\popboxhead}[3]{\poppill{#1}{#2}{white}\ifx\relax#3\relax\else\hspace{6pt}{\popxbold\fontsize{10.6}{12}\selectfont\color{popink}#3}\fi\par\vspace{7pt}}

\newtcolorbox{aretenir}[1][]{popbase,colframe=popgreen,before upper={\popboxhead{À retenir}{popgreen}{#1}}}
% Texte en espagnol (document de lecture, dialogue, poème, extrait) : cadre
% d'étude. Un titre optionnel (source, auteur) s'affiche dans l'en-tête.
\newtcolorbox{textebox}[1][]{popbase,colframe=popblue,before upper={\popboxhead{Texto}{popblue}{#1}\begin{otherlanguage}{spanish}},after upper={\end{otherlanguage}}}
% Marques ✓ / ✗ (forme correcte / fautive) souvent utilisées par les modèles
\providecommand{\cross}{\textcolor{poppink}{\ensuremath{\times}}}
\providecommand{\xmark}{\cross}\providecommand{\crossmark}{\cross}\providecommand{\cmark}{\ensuremath{\checkmark}}
% Ligne de réponse / justification à compléter (inventée par les modèles)
\providecommand{\justification}[1]{\par\noindent\textit{Justification~:} \dotfill\par}
% \textipa (package tipa non chargé) : la police ne couvre pas l'API -> simple passage
\providecommand{\textipa}[1]{#1}
% Soulignement (ulem non chargé) : \uline, \ul
\providecommand{\uline}[1]{\underline{#1}}\providecommand{\ul}[1]{\underline{#1}}
% \glossaire{\item ...} inventée par les modèles : liste de glossaire
\providecommand{\glossaire}[1]{\begin{itemize}#1\end{itemize}}
% \alert (beamer) inventée par les modèles : texte mis en évidence
\providecommand{\alert}[1]{\textbf{\textcolor{poppink}{#1}}}
% variantes ASCII d'ordinaux écrites en macro
\providecommand{\iere}{\textsuperscript{re}}\providecommand{\ieme}{\textsuperscript{e}}\providecommand{\ere}{\textsuperscript{re}}\providecommand{\eme}{\textsuperscript{e}}\providecommand{\er}{\textsuperscript{er}}\providecommand{\ie}{\textsuperscript{e}}
% Ordinaux français écrits en macro par les modèles : 1\ière, 3\ième
\newcommand{\ière}{\textsuperscript{re}}\newcommand{\ième}{\textsuperscript{e}}
% \exemple (inventée par les modèles) : étiquette « Exemple : »
\providecommand{\exemple}{\textbf{Exemple~:}\ }
% Mot ou phrase en espagnol à l'intérieur d'un paragraphe français
\newcommand{\esp}[1]{\ifmmode\text{\textit{#1}}\else\begingroup\selectlanguage{spanish}\itshape #1\endgroup\fi}
\newtcolorbox{exemplebox}[1][]{popbase,colframe=poporange,before upper={\popboxhead{Exemple}{poporange}{#1}}}
\newtcolorbox{definition}[1][]{popbase,colframe=popblue,before upper={\popboxhead{Définition}{popblue}{#1}}}
\newtcolorbox{propriete}[1][]{popbase,colframe=poppurple,before upper={\popboxhead{Propriété}{poppurple}{#1}}}
\newtcolorbox{methode}[1][]{popbase,colframe=popgold,before upper={\popboxhead{Méthode}{popgold}{#1}}}
\newtcolorbox{attention}[1][]{popbase,colframe=poppink,before upper={\popboxhead{Attention}{poppink}{#1}}}
\newtcolorbox{experience}[1][]{popbase,colframe=popblue,colback=popblueL,before upper={\popboxhead{Expérience}{popblue}{#1}}}
% « Le savais-tu ? » : carte violette pleine (contexte, culture, Cameroun)
\newtcolorbox{savaistu}[1][]{popbase,colback=poppurple,colframe=popink,
  fontupper=\popsemi\fontsize{10.4}{14.2}\selectfont\color{white},
  before upper={\poppill{Le savais-tu\,?}{white}{poppurple}\ifx\relax#1\relax\else\hspace{6pt}{\popxbold\color{white}#1}\fi\par\vspace{7pt}}}
% Exercice résolu : énoncé en haut, \tcblower puis la solution sur fond crème
\newcounter{popexo}\newcounter{popexr}
\newtcolorbox{exoresolu}[1][]{popbase,colframe=poporange,colbacklower=popcream,
  segmentation style={popink,line width=1.2pt,dashed},
  before upper={\stepcounter{popexr}\popboxhead{Exercice résolu \thepopexr}{poporange}{#1}},
  before lower={\poppill{Solution}{popink}{popcream}\par\vspace{6pt}}}
% Exercice (non résolu en place) — la correction est dans la partie Corrigés
\newtcolorbox{exercice}[1][]{popbase,colframe=popink,
  before upper={\stepcounter{popexo}\popboxhead{Exercice \thepopexo}{popink}{#1}}}
% Corrigé
\newtcolorbox{corrige}[1][]{popbase,colframe=popgreen,colback=popgreenL,
  before upper={\popboxhead{Corrigé}{popgreen}{#1}}}
% Objectifs / prérequis / bilan
\newtcolorbox{objectifs}{popbase,colframe=popink,colback=poporangeL,
  before upper={\popboxhead{Objectifs de la leçon}{popink}{}}}
\newtcolorbox{prerequis}{popbase,colframe=popink,colback=popblueL,
  before upper={\popboxhead{Prérequis}{popblue}{}}}
\newtcolorbox{bilan}{popbase,colframe=popink,colback=white,boxrule=2.8pt,
  before upper={\popboxhead{Fiche bilan}{poporange}{L'essentiel à connaître pour l'examen}}}
% Figure encadrée avec légende
\newtcolorbox{popfigure}[1][]{enhanced,breakable=false,colback=white,colframe=popline,boxrule=1.4pt,arc=8pt,
  left=10pt,right=10pt,top=10pt,bottom=8pt,before skip=10pt,after skip=12pt,halign=center,
  lower separated=false,halign lower=center,
  fontlower=\popsemi\fontsize{9}{11.5}\selectfont\color{popmuted},
  #1}
\newcommand{\legende}[1]{\par\vspace{6pt}{\popsemi\fontsize{9}{11.5}\selectfont\color{popmuted}#1}}

% ============================================================
% Signature OnBuch+
% ============================================================
\newcommand{\poplogoinline}[1]{{\poplogo\fontsize{#1}{#1}\selectfont\color{popink}OnBuch\color{poporange}+}}
\newcommand{\signatureonbuch}{%
  \begin{tcolorbox}[enhanced,colback=popink,colframe=popink,boxrule=2.2pt,arc=10pt,
      drop shadow=poporange,left=14pt,right=14pt,top=10pt,bottom=10pt,before skip=0pt,after skip=0pt]
    \tikz[baseline=-0.7ex]\node[circle,fill=popgreen,draw=popcream,line width=1.3pt,minimum size=22pt,inner sep=0]{\popcheck{white}};%
    \hspace{9pt}\begin{minipage}[c]{0.62\linewidth}
      {\popxbold\fontsize{12}{13}\selectfont\color{popcream}Signé\ \textbullet\ {\poplogo OnBuch\color{poporange}+}}\par\vspace{2pt}
      {\raggedright\popsemi\fontsize{8}{10.5}\selectfont\color{popline}\MakeUppercase{Équipe pédagogique OnBuch+}\\\MakeUppercase{Conforme au programme officiel MINESEC}\par}
    \end{minipage}\hfill
    {\poplogo\fontsize{22}{22}\selectfont\color{popcream}OnBuch\color{poporange}+}
  \end{tcolorbox}%
}

% ============================================================
% Page de garde
% #1 titre · #2 matière · #3 niveau · #4 module · #5 n° leçon · #6 durée ·
% #7 description / objectifs (texte ou itemize)
% ============================================================
\newcommand{\pagegarde}[7]{%
  \thispagestyle{empty}
  \newgeometry{a4paper,margin=1.7cm,top=1.6cm,bottom=1.4cm}%
  \begin{tikzpicture}[remember picture,overlay]
    \fill[poporange!18] ([xshift=-3.2cm,yshift=-3.6cm]current page.north east) circle (4.6cm);
    \fill[poppurple!14] ([xshift=2.4cm,yshift=7.6cm]current page.south west) circle (3.8cm);
  \end{tikzpicture}%
  {\poplogo\fontsize{30}{30}\selectfont\color{popink}OnBuch\color{poporange}+}\hfill
  \raisebox{0.6ex}{\poppill{Cours signé \textbullet\ Premium}{popink}{popcream}}\par
  \vspace{1pt}\popsquiggle{poppurple}\par
  \vspace{8pt}
  \begin{tcolorbox}[enhanced,colback=poporange,colframe=popink,boxrule=2.8pt,arc=13pt,
      drop shadow=popink,left=18pt,right=18pt,top=12pt,bottom=13pt,after skip=10pt]
    \poppill{#2 \textbullet\ #3}{popink}{popcream}\hspace{5pt}\poppill{#4}{white}{popink}\par\vspace{8pt}
    {\popdisplay\fontsize{14}{14}\selectfont\color{popink}LEÇON #5}\par\vspace{4pt}
    {\raggedright\hyphenpenalty=10000\exhyphenpenalty=10000\popdisplay\fontsize{25}{27}\selectfont\color{popcream}\MakeUppercase{#1}\par}
    \vspace{8pt}
    {\popxbold\fontsize{9.5}{9.5}\selectfont\color{popink}\MakeUppercase{Durée conseillée : #6}}
  \end{tcolorbox}
  \begin{tcolorbox}[enhanced,colback=white,colframe=popink,boxrule=2.2pt,arc=10pt,
      drop shadow=popink,left=16pt,right=16pt,top=10pt,bottom=10pt,after skip=8pt]
    \poppill{Dans cette leçon}{poppurple}{white}\par\vspace{5pt}
    \popsemi\fontsize{9.3}{12.6}\selectfont\color{popink}#7
  \end{tcolorbox}
  \noindent\begin{minipage}[t]{0.487\linewidth}
    \begin{tcolorbox}[enhanced,colback=popgreen,colframe=popink,boxrule=2.2pt,arc=10pt,
        drop shadow=popink,left=11pt,right=11pt,top=10pt,bottom=10pt,height=3.1cm,valign=center]
      \tcbox[colback=white,colframe=popink,boxrule=1.3pt,arc=5pt,boxsep=0pt,nobeforeafter,
        left=3pt,right=3pt,top=3pt,bottom=3pt,valign=center]{\qrcode[height=1.9cm]{https://play.google.com/store/apps/details?id=com.luvvix.onbuch&hl=fr}}%
      \hspace{8pt}\begin{minipage}[c]{0.5\linewidth}
        {\popdisplay\fontsize{11}{12.5}\selectfont\color{white}\MakeUppercase{Télécharge OnBuch}}\par\vspace{4pt}
        {\raggedright\popsemi\fontsize{8.4}{11}\selectfont\color{white}Gratuit \textbullet\ Google Play\\Tuteur IA, annales, concours, résultats\par}
      \end{minipage}
    \end{tcolorbox}
  \end{minipage}\hfill
  \begin{minipage}[t]{0.487\linewidth}
    \begin{tcolorbox}[enhanced,colback=poppurple,colframe=popink,boxrule=2.2pt,arc=10pt,
        drop shadow=popink,left=11pt,right=11pt,top=10pt,bottom=10pt,height=3.1cm,valign=center]
      \tikz[baseline=-0.7ex]\node[circle,fill=white,draw=popink,line width=1.3pt,minimum size=1.6cm,inner sep=0]{\popdisplay\fontsize{24}{24}\selectfont\color{poppurple}+};%
      \hspace{8pt}\begin{minipage}[c]{0.6\linewidth}
        {\popdisplay\fontsize{11}{12.5}\selectfont\color{white}\MakeUppercase{Passe à OnBuch+}}\par\vspace{4pt}
        {\raggedright\popsemi\fontsize{8.4}{11}\selectfont\color{white}Tous les cours signés, Tuteur IA illimité, fiches PDF — onglet OnBuch+ de l'app\par}
      \end{minipage}
    \end{tcolorbox}
  \end{minipage}
  \vspace{8pt}
  \popcontenu
  \vfill
  \signatureonbuch
  \restoregeometry
  \newpage
}

% Rangée « Ce cours contient » (page de garde)
\newcommand{\poptile}[3]{% #1 couleur #2 gros texte #3 légende
  \begin{tcolorbox}[enhanced,colback=#1,colframe=popink,boxrule=2pt,arc=9pt,shadow={2.5pt}{-2.5pt}{0pt}{popink},
      left=6pt,right=6pt,top=9pt,bottom=9pt,halign=center,valign=center,height=1.75cm,width=\linewidth,nobeforeafter]
    {\popdisplay\fontsize{12.5}{13}\selectfont\color{white}\MakeUppercase{#2}}\par\vspace{3pt}
    {\centering\popsemi\fontsize{7.8}{9.5}\selectfont\color{white}#3\par}
  \end{tcolorbox}}
\newcommand{\popcontenu}{%
  \par\noindent{\popxboldls\fontsize{7.5}{7.5}\selectfont\color{popmuted}CE COURS SIGNÉ CONTIENT}\par\vspace{7pt}
  \noindent\begin{minipage}[t]{0.235\linewidth}\poptile{popblue}{Cours}{complet, illustré, pas à pas}\end{minipage}\hfill
  \begin{minipage}[t]{0.235\linewidth}\poptile{poporange}{Méthodes}{et exercices résolus}\end{minipage}\hfill
  \begin{minipage}[t]{0.235\linewidth}\poptile{poppink}{Exercices}{type examen + corrigés}\end{minipage}\hfill
  \begin{minipage}[t]{0.235\linewidth}\poptile{popgreen}{Fiche bilan}{l'essentiel à retenir}\end{minipage}\par}

% Sommaire
\newtcolorbox{sommaire}{enhanced,colback=white,colframe=popink,boxrule=2.2pt,arc=10pt,
  drop shadow=popink,left=18pt,right=18pt,top=13pt,bottom=13pt,before skip=6pt,after skip=20pt,
  before upper={\poppill{Sommaire}{poppurple}{white}\par\vspace{9pt}\popsemi\fontsize{10.6}{14.5}\selectfont\color{popink}}}

% Signature de fin de cours — poussée tout en bas de la dernière page
\newcommand{\findecours}{%
  \par\vfill\nopagebreak
  \begin{center}\popsquiggle{poporange}\end{center}\vspace{4pt}
  \signatureonbuch
}
\AtBeginDocument{\def\llbracket{\mathopen{[\mkern-3.5mu[}}\def\rrbracket{\mathclose{]\mkern-3.5mu]}}} % intervalles d'entiers (glyphe ⟦ absent de la police)
% Glyphes absents de Fira Math : redéfinis à partir de glyphes disponibles
\AtBeginDocument{%
  \def\setminus{\mathbin{\backslash}}\let\smallsetminus\setminus
  \def\vdots{\mathord{\vbox{\baselineskip3.2pt\lineskiplimit0pt\kern4pt\hbox{.}\hbox{.}\hbox{.}}}}%
  \def\ddots{\mathinner{\mkern1mu\raise6.5pt\hbox{.}\mkern2mu\raise3.6pt\hbox{.}\mkern2mu\raise0.7pt\hbox{.}\mkern1mu}}%
  \def\triangle{\mathord{\scalebox{0.9}{$\symup{\Delta}$}}}%
  \def\nabla{\mathord{\raisebox{\depth}{\scalebox{1}[-1]{$\symup{\Delta}$}}}}%
  \def\checkmark{\popcheck{popdark}}%
  \def\star{\mathbin{\ast}}\def\bigstar{\mathord{\ast}}%
  \def\diamond{\mathbin{\scalebox{0.6}{$\Diamond$}}}%
  \def\bigcup{\mathop{\mathchoice{\raisebox{-0.3em}{\scalebox{1.7}{$\cup$}}}{\scalebox{1.25}{$\cup$}}{\cup}{\cup}}\displaylimits}%
  \def\bigcap{\mathop{\mathchoice{\raisebox{-0.3em}{\scalebox{1.7}{$\cap$}}}{\scalebox{1.25}{$\cap$}}{\cap}{\cap}}\displaylimits}%
  \def\lesseqqgtr{\mathrel{\lessgtr}}\def\bigoplus{\mathop{\oplus}\displaylimits}%
}
\providecommand{\ssi}{si et seulement si} % abréviation fréquente
\AtBeginDocument{\providecommand{\wideparen}{\overparen}} % arc de cercle

%% FIGFIX-BEGIN v3 — correctifs globaux des figures (docs/onbuch-generation/tools/figfix)
\usepackage{adjustbox}\usepackage{etoolbox}
% toute figure TikZ/pgfplots plus large que la ligne est réduite à la largeur disponible
\BeforeBeginEnvironment{tikzpicture}{\begin{adjustbox}{max width=\linewidth}}
\AfterEndEnvironment{tikzpicture}{\end{adjustbox}}
% légendes pgfplots lisibles par-dessus les courbes
\pgfplotsset{every axis/.append style={legend style={fill=white,fill opacity=0.92,text opacity=1,draw=black!15,inner sep=3pt}}}
% étiquettes d'arêtes (syntaxe edge["..."]) avec fond blanc
\tikzset{every edge quotes/.append style={fill=white,inner sep=1.2pt}}
% bulles de cartes mentales : texte à la ligne dans la bulle, bulle un peu plus grande
\tikzset{every concept/.append style={text width=2.0cm,align=center,minimum size=2.3cm,inner sep=2pt}}
%% FIGFIX-END
\begin{document}

\pagegarde{Lectura y comentario : la familia y sus problemas (viudez, herencia) ; juego de rol}{Espagnol}{Tle A}{Módulo 1 : Vie familiale et intégration sociale}{E05}{2 h}{Cette leçon propose la lecture et le commentaire d'un texte sur les problèmes familiaux liés au veuvage et à l'héritage, l'acquisition du lexique du deuil et de la succession, puis un jeu de rôle : une conversation familiale sur le partage d'un héritage.
\vspace{4pt}
\begin{multicols}{2}\small
\begin{itemize}
\item À la fin de cette leçon, je sais lire et comprendre globalement et finement un texte sur les problèmes familiaux (veuvage, héritage).
\item À la fin de cette leçon, je sais identifier les personnages, la situation et les enjeux d'un récit familial.
\item À la fin de cette leçon, je sais utiliser correctement le lexique du deuil et de l'héritage (viudez, herencia, testamento, huérfano...).
\item À la fin de cette leçon, je sais dégager les thèmes et la problématique d'un texte pour en faire un commentaire organisé.
\item À la fin de cette leçon, je sais rédiger un commentaire de texte en espagnol avec introduction, développement et conclusion.
\item À la fin de cette leçon, je sais exprimer un désaccord, une émotion et une proposition de réconciliation en espagnol.
\end{itemize}\end{multicols}}

\begin{sommaire}
\begin{multicols}{2}
\begin{enumerate}
\item Texte support : « La herencia de don Ramón »
\item Compréhension du texte
\item Lexique thématique : deuil et héritage
\item Méthode du commentaire de texte
\item Jeu de rôle : la conversation familiale
\item Production guidée et bilan
\item Activité d'intégration
\item Méthodes et exercices résolus
\item Exercices
\item Corrigés des exercices
\item Fiche bilan
\end{enumerate}
\end{multicols}
\end{sommaire}

\begin{objectifs}
\begin{itemize}
\item À la fin de cette leçon, je sais lire et comprendre globalement et finement un texte sur les problèmes familiaux (veuvage, héritage).
\item À la fin de cette leçon, je sais identifier les personnages, la situation et les enjeux d'un récit familial.
\item À la fin de cette leçon, je sais utiliser correctement le lexique du deuil et de l'héritage (viudez, herencia, testamento, huérfano...).
\item À la fin de cette leçon, je sais dégager les thèmes et la problématique d'un texte pour en faire un commentaire organisé.
\item À la fin de cette leçon, je sais rédiger un commentaire de texte en espagnol avec introduction, développement et conclusion.
\item À la fin de cette leçon, je sais exprimer un désaccord, une émotion et une proposition de réconciliation en espagnol.
\item À la fin de cette leçon, je sais jouer une scène de conversation familiale (partage d'héritage, réconciliation) avec un lexique et des structures appropriés.
\item À la fin de cette leçon, je sais comparer les pratiques successorales du Cameroun et du monde hispanophone.
\end{itemize}
\end{objectifs}

\begin{prerequis}
\begin{itemize}
\item Lexique de la famille et des liens de parenté (Première : padres, hermanos, tíos, abuelos, cónyuge).
\item Passé simple et imparfait de l'indicatif pour raconter des événements passés.
\item Subjonctif présent après les verbes de sentiment et de volonté (espero que..., quiero que...).
\item Expressions de l'opinion et de l'accord/désaccord (creo que, me parece que, no estoy de acuerdo).
\item Technique de la question et de la réponse en compréhension écrite (Première).
\end{itemize}
\end{prerequis}

\newpage

\coursec{Texte support : « La herencia de don Ramón »}

\courssub{Mise en situation et problématique}

À Douala comme à Madrid, la mort d'un père de famille soulève souvent des questions délicates : qui hérite de la maison ? Comment partager les biens entre les enfants ? Au Cameroun, les conflits successoraux déchirent parfois les familles : au village ou à Yaoundé, le partage des terres et des maisons du défunt peut opposer les frères et sœurs pendant des années. En Espagne aussi, le \esp{reparto de la herencia} (le partage de l'héritage) est une source fréquente de tensions familiales. Aujourd'hui, nous lisons le récit d'une famille espagnole confrontée au veuvage et à l'héritage, puis nous jouerons nous-mêmes une scène de partage. À vous de trouver les mots justes... en espagnol !

\textbf{Problématique :} comment une famille déchirée par la mort et l'argent peut-elle se réconcilier ? Quel vocabulaire espagnol nous faut-il pour raconter et discuter ces situations ?

\courssub{Présentation du texte}

\begin{itemize}
\item \textbf{Genre :} récit narratif à la troisième personne, en Espagne contemporaine.
\item \textbf{Personnages :} don Ramón (le père, récemment décédé), Carmen (sa veuve), et leurs trois enfants : Luis (l'aîné), Ana (la cadette) et Pablo (le benjamin).
\item \textbf{Situation initiale :} la famille se réunit pour la lecture du testament.
\item \textbf{Conflit :} Luis veut vendre la maison familiale ; Ana veut la garder ; Pablo réclame sa part en argent.
\item \textbf{Dénouement :} ouvert — une réunion familiale où la réconciliation semble possible.
\end{itemize}

\courssub{Le texte}

\begin{textebox}[La herencia de don Ramón]
Cuando don Ramón murió, la familia entera quedó sumida en el luto. Su viuda, Carmen, y sus tres hijos se reunieron en el salón de la casa familiar para la lectura del testamento. El notario abrió el sobre con solemnidad y leyó las últimas voluntades del difunto: la casa, los ahorros y el pequeño terreno del pueblo debían repartirse entre los tres herederos.

Luis, el mayor, habló el primero:
—La casa es grande y nadie vive en ella. Lo mejor es venderla y repartir el dinero.
Ana se negó enérgicamente:
—¡Jamás! Esta casa es el hogar de nuestra infancia. Aquí jugábamos, aquí celebrábamos las fiestas con papá. No se puede vender.
Pablo, el menor, solamente pedía su parte de la herencia:
—Yo necesito el dinero para mis estudios. Que cada uno reciba lo que le corresponde.

La tensión crecía. Los hermanos se miraban con rencor y las palabras amables se convertían en reproches. Carmen escuchaba en silencio, con los ojos llenos de lágrimas. De pronto, tomó la palabra:
—Hijos míos, vuestro padre trabajó toda su vida para veros unidos, no enfrentados. Esta casa no vale nada si destruye a la familia. Sentémonos, hablemos con calma y busquemos una solución juntos.

Los tres hermanos bajaron la voz. Por primera vez desde la muerte de su padre, se miraron no como rivales, sino como hermanos. La reconciliación, tal vez, era posible.
\end{textebox}

\textbf{Glossaire :}

\begin{center}
\begin{tabularx}{\linewidth}{|l|X|l|X|}
\hline
\rowcolor{popblueL} \textbf{Español} & \textbf{Français} & \textbf{Español} & \textbf{Français} \\
\hline
el luto & le deuil & el testamento & le testament \\
\hline
la viuda & la veuve & el difunto & le défunt \\
\hline
el notario & le notaire & los ahorros & les économies \\
\hline
el terreno & le terrain & el heredero & l'héritier \\
\hline
repartirse & se partager & negarse & refuser \\
\hline
el hogar & le foyer & el rencor & la rancune \\
\hline
el reproche & le reproche & enfrentados & opposés \\
\hline
la reconciliación & la réconciliation & tomar la palabra & prendre la parole \\
\hline
\end{tabularx}
\end{center}

\begin{definition}[La viudez]
La \cle{viudez} (le veuvage) est la situation d'une personne dont le conjoint est décédé. L'homme est \esp{viudo} (veuf), la femme est \esp{viuda} (veuve). Un enfant qui a perdu son père ou sa mère est \esp{huérfano/a} (orphelin/orpheline).
\esp{Carmen enviudó a los cincuenta años.} « Carmen est devenue veuve à cinquante ans. »
\end{definition}

\courssub{Première lecture : démarche}

\begin{methode}[Comment aborder le texte]
\begin{enumerate}
\item \textbf{Lecture silencieuse individuelle} (3 minutes) : souligne les mots inconnus sans chercher dans le glossaire.
\item \textbf{Lecture à voix haute par l'enseignant} : écoute le modèle de prononciation (attention : \esp{herencia} se prononce [eˈɾenθja] / [eˈɾensja], le \emph{h} est muet ; \esp{luto} [ˈluto]).
\item \textbf{Hypothèses de sens} : à partir du contexte, devine le sens des mots soulignés avant de vérifier dans le glossaire. Exemple : dans \esp{quedó sumida en el luto}, le contexte (\esp{murió}) permet de deviner que \esp{luto} désigne le deuil.
\item \textbf{Vérification} avec le glossaire, puis relecture globale pour saisir le conflit.
\end{enumerate}
\end{methode}

\begin{exoresolu}[Hypothèses de sens à partir du contexte]
Trouve le sens des mots soulignés sans dictionnaire :
1. \esp{Su \textbf{viuda}, Carmen, y sus tres hijos se reunieron.} (don Ramón vient de mourir)
2. \esp{El notario leyó las últimas voluntades del \textbf{difunto}.}
3. \esp{La casa debía \textbf{repartirse} entre los tres herederos.}
\tcblower
1. \esp{Viuda} = veuve : c'est la femme de don Ramón, qui vient de mourir.
2. \esp{Difunto} = défunt : celui dont on lit les « dernières volontés », donc le mort.
3. \esp{Repartirse} = se partager : la maison est divisée entre trois personnes, les héritiers.
\end{exoresolu}

\courssub{Les personnages et le conflit}

\begin{popfigure}
\begin{tikzpicture}[
box/.style={draw=popblue, fill=popblueL, rounded corners, align=center, font=\small, inner sep=4pt},
flecha/.style={-{Stealth[length=2.5mm]}, thick, popdark}]
\node[box, fill=popgoldL, draw=popgold, align=center] (ramon) at (0,2.4){don Ramón $\dagger$\\(el difunto)};
\node[box, fill=poppinkL, draw=poppink, align=center] (carmen) at (4.4,2.4){Carmen\\(la viuda)};
\draw[thick, popdark] (ramon) -- (carmen) node[midway, above, font=\scriptsize] {matrimonio};
\node[box, align=center] (luis) at (0,0){Luis\\el mayor\\quiere vender};
\node[box, align=center] (ana) at (4.4,0){Ana\\la hija\\quiere quedarse la casa};
\node[box, align=center] (pablo) at (8.8,0){Pablo\\el menor\\pide su parte en dinero};
\draw[flecha] (2.2,2.1) -- (luis.north);
\draw[flecha] (2.2,2.1) -- (ana.north);
\draw[flecha] (2.2,2.1) -- (pablo.north);
\draw[flecha, poporange] (luis.east) -- node[above, font=\scriptsize, poporange]{conflicto} (ana.west);
\draw[flecha, poporange] (ana.east) -- node[above, font=\scriptsize, poporange]{conflicto} (pablo.west);
\end{tikzpicture}
\legende{Arbre familial et flèches du conflit : les trois enfants s'opposent sur le sort de la maison.}
\end{popfigure}

\courssub{Carte mentale du thème}

\begin{popfigure}
\begin{tikzpicture}[mindmap, grow cyclic, every node/.style={concept, align=center, font=\small},
root concept/.append style={concept color=popblue, font=\bfseries},
level 1/.append style={level distance=3.4cm, sibling angle=90},
level 2/.append style={level distance=2.2cm, sibling angle=45, font=\scriptsize}]
\node[root concept]{la herencia}
child[concept color=popgreen]{node{personajes}
child{node{la viuda}}
child{node{los herederos}}
child{node{el notario}}}
child[concept color=poporange]{node{bienes}
child{node{la casa}}
child{node{los ahorros}}
child{node{el terreno}}}
child[concept color=poppink]{node{conflictos}
child{node{vender / guardar}}
child{node{los reproches}}
child{node{el rencor}}}
child[concept color=popgold]{node{solución}
child{node{dialogar}}
child{node{reconciliarse}}};
\end{tikzpicture}
\legende{Carte mentale lexicale autour de \esp{la herencia} : personnages, biens, conflits, solution.}
\end{popfigure}

\courssub{Lexique fondamental : deuil et héritage}

\begin{center}
\begin{tabularx}{\linewidth}{|l|X|X|}
\hline
\rowcolor{popblueL} \textbf{Español} & \textbf{Français} & \textbf{Exemple} \\
\hline
la viudez & le veuvage & \esp{La viudez es una etapa difícil.} \\
\hline
viudo / viuda & veuf / veuve & \esp{Carmen es viuda desde hace un mes.} \\
\hline
el luto & le deuil & \esp{La familia está de luto.} \\
\hline
huérfano/a & orphelin/e & \esp{Los niños quedaron huérfanos de padre.} \\
\hline
la herencia & l'héritage (les biens) & \esp{La herencia incluye una casa.} \\
\hline
el/la heredero/a & l'héritier / l'héritière & \esp{Los tres hijos son herederos.} \\
\hline
el testamento & le testament & \esp{El notario leyó el testamento.} \\
\hline
el reparto & le partage & \esp{El reparto de los bienes fue justo.} \\
\hline
reconciliarse & se réconcilier & \esp{Los hermanos se reconciliaron.} \\
\hline
el conflicto familiar & le conflit familial & \esp{El dinero causa conflictos familiares.} \\
\hline
\end{tabularx}
\end{center}

\begin{attention}[Faux amis et confusions fréquentes]
\esp{Herencia} signifie « héritage » au sens des \textbf{biens transmis}, pas la personne ! Ne confonds pas :
\begin{itemize}
\item \esp{la herencia} = l'héritage (les biens) : \esp{Recibió una gran herencia.} « Il a reçu un grand héritage. »
\item \esp{el heredero / la heredera} = l'héritier / l'héritière (la personne) : \esp{Luis es el heredero mayor.} « Luis est l'héritier aîné. »
\item \esp{la sucesión} = la succession (la procédure juridique) : \esp{La sucesión se abrió ante el notario.}
\end{itemize}
Autre piège : \esp{heredar} = hériter (de) : \esp{Heredó la casa de su padre.} « Il a hérité de la maison de son père. » En français on dit « hériter de quelque chose », en espagnol aussi : \esp{heredar algo de alguien}.
\end{attention}

\begin{savaistu}[La legítima : l'héritage protégé]
En Espagne comme au Cameroun, le partage de l'héritage est une cause fréquente de conflits familiaux. Mais les droits diffèrent : en droit espagnol, une part appelée \esp{la legítima} est obligatoirement réservée aux enfants — on ne peut pas déshériter totalement ses descendants. Au Cameroun, coexistent le droit moderne et les coutumes : dans certaines traditions, ce sont les neveux ou les oncles maternels qui héritent, ce qui provoque parfois des disputes entre la famille nucléaire et la famille élargie. D'où l'importance du dialogue... et du notaire !
\end{savaistu}

\courssub{Point de grammaire : Le récit au passé (Indéfini vs Imparfait)}

\begin{propriete}[Choix entre Pretérito Indefinido et Pretérito Imperfecto]
Dans un récit, l'espagnol distingue deux passés principaux :
\begin{itemize}
\item \textbf{Indéfini (\emph{pretérito indefinido})} : actions \textbf{ponctuelles, achevées, qui font avancer l'intrigue} (suite d'événements). Marqueurs : \esp{ayer, anoche, el otro día, de repente, entonces, primero}.
\item \textbf{Imparfait (\emph{pretérito imperfecto})} : \textbf{descriptions, habitudes, actions en cours, contexte} (arrière-plan). Marqueurs : \esp{mientras, antes, siempre, a menudo, en aquel tiempo}.
\end{itemize}
\end{propriete}

\begin{exemplebox}[Analyse du texte]
\begin{center}
\begin{tabularx}{\linewidth}{|l|X|l|}
\hline
\rowcolor{popblueL} \textbf{Verbe (texte)} & \textbf{Infinitif / Temps} & \textbf{Justification} \\
\hline
\esp{murió} & morir (Indéfini) & Action ponctuelle, point de départ du récit. \\
\hline
\esp{quedó} & quedar (Indéfini) & Résultat immédiat, changement d'état. \\
\hline
\esp{se reunieron} & reunirse (Indéfini) & Action achevée, événement principal. \\
\hline
\esp{abrió / leyó} & abrir / leer (Indéfini) & Suite d'actions ponctuelles (notaire). \\
\hline
\esp{debían repartirse} & deber + inf. (Imparfait) & Obligation existante à ce moment (contexte). \\
\hline
\esp{habló / negó / pedía} & hablar / negar / pedir & \esp{habló, negó} = actions ponctuelles ; \esp{pedía} = attitude/volonté en cours (imparfait). \\
\hline
\esp{crecía / miraban / se convertían / escuchaba} & crecer / mirar / convertirse / escuchar (Imparfait) & Description de l'atmosphère, actions en cours simultanées. \\
\hline
\esp{tomó / trabajó / veros / enfrentados / vale / destruye / sentémonos / hablemos / busquemos / bajaron / miraron / era} & Mixte & \esp{tomó, trabajó, bajaron, miraron} = actions ponctuelles (Indéfini). \esp{veros, enfrentados} = infinitif / participe. \esp{vale, destruye} = présent (dialogue direct). \esp{sentémonos, hablemos, busquemos} = \textbf{subjonctif présent} (impératif \emph{nosotros} / exhortation). \esp{era} = imparfait (appréciation finale). \\
\hline
\end{tabularx}
\end{center}
\end{exemplebox}

\begin{attention}[Piège francophone : « Il mourait » vs « Il est mort »]
En français, l'imparfait narratif (\emph{« il mourait »}) n'existe pas en espagnol pour une action principale. \esp{Moría} implique une agonie longue ou une habitude (\emph{« chaque hiver il faiblissait »}). Pour dire « il est mort » (événement unique) : \esp{murió} (Indéfini).
\end{attention}

\begin{exercice}[Grammaire : Indéfini ou Imparfait ?]
Conjugue le verbe entre parenthèses au temps approprié.
1. Ayer \_\_\_\_\_\_ (llegar) el notario a las diez.
2. Mientras él \_\_\_\_\_\_ (leer) el testamento, los hijos \_\_\_\_\_\_ (escuchar) en silencio.
3. De pronto, Carmen \_\_\_\_\_\_ (empezar) a llorar.
4. Antes, la familia \_\_\_\_\_\_ (ser) muy unida.
\end{exercice}

\begin{corrige}[Exercice Grammaire]
1. \esp{llegó} (action ponctuelle, \emph{ayer}). 2. \esp{leía} (action en cours, \emph{mientras}) / \esp{escuchaban} (action simultanée, description). 3. \esp{empezó} (rupture, \emph{de pronto}). 4. \esp{era} (description état passé, \emph{antes}).
\end{corrige}

\courssub{Méthode : Le commentaire de texte (Bac)}

\begin{methode}[Étapes du commentaire composé]
\begin{enumerate}
\item \textbf{Introduction} (1 paragraphe) :
\begin{itemize}
\item \textbf{Accroche} : thème général (famille, héritage, conflits).
\item \textbf{Présentation} : titre (entre guillemets), auteur (ici : texte d'examen / anonyme), genre (récit), thème principal.
\item \textbf{Problématique} : question à laquelle le texte répond (ex: \emph{Comment le dialogue permet-il de surmonter le conflit successoral ?}).
\item \textbf{Annonce du plan} : 2 ou 3 parties logiques.
\end{itemize}
\item \textbf{Développement} (2 ou 3 parties, chacune = 1 idée directrice) :
\begin{itemize}
\item \textbf{Idée directrice} (phrase d'intro de paragraphe).
\item \textbf{Argument} : citation courte (\esp{« ... »}) + analyse (vocabulaire, temps, figures, point de vue).
\item \textbf{Mini-synthèse} : lien avec la problématique.
\end{itemize}
\item \textbf{Conclusion} :
\begin{itemize}
\item \textbf{Bilan} : réponse à la problématique.
\item \textbf{Ouverture} : lien avec un autre texte, l'actualité, le contexte camerounais, expérience personnelle.
\end{itemize}
\end{enumerate}
\end{methode}

\begin{exemplebox}[Plan type pour « La herencia de don Ramón »]
\textbf{Problématique :} \emph{En quoi la figure maternelle incarne-t-elle la solution au conflit successoral ?}

\textbf{I. Une famille déchirée par l'argent et le deuil}
\begin{itemize}
\item A. Le cadre : la lecture du testament (formel, froid).
\item B. L'opposition des trois héritiers : vendre (Luis) / garder (Ana) / argent (Pablo).
\item C. Lexique du conflit : \esp{rencor, reproches, enfrentados, tensión}.
\end{itemize}

\textbf{II. Carmen, figure médiatrice et gardienne de l'unité}
\begin{itemize}
\item A. Le silence puis la parole : \esp{escuchaba en silencio} $\rightarrow$ \esp{tomó la palabra}.
\item B. L'argument affectif vs matériel : \esp{veros unidos, no enfrentados} ; \esp{no vale nada si destruye a la familia}.
\item C. L'usage de l'impératif exhortatif (subjonctif) : \esp{sentémonos, hablemos, busquemos}.
\end{itemize}

\textbf{III. Une réconciliation possible mais fragile}
\begin{itemize}
\item A. Changement d'attitude : \esp{bajaron la voz, se miraron no como rivales, sino como hermanos}.
\item B. L'incertitude finale : \esp{tal vez, era posible} (imparfait d'hypothèse / probabilité).
\end{itemize}

\textbf{Conclusion :} Le texte montre que l'héritage matériel ne doit pas primer sur l'héritage affectif. Ouverture : comparaison avec les traditions successorales au Cameroun (droit coutumier vs droit moderne).
\end{exemplebox}

\courssub{Jeu de rôle : La reunión familiar}

\begin{experience}[Mise en scène : Le partage de l'héritage]
\textbf{Objectif :} Réinvestir le lexique (héritage, conflit, réconciliation) et la grammaire (passés, subjonctif d'exhortation, conditionnel de politesse) dans une interaction orale.

\textbf{Distribution des rôles (groupes de 5) :}
\begin{itemize}
\item \textbf{El Notario} : formel, neutre, lit le testament, gère le temps de parole.
\item \textbf{Carmen (la viuda)} : émue, pacificatrice, utilise \emph{nosotros} + subjonctif (\esp{hablemos, busquemos}).
\item \textbf{Luis (el mayor)} : pragmatique, veut vendre, arguments financiers. Utilise \esp{creo que debemos vender...}, \esp{es mejor que vendamos...} (subjonctif après opinion).
\item \textbf{Ana (la hija)} : attachée au patrimoine affectif, refuse la vente. \esp{No quiero que vendáis la casa}, \esp{exijo que se respete la memoria de papá}.
\item \textbf{Pablo (el menor)} : étudiant, a besoin d'argent, pressé. \esp{Necesito mi parte ahora}, \esp{Quiero que me den mi dinero}.
\end{itemize}

\textbf{Déroulement (10 min préparation + 5 min jeu) :}
\begin{enumerate}
\item \textbf{Phase 1 : Lecture du testament} (Notario). Les autres écoutent, prennent des notes.
\item \textbf{Phase 2 : Tour de table des positions} (Luis, Ana, Pablo). Chacun expose son point de vue (1 min chacun).
\item \textbf{Phase 3 : Médiation de Carmen} : Elle propose une solution (ex: location de la maison, partage des loyers, ou vente différée).
\item \textbf{Phase 4 : Négociation finale} : Accord ou désaccord. Utiliser \esp{Estoy de acuerdo / No estoy de acuerdo / Propongo que... / Acepto si...}.
\end{enumerate}

\textbf{Expressions utiles (à afficher au tableau) :}
\begin{center}
\begin{tabularx}{\linewidth}{|X|X|}
\hline
\rowcolor{popblueL} \textbf{Exprimer son avis / Exiger} & \textbf{Proposer / Négocier} \\
\hline
\esp{En mi opinión, deberíamos...} & \esp{Propongo que...} \\
\esp{Exijo que se respete...} & \esp{¿Y si... vendemos / alquilamos...?} \\
\esp{No pienso ceder en...} & \esp{Acepto, con la condición de que...} \\
\esp{Es fundamental para mí que...} & \esp{Busquemos una solución justa.} \\
\hline
\end{tabularx}
\end{center}

\textbf{Évaluation par les pairs (grille) :} Prononciation / Richesse lexicale (héritage) / Grammaire (subjonctif, passés) / Interaction (écoute, réactivité) / Respect du rôle.
\end{experience}

\courssub{Production écrite guidée}

\begin{exercice}[Choisissez UN sujet (150--180 mots)]
\begin{enumerate}
\item \textbf{Suite du récit} : Écrivez la suite de la réunion familiale. Les frères et sœurs trouvent un accord (ou pas). Utilisez le dialogue direct et le narrateur. Temps : Indéfini / Imparfait / Subjonctif.
\item \textbf{Lettre formelle} : Vous êtes l'avocat de la famille. Écrivez une lettre à \esp{Doña Carmen} pour lui expliquer la procédure de \esp{partición de herencia} et les droits de \esp{la legítima}. Registre formel (\esp{usted}, conditionnel, futur).
\item \textbf{Article d'opinion} : « L'argent détruit-il la famille ? » Rédigez un article pour le journal du lycée (\esp{El Eco del Liceo}) en vous appuyant sur le texte et des exemples camerounais (terres, maisons, coutumes). Structure : Thèse -- Antithèse -- Synthèse.
\end{enumerate}
\end{exercice}

\begin{corrige}[Pistes de correction - Sujet 1 (Suite du récit)]
\textbf{Extrait modèle :}
\esp{—Propongo —dijo Luis, suavizando el tono— que alquilemos la casa. Los ingresos se reparten en tres partes y Pablo paga sus estudios. Dentro de cinco años, decidimos si vendemos.}
\esp{Ana lo pensó un momento. Las lágrimas se secaron en sus mejillas. —De acuerdo —susurró—. Pero la habitación de papá queda intacta. Es... nuestro rincón de memoria.}
\esp{Pablo sonrió, aliviado. —Me parece justo. Gracias, hermanos.}
\esp{Carmen los abrazó a los tres. —Vuestro padre estaría orgulloso. La legítima no es el dinero, es el amor.}
\esp{El notario cerró la carpeta. Afuera, el sol iluminaba el jardín de la casa familiar. Por ahora, la paz había vuelto.}
\end{corrige}

\courssub{Bilan et auto-évaluation}

\begin{aretenir}[Ce que je dois maîtriser à la fin de cette leçon]
\begin{itemize}
\item \textbf{Lexique :} \esp{viudez, viudo/a, luto, huérfano/a, herencia, heredero/a, testamento, notario, difunto, reparto, rencor, reconciliarse, legítima}.
\item \textbf{Grammaire :} Distinction Indéfini / Imparfait dans le récit ; Subjonctif présent pour l'exhortation (\esp{hablemos, busquemos}) et la volonté (\esp{quiero que...}, \esp{exijo que...}) ; Impératif \emph{nosotros} (\esp{sentémonos}).
\item \textbf{Méthode :} Structure du commentaire composé (Intro / Développement / Conclusion) ; Analyse de citations (vocabulaire, temps, syntaxe).
\item \textbf{Compétences :} Lire/comprendre un texte narratif ; Jouer un rôle (interaction orale) ; Écrire une suite narrative / lettre formelle / article d'opinion.
\item \textbf{Culture :} Différences droit successoral Espagne (legítima) / Cameroun (coutumes / droit moderne) ; Rôle du notaire.
\end{itemize}
\end{aretenir}

\begin{exercice}[Auto-évaluation : Cochez votre niveau]
\begin{center}
\begin{tabularx}{\linewidth}{|X|c|c|c|}
\hline
\rowcolor{popblueL} \textbf{Compétence} & \textbf{\textcolor{popgreen}{✓ Acquis}} & \textbf{\textcolor{poporange}{~ En cours}} & \textbf{\textcolor{poppink}{✗ À revoir}} \\
\hline
Je définis \esp{herencia} vs \esp{heredero} sans confusion. & & & \\
\hline
Je choisis correctement Indéfini / Imparfait dans un récit. & & & \\
\hline
J'utilise le subjonctif après \esp{quiero que / exijo que / propongo que}. & & & \\
\hline
Je structure un commentaire de texte (Intro / Dev / Concl). & & & \\
\hline
Je participe à un jeu de rôle en respectant mon personnage. & & & \\
\hline
Je connais les bases de \esp{la legítima} (Espagne) et coutumes (Cameroun). & & & \\
\hline
\end{tabularx}
\end{center}
\end{exercice}

\coursec{Compréhension du texte}

Dans la section précédente, nous avons lu et découvert le texte support \esp{« La herencia de don Ramón »}, l'histoire d'une famille réunie après la mort du père pour le partage de son héritage. Nous allons maintenant vérifier et approfondir notre compréhension de ce texte : c'est exactement l'exercice que l'on vous demandera au baccalauréat dans la partie \esp{comprensión lectora}. Savoir répondre à une question de compréhension est une compétence qui s'apprend : il ne suffit pas d'avoir compris, il faut savoir le \emph{prouver} par écrit, en espagnol, avec des citations.

\courssub{Rappel du texte support}

Avant de répondre aux questions, relisons une version condensée du récit pour raviver notre mémoire.

\begin{textebox}[La herencia de don Ramón (resumen)]
Don Ramón falleció una mañana de octubre, dejando a su esposa Carmen y a sus tres hijos: Luis, el mayor, empresario en la ciudad; Ana, profesora en un pueblo cercano; y Pablo, el menor, que vivía aún en la casa familiar. Tras el funeral, la familia quedó sumida en el luto, pero pronto apareció la cuestión de la herencia.

El notario leyó el testamento: la casa, los campos de cacao y una cuenta bancaria debían repartirse entre los tres hermanos. Luis quería venderlo todo y dividir el dinero. Ana deseaba conservar la casa, llena de recuerdos de su infancia. Pablo solo pedía su parte, pues había cuidado a su padre durante su larga enfermedad y se sentía olvidado por sus hermanos.

La tensión creció durante la reunión. Luis acusó a Pablo de querer quedarse con la casa sin pagar nada; Ana lloró en silencio, recordando los domingos en familia. Entonces Carmen, la madre, se levantó y habló con calma: «Vuestro padre no trabajó toda su vida para veros pelear por unas paredes.» Propuso que Pablo viviera en la casa, que los campos se repartieran entre los tres y que el dinero sirviera para la educación de los nietos. Poco a poco, los hermanos aceptaron. Aquella noche, por primera vez desde la muerte de don Ramón, la familia cenó junta en paz.
\end{textebox}

\textbf{Glossaire :} \esp{fallecer} = décéder ; \esp{el luto} = le deuil ; \esp{el notario} = le notaire ; \esp{el testamento} = le testament ; \esp{repartir} = partager, répartir ; \esp{olvidado} = oublié ; \esp{pelear} = se disputer ; \esp{las paredes} = les murs ; \esp{los nietos} = les petits-enfants ; \esp{la viudez} = la veuvage ; \esp{el heredero} = l'héritier ; \esp{el huérfano} = l'orphelin ; \esp{el conflicto familiar} = le conflit familial.

\begin{propriete}[Lexique thématique : deuil et héritage (programme officiel)]
Voici les sept termes clés du module, à connaître parfaitement (orthographe, genre, sens) pour le Bac :
\begin{center}
\begin{tabularx}{\linewidth}{|X|X|X|}
\hline
\rowcolor{popblueL} \textbf{Palabra} & \textbf{Definición / Contexto} & \textbf{Traduction} \\
\hline
\esp{la viudez} & \esp{Estado de la persona cuyo cónyuge ha muerto.} & le veuvage \\
\hline
\esp{la herencia} & \esp{Bienes, derechos y obligaciones que se transmiten tras la muerte.} & l'héritage \\
\hline
\esp{el heredero} & \esp{Persona que tiene derecho a recibir la herencia.} & l'héritier \\
\hline
\esp{el testamento} & \esp{Acto por el que una persona dispone de sus bienes para después de su muerte.} & le testament \\
\hline
\esp{el conflicto familiar} & \esp{Disputa entre miembros de una familia por intereses o sentimientos.} & le conflit familial \\
\hline
\esp{el luto} & \esp{Duelo, pena por la muerte de un ser querido; también ropa negra.} & le deuil \\
\hline
\esp{el huérfano} & \esp{Persona que ha perdido a su padre, a su madre o a ambos.} & l'orphelin \\
\hline
\end{tabularx}
\end{center}
\end{propriete}

\courssub{Compréhension globale : les questions fondamentales}

Toute lecture méthodique commence par trois questions simples : \emph{qui ?}, \emph{où et quand ?}, \emph{de quoi parle-t-on ?} Observons d'abord, puis répondons.

\begin{exemplebox}[Compréhension globale : modèles de réponses]
\begin{itemize}
\item \esp{¿Quiénes son los personajes?} $\rightarrow$ \esp{Los personajes son don Ramón (el padre fallecido), Carmen (la viuda) y sus tres hijos: Luis, Ana y Pablo.} « Les personnages sont don Ramón (le père décédé), Carmen (la veuve) et leurs trois enfants : Luis, Ana et Pablo. »
\item \esp{¿Dónde y cuándo ocurre la escena?} $\rightarrow$ \esp{La escena ocurre en la casa familiar, después del funeral de don Ramón, una mañana de octubre.} « La scène se passe dans la maison familiale, après les funérailles de don Ramón, un matin d'octobre. »
\item \esp{¿De qué habla el texto?} $\rightarrow$ \esp{El texto habla de los problemas de una familia que debe repartirse la herencia del padre.} « Le texte parle des problèmes d'une famille qui doit se partager l'héritage du père. »
\end{itemize}
\end{exemplebox}

\begin{methode}[Répondre à une question de compréhension]
\begin{enumerate}
\item \textbf{Repérer} les mots-clés de la question (\esp{quién, qué, dónde, por qué, cómo}) et retrouver le passage du texte qui y répond.
\item \textbf{Rédiger} une phrase complète en espagnol, en reprenant les mots de la question : jamais un simple « oui » ou « non », jamais un mot isolé.
\item \textbf{Justifier} par une courte citation du texte entre guillemets, introduite par \esp{como dice el texto} ou \esp{según el texto}.
\item \textbf{Vérifier} l'accord des temps (le texte est au passé : \esp{falleció, quería, propuso}) et l'orthographe des accents.
\end{enumerate}
\end{methode}

\begin{exemplebox}[Exemples de phrases du texte traduites]
\esp{La familia quedó sumida en el luto.} = « La famille fut plongée dans le deuil. »\par
\esp{Pablo solo pedía su parte.} = « Pablo demandait seulement sa part. »\par
\esp{Ana lloró en silencio, recordando los domingos en familia.} = « Ana pleura en silence, en se souvenant des dimanches en famille. »
\end{exemplebox}

\courssub{Compréhension fine : les motivations de chaque personnage}

La compréhension fine consiste à aller au-delà des faits : il faut comprendre les \emph{intentions} et les \emph{sentiments} des personnages.

\begin{itemize}
\item \textbf{Pourquoi la famille est-elle réunie ?} Parce que le notaire doit lire le testament de don Ramón : \esp{El notario leyó el testamento.}
\item \textbf{Que veut Luis ?} Vendre la maison et les champs pour partager l'argent : \esp{Luis quería venderlo todo y dividir el dinero.} C'est l'aîné, pragmatique, éloigné de la vie familiale.
\item \textbf{Que veut Ana ?} Garder la maison pour sa valeur sentimentale : \esp{Ana deseaba conservar la casa, llena de recuerdos de su infancia.}
\item \textbf{Que veut Pablo ?} Obtenir sa part et être reconnu : il a soigné son père malade et se sent lésé : \esp{se sentía olvidado por sus hermanos.}
\item \textbf{Quel est le rôle de Carmen ?} C'est la médiatrice : veuve (\esp{viudez}) et mère, elle apaise le conflit et propose une solution équitable qui réconcilie les frères et sœurs.
\end{itemize}

\begin{popfigure}
\begin{center}
\begin{tabular}{|p{2.2cm}|p{5.2cm}|p{5.2cm}|}
\hline
\rowcolor{popblueL} \textbf{Personaje} & \textbf{Lo que quiere} & \textbf{Por qué (motivación)} \\
\hline
Luis & Vender la casa y los campos, dividir el dinero & Es empresario; piensa en lo práctico, no en los recuerdos \\
\hline
Ana & Conservar la casa familiar & La casa está llena de recuerdos de su infancia \\
\hline
Pablo & Recibir su parte y ser reconocido & Cuidó a su padre enfermo; se sentía olvidado \\
\hline
Carmen & La paz y un reparto justo & Es la madre; quiere proteger la unión de la familia \\
\hline
\end{tabular}
\end{center}
\legende{Lo que quiere cada personaje : tableau récapitulatif des motivations}
\end{popfigure}

\courssub{Questions de compréhension et Vrai/Faux justifié}

\begin{exoresolu}[Questions de compréhension type Bac]
Réponds par des phrases complètes en espagnol : 1) \esp{¿Quién ha muerto?} 2) \esp{¿Qué quieren Luis, Ana y Pablo?} 3) \esp{¿Por qué hay tensión en la familia?} 4) \esp{¿Qué hace Carmen al final?}
\tcblower
1) \esp{Ha muerto don Ramón, el padre de la familia, como indica el texto: «Don Ramón falleció una mañana de octubre».}\par
2) \esp{Luis quiere venderlo todo y dividir el dinero; Ana quiere conservar la casa por sus recuerdos; Pablo solo pide su parte porque cuidó a su padre.}\par
3) \esp{Hay tensión porque los hermanos no están de acuerdo sobre el reparto de la herencia; el texto dice: «La tensión creció durante la reunión».}\par
4) \esp{Al final, Carmen propone una solución justa: Pablo vivirá en la casa, los campos se repartirán entre los tres y el dinero servirá para la educación de los nietos. Así, la familia se reconcilia.}
\end{exoresolu}

\begin{exoresolu}[Vrai ou faux : justifie par une citation]
\esp{1) Luis quiere quedarse con la casa familiar. 2) Pablo cuidó a su padre durante su enfermedad. 3) La familia termina peleada para siempre.}
\tcblower
1) \textbf{Falso.} \esp{Luis no quiere la casa: «Luis quería venderlo todo y dividir el dinero».}\par
2) \textbf{Verdadero.} \esp{El texto afirma que Pablo «había cuidado a su padre durante su larga enfermedad».}\par
3) \textbf{Falso.} \esp{Al contrario, se reconcilian: «la familia cenó junta en paz».}
\end{exoresolu}

\begin{attention}[Au Bac : la citation est obligatoire]
Une réponse recopiée du texte sans justification, ou une réponse « Vrai » / « Faux » sans citation, vaut moins de points. La règle d'or : \textbf{toujours citer le texte entre guillemets}. Attention aussi à ne pas traduire mot à mot : \esp{¿Quién ha muerto?} se répond en espagnol, pas en français. Et méfiez-vous du faux ami \esp{pretender} (= prétendre à, chercher à obtenir), qui ne signifie pas « prétendre » au sens d'« affirmer ».
\end{attention}

\courssub{Champs lexicaux du deuil et du conflit}

Repérer un \cle{champ lexical} dans un texte, c'est retrouver tous les mots qui appartiennent à une même famille de sens. C'est une question fréquente au baccalauréat.

\begin{center}
\begin{tabularx}{\linewidth}{|X|X|X|}
\hline
\rowcolor{popblueL} \textbf{Campo lexical} & \textbf{Palabras del texto} & \textbf{Traduction} \\
\hline
El luto (le deuil) & falleció, funeral, luto, lloró, muerte & décéder, funérailles, deuil, pleurer, mort \\
\hline
La herencia (l'héritage) & herencia, testamento, notario, repartir, parte & héritage, testament, notaire, partager, part \\
\hline
El conflicto (le conflit) & tensión, acusó, pelear, olvidado & tension, accuser, se disputer, oublié \\
\hline
La reconciliación (la réconciliation) & con calma, propuso, aceptaron, en paz & calmement, proposer, accepter, en paix \\
\hline
\end{tabularx}
\end{center}

\begin{popfigure}
\begin{center}
\begin{tikzpicture}[mindmap, grow cyclic, every node/.style={concept, align=center, font=\small},
root concept/.append style={concept color=popblue, font=\small\bfseries},
level 1/.append style={level distance=3.2cm, sibling angle=90},
level 2/.append style={level distance=2.2cm, sibling angle=45, font=\footnotesize}]
\node[root concept]{El texto}
  child[concept color=poporange]{ node {El luto} child { node {fallecer} } child { node {llorar} } }
  child[concept color=popgreen]{ node {La herencia} child { node {testamento} } child { node {repartir} } }
  child[concept color=poppink]{ node {El conflicto} child { node {tensión} } child { node {acusar} } }
  child[concept color=popgold]{ node {La paz} child { node {proponer} } child { node {aceptar} } };
\end{tikzpicture}
\end{center}
\legende{Carte mentale des champs lexicaux du texte}
\end{popfigure}

\courssub{Reformulation : résumer le texte en trois phrases}

Résumer, c'est dire la même chose \emph{avec ses propres mots}, sans recopier. On garde les étapes essentielles : la mort, le conflit, la solution.

\begin{exemplebox}[Modèle de résumé en trois phrases]
1) \esp{Tras la muerte de don Ramón, su viuda y sus tres hijos se reúnen para leer el testamento.} « Après la mort de don Ramón, sa veuve et ses trois enfants se réunissent pour lire le testament. »\par
2) \esp{Los hermanos discuten porque cada uno desea algo distinto de la herencia.} « Les frères et sœurs se disputent parce que chacun désire quelque chose de différent de l'héritage. »\par
3) \esp{Gracias a la sabia propuesta de Carmen, la madre, la familia encuentra un acuerdo y recupera la paz.} « Grâce à la sage proposition de Carmen, la mère, la famille trouve un accord et retrouve la paix. »
\end{exemplebox}

\courssub{Question d'ouverture et entraînement type Bac}

La question d'ouverture invite à donner son \emph{opinion personnelle} : \esp{¿Crees que el dinero puede destruir una familia? Justifica tu respuesta.} « Penses-tu que l'argent peut détruire une famille ? Justifie ta réponse. » Pour y répondre, utilisez les formules d'opinion : \esp{creo que, pienso que, en mi opinión, me parece que} + indicatif ; \esp{no creo que} + subjonctif.

\begin{exemplebox}[Modèle de réponse d'ouverture]
\esp{En mi opinión, el dinero puede destruir una familia si los parientes solo piensan en su interés. Sin embargo, creo que el diálogo y el amor, como muestra Carmen en el texto, pueden salvar la unión familiar.} « À mon avis, l'argent peut détruire une famille si les proches ne pensent qu'à leur intérêt. Cependant, je crois que le dialogue et l'amour, comme le montre Carmen dans le texte, peuvent sauver l'union familiale. »
\end{exemplebox}

\begin{textebox}[Grille de questions type Bac]
1) Identifica a los personajes del texto y di qué relación tienen entre ellos.\par
2) ¿Qué conflicto divide a la familia? Justifica tu respuesta con una cita.\par
3) Cita dos palabras del campo lexical del luto.\par
4) ¿Cuál es tu opinión sobre el reparto que propone Carmen? ¿Te parece justo?
\end{textebox}

\begin{exercice}[À vous de jouer]
1) Réponds en espagnol : \esp{¿Por qué Pablo se siente olvidado?}\par
2) Vrai ou faux (justifie) : \esp{Ana quiere vender la casa.}\par
3) Résume le texte en trois phrases avec tes propres mots.\par
4) Donne ton opinion : \esp{¿Es importante conservar la casa familiar? ¿Por qué?}
\end{exercice}


\begin{corrige}[Exercice]
1) \esp{Pablo se siente olvidado porque cuidó a su padre durante su enfermedad y sus hermanos no le ayudaron.}\par
2) \textbf{Falso.} \esp{Ana «deseaba conservar la casa, llena de recuerdos de su infancia».}\par
3) Exemple : \esp{Don Ramón muere y deja una herencia. Sus hijos discuten sobre el reparto. Su madre Carmen propone una solución y la familia se reconcilia.}\par
4) Réponse libre, par exemple : \esp{Sí, creo que es importante porque la casa guarda la memoria y la historia de la familia.}
\end{corrige}

\coursec{Méthode du commentaire de texte}

Au baccalauréat, l'épreuve de \esp{comprensión lectora} peut exiger un commentaire structuré (introduction, développement, conclusion) et non plus seulement des réponses à des questions. Voici la méthode rigoureuse à appliquer à tout texte narratif ou argumentatif.

\begin{methode}[Commentaire de texte : 5 étapes]
\begin{enumerate}
\item \textbf{Lecture active (10 min)} : Lire 2 fois. Souligner : personnages, temps, lieu, verbes d'action, connecteurs logiques, champs lexicaux. Noter la structure (début, milieu, fin).
\item \textbf{Analyse du contenu (quoi ?)} : Résumer l'intrigue en 2 phrases. Identifier le thème central (ici : héritage et réconciliation) et le problème (conflit fratrie).
\item \textbf{Analyse de la forme (comment ?)} : Type de texte (narratif), temps verbaux (passé narratif : \esp{pretérito indefinido} actions, \esp{imperfecto} descriptions/habitudes), point de vue (narrateur omniscient), figures de style (dialogue direct, métaphore \esp{« por unas paredes »}).
\item \textbf{Construction du plan (dialectique)} : \textbf{Thèse} : La mort déclenche la cupidité/le conflit. \textbf{Antithèse} : La figure maternelle (Carmen) incarne la sagesse/la loi du cœur. \textbf{Synthèse} : Le compromis restaure l'unité familiale ; l'héritage matériel devient héritage moral.
\item \textbf{Rédaction soignée (20 min)} : \textbf{Introduction} : Accroche, présentation du texte (titre, auteur fictif, thème), problématique, annonce du plan. \textbf{Développement} : 2-3 paragraphes arguments + citations. \textbf{Conclusion} : Bilan, réponse à la problématique, ouverture (ex: rôle de la femme/mère dans la société hispanique).
\end{enumerate}
\end{methode}

\begin{attention}[Pièges à éviter au commentaire]
\begin{itemize}
\item \textbf{Le résumé pur :} Ne pas raconter l'histoire, mais \emph{analyser} comment le texte construit son sens.
\item \textbf{L'absence de citations :} Toute affirmation sur le texte doit être étayée par un exemple précis (mot, phrase).
\item \textbf{L'oubli du contexte culturel :} Ici, \esp{la viudez}, le rôle de la \esp{madre} comme pilier familial, la \esp{casa familiar} comme patrimoine symbolique (terre, cacao, racines).
\item \textbf{La langue :} Écrire en espagnol correct, phrases complètes, connecteurs (\esp{en primer lugar, por tanto, sin embargo, en conclusión}).
\end{itemize}
\end{attention}

\coursec{Jeu de rôle : la conversation familiale}

\courssub{Consignes et mise en situation}

\begin{experience}[Jeu de rôle : « La reunión de herederos »]
\textbf{Objectif :} Réinvestir le lexique (héritage, deuil, conflit, réconciliation) et les structures (donner son opinion, exprimer un accord/désaccord, proposer, négocier) dans une interaction orale authentique.

\textbf{Mise en scène :} Salon de la maison familiale, après la lecture du testament par le notaire (sorti). Les 4 personnages sont présents.

\textbf{Rôles (4 élèves + 1 observateur) :}
\begin{enumerate}
\item \textbf{Luis (l'aîné, pragmatique) :} Veut vendre vite. Arguments : « \esp{No tiene sentido mantener la casa vacía}, \esp{el dinero nos sirve para invertir} ». Utilise \esp{creo que debemos vender}, \esp{es lo más lógico}.
\item \textbf{Ana (la sentimentale) :} Veut garder la maison. Arguments : « \esp{Esta casa es nuestra historia}, \esp{no se puede vender la memoria} ». Utilise \esp{me duele pensar en vender}, \esp{propongo conservar}.
\item \textbf{Pablo (le cadet, soignant) :} Veut justice/reconnaissance. Arguments : « \esp{Yo me quedé cuidando a papá}, \esp{me corresponde la casa o mi parte justa} ». Utilise \esp{exijo mi derecho}, \esp{me siento marginado}.
\item \textbf{Carmen (la mère, médiatrice) :} Cherche la paix. Arguments : « \esp{No quiero gritos}, \esp{vuestro padre quería la unión} ». Utilise \esp{propongo un acuerdo}, \esp{escuchadme}.
\item \textbf{Observateur :} Note le vocabulaire utilisé, les temps (subjonctif après \esp{quiero que, propongo que}), la politesse (\esp{usted} vs \esp{tú} selon région), la fluidité.
\end{enumerate}

\textbf{Déroulement (10 min prep + 5 min jeu + 5 min débrief) :}
\begin{enumerate}
\item Préparation individuelle : chaque élève écrit 3-4 phrases-clés pour son rôle (utiliser fiche lexique).
\item Jeu libre : discussion animée, objectif = trouver un accord (comme dans le texte ou différent).
\item Débrief : l'observateur donne son retour ; le professeur corrige 2-3 erreurs majeures (ex: \esp{*quiero que vendes} $\rightarrow$ \esp{quiero que vendas}).
\end{enumerate}
\end{experience}

\begin{aretenir}[Expressions utiles pour la négociation familiale]
\begin{center}
\begin{tabularx}{\linewidth}{|X|X|}
\hline
\rowcolor{popblueL} \textbf{Fonction} & \textbf{Exemples en espagnol} \\
\hline
Donner son avis & \esp{En mi opinión... ; Pienso que... ; Me parece que...} \\
\hline
Exprimer un désaccord & \esp{No estoy de acuerdo ; Eso no es justo ; Pero...} \\
\hline
Proposer une solution & \esp{Propongo que + subj. ; ¿Y si... + imperf. subj.? ; Podríamos + inf.} \\
\hline
Insister / Négocier & \esp{Te ruego que... ; A cambio, yo... ; Es mi última oferta.} \\
\hline
Accepter / Conclure & \esp{De acuerdo ; Me parece bien ; Trato hecho ; Firmemos.} \\
\hline
\end{tabularx}
\end{center}
\end{aretenir}

\coursec{Production guidée et bilan}

\courssub{Expression écrite : rédaction d'un article d'opinion}

\begin{exercice}[Sujet de rédaction type Bac (150-180 mots)]
\textbf{Contexte :} Vous êtes journaliste pour le journal des jeunes \esp{« Voz Joven »} au Cameroun. Écrivez un article d'opinion intitulé : \esp{« La herencia: ¿fuente de conflicto o de unión familiar? »}.

\textbf{Contraintes :}
\begin{itemize}
\item Structure : Titre, introduction (accroche + thèse), 2 arguments illustrés (ex: texte lu, exemple camerounais), conclusion (ouverture).
\item Vocabulaire imposé : \esp{viudez, heredero, testamento, luto, repartir, reconciliación}.
\item Grammaire : Au moins 3 subjonctifs (ex: \esp{es importante que...}, \esp{quiero que...}, \esp{propongo que...}), connecteurs logiques.
\item Registre : Formel / Standard.
\end{itemize}
\end{exercice}


\begin{corrige}[Exercice - Modèle de rédaction]
\textbf{La herencia: ¿fuente de conflicto o de unión familiar?}

La muerte de un padre siempre sumbe a la familia en el luto, pero a menudo la viudez y el reparto de la herencia sacan a la luz tensiones ocultas. En mi opinión, la herencia no tiene por qué dividir; todo depende del diálogo y de la justicia.

En primer lugar, la ausencia de testamento claro genera disputas entre los herederos. Como vimos en el texto de don Ramón, Luis quería vender los campos de cacao, Ana conservar la casa y Pablo exigía reconocimiento por su cuidado. El dinero y los bienes materiales ciegan el afecto. En Camerún, como en España o América Latina, los conflictos por la tierra o la casa familiar son frecuentes y rompen la unidad.

Sin embargo, la reconciliación es posible si prima el bien común. Carmen, la madre, propone un reparto equitativo: la casa para Pablo, los campos para los tres, el dinero para los nietos. Es fundamental que los hermanos acepten ceder. Propongo que las familias redacten testamentos en vida y hablen de la muerte sin tabú.

En conclusión, la herencia será fuente de unión si el amor y el respeto guían el reparto. ¿Y tú, qué harías para evitar la guerra entre hermanos?
\end{corrige}

\courssub{Bilan de la leçon}

\begin{aretenir}[Ce qu'il faut retenir de la leçon E05]
\begin{itemize}
\item \textbf{Texte support :} \esp{« La herencia de don Ramón »} (récit narratif, passé \esp{indefinido/imperfecto}, dialogue direct, dénouement pacifié).
\item \textbf{Lexique clé (7 mots) :} \esp{viudez, herencia, heredero, testamento, conflicto familiar, luto, huérfano}. Orthographe, genre, emploi en contexte.
\item \textbf{Compréhension :} Réponse en phrase complète + citation obligatoire (\esp{« ... »}). Distinction global / fin / ouverture.
\item \textbf{Champs lexicaux :} Identifier 4 champs (luto, herencia, conflicto, paz) et citer les mots du texte.
\item \textbf{Grammaire réinvestie :} Subjonctif après verbes de volonté/émotion (\esp{propongo que viva, quiero que repartan}), passé narratif, \esp{ser/estar} (\esp{está llena de recuerdos} vs \esp{es empresario}).
\item \textbf{Compétence orale :} Jeu de rôle argumentatif (négociation, médiation) avec fiches de rôle.
\item \textbf{Compétence écrite :} Commentaire structuré (Intro/Développement/Conclusion) et Article d'opinion (thèse/arguments/exemples/conclusion).
\end{itemize}
\end{aretenir}

\coursec{Lexique thématique : deuil et héritage}

Après avoir analysé le texte \emph{« La herencia de don Ramón »} et vérifié votre compréhension globale, nous allons maintenant fixer le vocabulaire essentiel pour parler de la mort, du deuil, de l'héritage, des conflits familiaux et des émotions qui les accompagnent. Ce lexique vous servira non seulement pour le commentaire de texte, mais aussi pour le jeu de rôle final et pour l'expression écrite (essai, lettre, article). Nous procéderons par \cle{champs lexicaux} : observation d'exemples authentiques, énoncé rigoureux, comparaison avec le français, exercices de fixation et pièges à éviter.

\courssub{Champ 1 — Le décès et le deuil : \emph{morir, fallecer, el luto, la viudez…}}

\begin{textebox}[Dialogue au cimetière : le deuil dans une famille camerounaise]
-- ¿Has visto a doña Amina? Está de luto desde que falleció su esposo.
-- Sí, la viudez le ha cambiado la vida. Lleva negro todos los días.
-- Es normal. Perder al compañero de toda la vida es un dolor inmenso.
-- Sus hijos, ahora huérfanos de padre, la apoyan mucho.
-- El funeral fue ayer. El entierro, muy emotivo. Todo el barrio de Mvog-Ada vino.
-- Que descanse en paz. La muerte de don Samuel nos deja a todos tristes.
\end{textebox}

\begin{center}
\begin{tabularx}{\linewidth}{|X|X|X|}
\hline
\rowcolor{popblueL}
\textbf{Espagnol} & \textbf{Français} & \textbf{Remarque} \\
\hline
\esp{morir (muerto)} & mourir (mort) & Verbe irrégulier (o $\to$ ue au présent indicatif/subjonctif : \emph{muero, mueras} ; participe passé \emph{muerto}). \\
\esp{fallecer} & décéder, trépasser & Registre \textbf{soutenu, respectueux, administratif}. Préféré dans les faire-part, les condoléances, la presse. \\
\hline
\esp{la muerte} & la mort & Nom féminin. \emph{La muerte de don Samuel}. \\
\esp{el luto} & le deuil & \emph{Estar de luto} = être en deuil (porter le deuil). \\
\esp{el funeral} & les funérailles (cérémonie) & Souvent au pluriel : \emph{los funerales}. \\
\esp{el entierro} & l'enterrement & Action d'enterrer. \\
\hline
\esp{la viudez} & le veuvage & Situation de la personne dont le conjoint est mort. \\
\esp{viudo / viuda} & veuf / veuve & Accord en genre. \emph{Don Samuel era viudo? No, doña Amina es viuda.} \\
\esp{huérfano / huérfana} & orphelin / orpheline & \emph{Huérfano de padre / de madre / de ambos}. \\
\hline
\end{tabularx}
\end{center}

\begin{attention}[Faux ami et registre : \emph{morir} vs \emph{fallecer}]
En français, « mourir » est neutre. En espagnol, \esp{morir} est le verbe courant, mais \esp{fallecer} est \textbf{obligatoire} dans le langage formel, journalistique, notarial et pour présenter ses condoléances.
\begin{itemize}
\item \esp{Mi abuelo murió anoche.} (familier, entre proches)
\item \esp{El Sr. López falleció ayer a los 89 años.} (formel, faire-part, journal)
\item \esp{¿Cuándo falleció su padre?} (politesse, respect)
\end{itemize}
Ne dites jamais \esp{« Su padre murió »} à une personne en deuil : c'est perçu comme brutal, voire grossier.
\end{attention}

\begin{definition}[Echar de menos — « manquer à quelqu'un »]
\esp{Echar de menos} est une locution verbale \textbf{intraduisible mot à mot}. Elle signifie « ressentir l'absence de quelqu'un/quelque chose », « regretter », « avoir la nostalgie ».
\begin{itemize}
\item Sujet = celui qui ressent l'absence (celui à qui « ça manque »).
\item Complément d'objet direct (souvent précédé de \emph{a} si c'est une personne) = la personne/la chose qui manque.
\end{itemize}
\esp{Echo de menos a mi padre.} = « Mon père me manque. » (Litt. : « Je jette de moins mon père. ») \\
\esp{¿Echas de menos a tu abuela?} = « Ta grand-mère te manque-t-elle ? »
\end{definition}

\begin{exemplebox}[Exemples variés — décès et deuil]
\esp{Don Ramón murió en paz a los 82 años.} — « Don Ramón est mort paisiblement à 82 ans. » \\
\esp{La noticia de su fallecimiento conmocionó al pueblo.} — « La nouvelle de son décès a bouleversé le village. » \\
\esp{Doña Carmen lleva tres años de luto.} — « Doña Carmen porte le deuil depuis trois ans. » \\
\esp{Los funerales serán mañana a las diez en la catedral.} — « Les funérailles auront lieu demain à dix heures à la cathédrale. » \\
\esp{Tras el entierro, la familia se reunió en la casa familiar.} — « Après l'enterrement, la famille s'est réunie dans la maison familiale. » \\
\esp{La viudez no es fácil, pero doña Amina es fuerte.} — « Le veuvage n'est pas facile, mais doña Amina est forte. » \\
\esp{Los niños, huérfanos de padre, necesitan apoyo psicológico.} — « Les enfants, orphelins de père, ont besoin d'un soutien psychologique. » \\
\esp{Te echo mucho de menos, abuelo.} — « Tu me manques beaucoup, grand-père. »
\end{exemplebox}

\courssub{Champ 2 — L'héritage : \emph{la herencia, el testamento, el notario…}}

\begin{textebox}[Extrait d'une conversation chez le notaire]
-- Buenas tardes, don Notario. Vengo por la herencia de mi padre, don Ramón.
-- Buenas. Tenemos el testamento aquí. Él dejó la casa en herencia a sus tres hijos.
-- ¿Y los bienes muebles? ¿El coche, las joyas, el dinero del banco?
-- Todo está en el testamento. La sucesión se abre hoy. Ustedes son los herederos legítimos.
-- Mi hermano mayor se niega a firmar. Dice que el reparto no es justo.
-- Entonces habrá conflicto. El notario no puede repartir si no hay acuerdo.
-- ¿Qué pasa si no llegamos a un acuerdo?
-- Juez, abogados, años de juicio... Lo mejor es reconciliarse y hacer las paces.
\end{textebox}

\begin{center}
\begin{tabularx}{\linewidth}{|X|X|X|}
\hline
\rowcolor{poporangeL}
\textbf{Espagnol} & \textbf{Français} & \textbf{Remarque} \\
\hline
\esp{la herencia} & l'héritage & Ensemble des biens transmis. \\
\esp{heredar} & hériter & Verbe régulier en \emph{-ar}. \esp{Heredé la casa de mi abuelo.} \\
\esp{el heredero / la heredera} & l'héritier / l'héritière & \emph{Herederos legítimos} = héritiers légaux (réserve héréditaire). \\
\esp{el testamento} & le testament & Document écrit devant notaire. \emph{Hacer testamento}. \\
\esp{dejar en herencia} & léguer, laisser en héritage & \esp{Dejó la finca en herencia a su hija menor.} \\
\esp{la sucesión} & la succession & Procédure juridique d'ouverture de l'héritage. \\
\esp{el notario} & le notaire & Officier public qui reçoit les actes. \\
\esp{los bienes} & les biens (meubles/immeubles) & \emph{Bienes muebles} (voiture, bijoux, argent) / \emph{bienes inmuebles} (terrain, maison). \\
\esp{la propiedad} & la propriété & Terrain, maison, appartement. \emph{La propiedad está a nombre de los tres.} \\
\esp{el reparto} & le partage & Action de répartir. \esp{El reparto de la herencia.} \\
\hline
\end{tabularx}
\end{center}

\begin{propriete}[Constructions clés du vocabulaire successoral]
\begin{enumerate}
\item \textbf{Heredar + objet} : \esp{Heredé la casa.} (J'ai hérité la maison.)
\item \textbf{Dejar en herencia + objet + a + bénéficiaire} : \esp{Dejó el coche en herencia a su nieto.} (Il a laissé la voiture en héritage à son petit-fils.)
\item \textbf{Ser heredero de} : \esp{Somos herederos de don Ramón.} (Nous sommes héritiers de don Ramón.)
\item \textbf{Abrir la sucesión} : expression figée = ouvrir la succession.
\item \textbf{Repartir la herencia / hacer el reparto} : partager l'héritage / faire le partage.
\end{enumerate}
\end{propriete}

\begin{exemplebox}[Exemples — héritage et notaire]
\esp{Mi tío hizo testamento ante notario el mes pasado.} — « Mon oncle a fait testament chez le notaire le mois dernier. » \\
\esp{La herencia incluye tres fincas de cacao cerca de Ebolowa.} — « L'héritage comprend trois plantations de cacao près d'Ebolowa. » \\
\esp{Los bienes muebles se vendieron en subasta.} — « Les biens meubles ont été vendus aux enchères. » \\
\esp{El notario leyó el testamento en voz alta.} — « Le notaire a lu le testament à haute voix. » \\
\esp{La sucesión se abrió oficialmente ayer.} — « La succession a été ouverte officiellement hier. » \\
\esp{¿Ya hiciste el reparto de la herencia?} — « As-tu déjà fait le partage de l'héritage ? »
\end{exemplebox}

\courssub{Champ 3 — Le conflit et la réconciliation : \emph{el conflicto, la disputa, reconciliarse…}}

\begin{textebox}[Dispute entre frères pour l'héritage]
-- ¡Esto es un robo! ¡Tú te quedas con la mejor parte!
-- No grites, hermano. El testamento es claro: la casa para los tres.
-- ¿Y el dinero? ¿Dónde está el dinero que papá guardaba en la caja fuerte?
-- No hay dinero. Las deudas del hospital se lo comieron todo.
-- ¡Mientes! Yo no me voy a callar. Voy a demandar.
-- ¿Demandar? ¿A tu propia familia? Piensa en mamá, por favor.
-- Mamá está destrozada. Si seguimos así, la matamos de pena.
-- Tienes razón. Basta. \textbf{Hagamos las paces}. Perdonémonos.
-- Vale. \textbf{Lleguemos a un acuerdo}. La casa se vende y repartimos lo que sobre.
-- Trato hecho. \textbf{Nos reconciliamos}. Por papá, por mamá, por nosotros.
\end{textebox}

\begin{center}
\begin{tabularx}{\linewidth}{|X|X|X|}
\hline
\rowcolor{poppurpleL}
\textbf{Espagnol} & \textbf{Français} & \textbf{Remarque} \\
\hline
\esp{el conflicto familiar} & le conflit familial & \\
\esp{la disputa} & la dispute, le litige & Plus formel que \emph{pelea}. \\
\esp{pelearse (con alguien)} & se disputer (avec qqn) & Verbe pronominal. \esp{Se pelean por la herencia.} \\
\esp{negarse a + infinitif} & refuser de & \esp{Se niega a firmar.} (Il refuse de signer.) \\
\esp{repartir} & répartir, partager & Verbe irrégulier (e $\to$ ie : \emph{reparto, repartes, reparte}…). \\
\esp{el reparto} & le partage (action/résultat) & Nom masculin. \\
\esp{llegar a un acuerdo} & parvenir à un accord & Locution verbale figée. \\
\esp{reconciliarse (con alguien)} & se réconcilier (avec qqn) & Verbe pronominal. \esp{Se reconciliaron ayer.} \\
\esp{perdonar} & pardonner & \esp{Perdoña a tu hermano.} \\
\esp{hacer las paces} & se réconcilier, faire la paix & \textbf{Expression idiomatique \#1} pour « se réconcilier ». Voir \emph{Le saviez-vous ?}. \\
\hline
\end{tabularx}
\end{center}

\begin{savaistu}[Hacer las paces — l'expression de la réconciliation quotidienne]
\esp{Hacer las paces} (litt. « faire les paix ») est l'expression la plus courante en Espagne et en Amérique latine pour dire « se réconcilier » après une dispute, entre amis, en couple, en famille. Elle s'emploie souvent au présent de l'indicatif pour proposer la réconciliation : \esp{« ¿Hacemos las paces? »} = « On fait la paix ? ». Au passé : \esp{« Hicieron las paces »} = « Ils se sont réconciliés ». Synonyme plus formel : \esp{reconciliarse}.
\end{savaistu}

\begin{exemplebox}[Exemples — conflit et réconciliation]
\esp{El conflicto familiar estalló tras la lectura del testamento.} — « Le conflit familial a éclaté après la lecture du testament. » \\
\esp{Los hermanos se pelean por la herencia.} — « Les frères se disputent pour l'héritage. » \\
\esp{Mi prima se niega a vender la casa familiar.} — « Ma cousine refuse de vendre la maison familiale. » \\
\esp{El abogado propuso un reparto equitativo.} — « L'avocat a proposé un partage équitable. » \\
\esp{Tras meses de juicio, llegaron a un acuerdo.} — « Après des mois de procès, ils sont parvenus à un accord. » \\
\esp{Por fin se reconciliaron con su padre.} — « Ils se sont enfin réconciliés avec leur père. » \\
\esp{Te perdono, pero no olvido.} — « Je te pardonne, mais j'oublie pas. » \\
\esp{¡Hagamos las paces! La vida es muy corta.} — « Faisons la paix ! La vie est trop courte. »
\end{exemplebox}

\courssub{Champ 4 — Les émotions : \emph{la tristeza, el dolor, la rabia, echar de menos…}}

\begin{center}
\begin{tabularx}{\linewidth}{|X|X|X|}
\hline
\rowcolor{popgreenL}
\textbf{Espagnol} & \textbf{Français} & \textbf{Remarque} \\
\hline
\esp{la tristeza} & la tristesse & Nom féminin. \esp{Sentir mucha tristeza.} \\
\esp{el dolor} & la douleur, le chagrin & Physique ou moral. \esp{El dolor de la pérdida.} \\
\esp{la rabia} & la colère, la rage & Plus fort que \emph{enfado}. \esp{Me da rabia su actitud.} \\
\esp{la alegría} & la joie & \esp{La alegría de reencontrarse.} \\
\esp{sentir} & ressentir, éprouver & Verbe à diphtongue (e $\to$ ie : \emph{siento, sientes, siente}…). \\
\esp{echar de menos} & manquer (à qqn), regretter & Voir \textbf{Définition} ci-dessus. Construction : \emph{echo de menos a mi madre}. \\
\hline
\end{tabularx}
\end{center}

\begin{attention}[Pièges francophones — \emph{sentir} vs \emph{sentirse} / \emph{echar de menos}]
\begin{enumerate}
\item \textbf{\emph{Sentir} (transitif) + nom} = ressentir une émotion. \esp{Siento mucha tristeza.} (Je ressens beaucoup de tristesse.) \\
\textbf{\emph{Sentirse} (pronominal) + adjectif} = se sentir (état). \esp{Me siento triste.} (Je me sens triste.)
\item \textbf{\emph{Echar de menos}} : le sujet est celui qui \textbf{subit} l'absence. \esp{Yo echo de menos a mi padre} = « Mon père me manque » (et non *« Je manque à mon père »*). Pour dire « Je manque à mon père », on inverse les rôles : \esp{Mi padre me echa de menos a mí.} (Mon père me manque à moi). \textbf{Attention au leïsme} : en Espagne, on entend souvent \esp{« Le echo de menos »} (pour un homme), mais la norme standard (et l'usage latino-américain/camerounais) exige \esp{« Lo echo de menos »} (COD masculin).
\item \textbf{\emph{Rabia}} vs \emph{enfado} : \emph{enfado} = agacement, colère passagère ; \emph{rabia} = colère intense, souvent mêlée d'impuissance ou d'injustice.
\end{enumerate}
\end{attention}

\begin{exemplebox}[Exemples — émotions du deuil et de l'héritage]
\esp{Siento un dolor profundo en el alma.} — « Je ressens une douleur profonde dans l'âme. » \\
\esp{La tristeza nos invade a todos.} — « La tristesse nous envahit tous. » \\
\esp{Me da rabia que mi hermano actúe así.} — « Ça me met en colère que mon frère agisse ainsi. » \\
\esp{¿Echas de menos los domingos en familia?} — « Les dimanches en famille te manquent-ils ? » \\
\esp{La alegría de la reconciliación borró la rabia.} — « La joie de la réconciliation a effacé la colère. » \\
\esp{Me siento solo desde que murió mi esposa.} — « Je me sens seul depuis que ma femme est morte. »
\end{exemplebox}

\courssub{Prononciation : points de vigilance}

\begin{center}
\begin{tabular}{|l|l|l|}
\hline
\rowcolor{poppinkL}
\textbf{Mot} & \textbf{API (approx.)} & \textbf{Observation} \\
\hline
\esp{viudo} & [\textipa{ˈbjuðo}] & \emph{d} intervocalique = \textipa{[ð]} (comme \emph{th} anglais \emph{this}), pas \textipa{[d]}. \\
\esp{viuda} & [\textipa{ˈbjuða}] & Même chose. Attention au genre : \emph{viud\textbf{o}} (masc.) / \emph{viud\textbf{a}} (fém.). \\
\esp{huérfano} & [\textipa{ˈweɾfano}] & \textbf{H muet} ! Ne pas aspirer. \emph{huérfano} = \emph{uérfano}. Accent sur \emph{é}. \\
\esp{herencia} & [\textipa{eˈɾenθja}] / [\textipa{eˈɾensja}] & \textbf{H muet}. \emph{c} devant \emph{i/e} = \textipa{[θ]} (Espagne) ou \textipa{[s]} (Amérique). \\
\esp{fallecer} & [\textipa{faˈʎeθeɾ}] / [\textipa{faˈʝeseɾ}] & \emph{ll} = \textipa{[ʎ]} (Espagne) ou \textipa{[ʝ]} (Amérique, \emph{yeísmo}). \\
\esp{testamento} & [\textipa{testaˈmento}] & Accent sur l'antépénultième (proparoxyton) $\to$ accent écrit obligatoire. \\
\hline
\end{tabular}
\end{center}

\courssub{Exercice résolu : fixation du lexique}

\begin{exoresolu}[Compléter le texte avec le vocabulaire adéquat]
\textbf{Énoncé :} Complétez le paragraphe suivant en choisissant le mot convenant parmi la liste : \emph{herencia, testamento, viudez, luto, herederos, falleció, reparto, reconciliaron, echan de menos, rabia}.

« Doña Carmen lleva dos años de \_\_\_\_\_\_ depuis que son esposo don Ramón \_\_\_\_\_\_. La \_\_\_\_\_\_ a été dure pour elle. Ses trois enfants, ahora \_\_\_\_\_\_, se reunieron chez le notaire pour lire le \_\_\_\_\_\_. La \_\_\_\_\_\_ comprend la maison de Yaoundé et une plantation de cacao. L'aîné, en colère, a provoqué un conflit : il \_\_\_\_\_\_ de signer le \_\_\_\_\_\_. Sa sœur, pleine de \_\_\_\_\_\_, a menacé de porter l'affaire au tribunal. Heureusement, ils ont fini par se parler et \_\_\_\_\_\_. Aujourd'hui, ils \_\_\_\_\_\_ beaucoup leur père, mais ils sourient en se rappelant ses blagues. »

\tcblower
\textbf{Solution détaillée :}
\begin{enumerate}
\item \esp{luto} (elle porte le deuil depuis deux ans) — \emph{estar de luto}.
\item \esp{falleció} (registre respectueux pour annoncer le décès).
\item \esp{viudez} (situation de veuve).
\item \esp{herederos} (héritiers légitimes).
\item \esp{testamento} (document lu chez le notaire).
\item \esp{herencia} (ensemble des biens).
\item \esp{se niega} (refuse de signer — \emph{negarse a} + infinitif).
\item \esp{reparto} (partage de l'héritage).
\item \esp{rabia} (colère intense).
\item \esp{se reconciliaron} / \esp{hicieron las paces} (réconciliation).
\item \esp{lo echan de menos} (sujet \emph{ellos}, COD \emph{lo} pour \emph{su padre}, norme standard).
\end{enumerate}
\end{exoresolu}

\courssub{Texte d'entraînement : « El último adiós en Mvog-Mbi »}

\begin{textebox}[Texte original — Contexte camerounais : deuil, héritage et réconciliation]
En el barrio Mvog-Mbi de Yaoundé, la campana de la parroquia de San José tañe lento. Es el funeral de don Patrice, antiguo jefe de estación de ferrocarril, fallecido a los 78 años tras una larga enfermedad. Doña Hélène, su viuda, vestida de negro riguroso, recibe los pésames con la cabeza alta, aunque sus ojos delatan una tristeza infinita. Sus cuatro hijos —dos varones y dos mujeres— están a su lado, pero la tensión se nota en los gestos, en los silencios, en las miradas que se cruzan.

-- ¿Has visto el testamento? —susurra el mayor, Pierre, a su hermana Marie.
-- Sí. La casa de Mvog-Mbi y la finca de cacao en Ebolowa para los cuatro a partes iguales. El coche y la cuenta bancaria para mamá.
-- ¿Y las joyas de la abuela? ¿Y el terreno de Nkolbisson que papá compró en 1995?
-- No aparecen en el testamento. Papá dijo que eso se arreglaría « entre hermanos ».
-- ¡Eso es una trampa! Pierre, no voy a permitir que te quedes con todo. Me niego a firmar el reparto si no hay claridad.

La voz de Marie sube. Doña Hélène se gira, cansada:
-- ¡Basta! Vuestro padre no querría esto. ¿Por la tierra y el dinero vais a romper la familia? ¡Haced las paces, por favor! Yo no quiero herencia, quiero paz.

El notario, don Armand, interviene con calma:
-- Señores, la sucesión está abierta. Si no hay acuerdo, iremos a juicio. Años, dinero, abogados... ¿Merece la pena?

Un silencio pesado. Luego, el menor, Jean, toma la mano de Pierre:
-- Hermano, papá nos enseñó a compartir. ¿Por qué no vendemos la casa de Yaoundé, repartimos el dinero y cada uno construye su vida? Mamá se queda con la finca de Ebolowa.

Pierre baja la mirada, avergonzado. Marie suspira, seca una lágrima:
-- Vale. Vendamos. Repartamos. \textbf{Nos reconciliamos}. Por mamá. Por papá.

Esa tarde, bajo el mango del patio familiar, firman el acuerdo. El notario sonríe. Doña Hélène, por primera vez en meses, esboza una sonrisa. La campana deja de tañer. El luto sigue, pero la guerra ha terminado.
\end{textebox}

\begin{center}
\begin{tabularx}{\linewidth}{|X|X|}
\hline
\rowcolor{popgoldL}
\textbf{Palabra / Expresión} & \textbf{Signification / Glossaire} \\
\hline
\esp{tañe} & sonne (cloche) — \emph{tañer} (verbe irrégulier, présent \emph{taño, tañes, tañe}…). \\
\esp{antiguo jefe de estación} & ancien chef de gare. \\
\esp{viuda} & veuve. \\
\esp{riguroso} & strict, rigoureux (noir rigoureux = noir complet). \\
\esp{delatan} & trahissent, révèlent — \emph{delatar}. \\
\esp{se nota} & se remarque, c'est visible. \\
\esp{susurra} & chuchote — \emph{susurrar}. \\
\esp{finca} & exploitation agricole, plantation (terme courant en zone cacayère). \\
\esp{a partes iguales} & à parts égales. \\
\esp{trampa} & piège, combine. \\
\esp{me niego a firmar} & je refuse de signer. \\
\esp{romper la familia} & briser la famille. \\
\esp{merece la pena} & ça en vaut la peine / ça vaut le coup. \\
\esp{avergonzado} & honteux, honteuse. \\
\esp{esboza una sonrisa} & esquisse un sourire. \\
\hline
\end{tabularx}
\end{center}

\begin{exercice}[Questions de compréhension et de lexique]
\begin{enumerate}
\item Relevez dans le texte cinq mots ou expressions du champ lexical de \textbf{l'héritage} et cinq du champ de \textbf{la réconciliation}.
\item Pourquoi Marie refuse-t-elle de signer le partage au début ?
\item Quel rôle joue doña Hélène dans la résolution du conflit ?
\item Traduisez en français : \esp{« ¿Merece la pena? »} et \esp{« Nos reconciliamos ».}
\item Transformez la phrase \esp{« Me niego a firmar »} à la 3\textsuperscript{e} personne du singulier (sujet : \emph{él}).
\item Mettez au féminin : \esp{el heredero, el viudo, el huérfano, el notario}.
\end{enumerate}
\end{exercice}

\courssub{Carte mentale : les quatre piliers du lexique}

\begin{popfigure}
\begin{tikzpicture}[
  mindmap,
  grow cyclic,
  every node/.style={concept, execute at begin node=\hskip0pt},
  root concept/.append style={concept color=popblue, minimum size=3.5cm, text=white, font=\large\bfseries},
  level 1 concept/.append style={sibling angle=90, minimum size=2.8cm, text=white, font=\bfseries},
  level 2 concept/.append style={sibling angle=45, minimum size=2.2cm, text=white, font=\small},
  level 3 concept/.append style={sibling angle=30, minimum size=1.8cm, text=white, font=\scriptsize},
  concept color=popblue,
  text=white
]
\node[root concept, align=center]{Familia \\ Problemas}
  child[concept color=poppurple] {
    node {El luto}
    child { node {morir / fallecer} }
    child { node {la muerte} }
    child { node {el funeral / entierro} }
    child { node {el luto / estar de luto} }
    child { node {la viudez} }
    child { node {viudo / viuda} }
    child { node {huérfano / a} }
    child { node {echar de menos} }
  }
  child[concept color=poporange] {
    node {La herencia}
    child { node {la herencia} }
    child { node {heredar} }
    child { node {heredero / a} }
    child { node {el testamento} }
    child { node {dejar en herencia} }
    child { node {la sucesión} }
    child { node {el notario} }
    child { node {los bienes / propiedad} }
    child { node {el reparto} }
  }
  child[concept color=poppink] {
    node {El conflicto}
    child { node {conflicto familiar} }
    child { node {la disputa} }
    child { node {pelearse} }
    child { node {negarse a} }
    child { node {repartir} }
    child { node {llegar a un acuerdo} }
    child { node {reconciliarse} }
    child { node {perdonar} }
    child { node {hacer las paces} }
  }
  child[concept color=popgreen] {
    node {Las emociones}
    child { node {la tristeza} }
    child { node {el dolor} }
    child { node {la rabia} }
    child { node {la alegría} }
    child { node {sentir} }
    child { node {sentirse} }
    child { node {echar de menos} }
  };
\end{tikzpicture}
\legende{Carte mentale des quatre champs lexicaux : deuil (violet), héritage (orange), conflit/réconciliation (rose), émotions (vert).}
\end{popfigure}

\courssub{Méthode pour mémoriser et réemployer ce lexique}

\begin{methode}[Apprendre le vocabulaire thématique par champs lexicaux]
\begin{enumerate}
\item \textbf{Classer} : Regroupez les mots par thème (comme les 4 champs ci-dessus). Utilisez des couleurs différentes sur vos fiches.
\item \textbf{Contextualiser} : Pour chaque mot, écrivez \textbf{une phrase personnelle} qui ancre le sens dans votre réalité (Cameroun, famille, quartier). Ex. : \esp{« Mi abuelo falleció en 2020, y mi abuela está de luto desde entonces. »}
\item \textbf{Associer} : Créez des paires \emph{nom + verbe} ou \emph{verbe + préposition} : \esp{hacer testamento, abrir la sucesión, llegar a un acuerdo, echar de menos a alguien}.
\item \textbf{Prononcer} : Lisez chaque mot à voix haute en respectant l'accent tonique et les particularités (h muet, \emph{ll}, \emph{d} intervocalique). Enregistrez-vous.
\item \textbf{Réemployer} : Utilisez au moins 5 mots de cette liste dans votre production écrite de la semaine (essai, lettre, dialogue du jeu de rôle).
\item \textbf{Réviser espacé} : Relisez vos fiches à J+1, J+3, J+7, J+15.
\end{enumerate}
\end{methode}

\begin{aretenir}[Lexique minimum à maîtriser pour l'évaluation]
\begin{itemize}
\item \textbf{Décès / deuil} : \esp{fallecer, la muerte, el funeral, el entierro, el luto, estar de luto, la viudez, viudo/a, huérfano/a, echar de menos}.
\item \textbf{Héritage} : \esp{la herencia, heredar, heredero/a, el testamento, dejar en herencia, la sucesión, el notario, los bienes, la propiedad, el reparto}.
\item \textbf{Conflit / paix} : \esp{el conflicto familiar, la disputa, pelearse, negarse a, repartir, llegar a un acuerdo, reconciliarse, perdonar, hacer las paces}.
\item \textbf{Émotions} : \esp{la tristeza, el dolor, la rabia, la alegría, sentir, sentirse, echar de menos}.
\item \textbf{Registre} : \esp{fallecer} (formel) vs \esp{morir} (courant) ; \esp{hacer las paces} (courant) vs \esp{reconciliarse} (formel).
\item \textbf{Prononciation} : \esp{h} muet (\emph{herencia, huérfano, hacer}) ; \esp{d} intervocalique \textipa{[ð]} (\emph{viudo, herencia, podido}) ; accent écrit sur proparoxytons (\emph{testamento, último, sábado}).
\end{itemize}
\end{aretenir}

\courssub{Jeu de rôle : la conversation familiale — partage d'héritage et réconciliation}

\begin{experience}[Jeu de rôle : « La lecture du testament chez maître Nkodo »]
\textbf{Situation :} Chez le notaire, maître Nkodo, à Yaoundé. La famille de don Patrice (veuve + 4 enfants) découvre le testament. Des biens ne sont pas mentionnés (terrain Nkolbisson, bijoux grand-mère). Tension entre l'aîné (Pierre) qui veut tout garder et la sœur (Marie) qui exige transparence. La mère (Doña Hélène) impose le calme. Le cadet (Jean) propose une solution équitable. Le notaire rappelle les conséquences judiciaires.

\textbf{Rôles (5 élèves) :}
\begin{enumerate}
\item \textbf{Maître Nkodo (Notaire) :} Neutre, juridique, calme. Utilise : \esp{« La sucesión está abierta », « Ustedes son herederos legítimos », « Si no hay acuerdo, iremos a juicio »}.
\item \textbf{Doña Hélène (Mère, veuve) :} Digne, fatiguée, cherche la paix. Utilise : \esp{« Basta », « Haced las paces », « Yo no quiero herencia, quiero paz »}.
\item \textbf{Pierre (Fils aîné) :} Au début égoïste, défensif, puis honteux. Utilise : \esp{« Me niego a firmar », « Es una trampa », « Tienes razón, me avergüenzo »}.
\item \textbf{Marie (Fille) :} Combative, juste, émue. Utilise : \esp{« No voy a permitir », « Exijo claridad », « Vale, vendamos »}.
\item \textbf{Jean (Fils cadet) :} Conciliateur, pratique. Utilise : \esp{« Papá nos enseñó a compartir », « Lleguemos a un acuerdo », « Nos reconciliamos »}.
\end{enumerate}

\textbf{Déroulement (10 min préparation + 5 min jeu) :}
\begin{enumerate}
\item Lecture silencieuse du testament (inventé par le notaire : maison Yaoundé, finca Ebolowa, compte banque, voiture pour les 4 ; terrain Nkolbisson et bijoux « à régler entre frères »).
\item Improvisation guidée : le notaire ouvre la séance. Pierre et Marie s'opposent. Hélène intervient. Jean propose. Accord final : vente maison Yaoundé, partage argent, finca à mère, terrain Nkolbisson partagé ou vendu.
\item Le notaire fait signer l'accord (\emph{acta de acuerdo}).
\end{enumerate}

\textbf{Expressions utiles au tableau :}
\begin{center}
\begin{tabularx}{\linewidth}{|X|X|}
\hline
\rowcolor{popblueL}
\esp{El testamento establece que…} & Le testament stipule que… \\
\esp{No estoy de acuerdo con…} & Je ne suis pas d'accord avec… \\
\esp{Propongo que vendamos…} & Je propose que nous vendions… (subjonctif !) \\
\esp{¿Podemos llegar a un acuerdo?} & Pouvons-nous trouver un accord ? \\
\esp{Por el bien de la familia…} & Pour le bien de la famille… \\
\esp{Firmemos el reparto.} & Signons le partage. \\
\hline
\end{tabularx}
\end{center}
\end{experience}

\courssub{Production guidée et bilan}

\begin{exercice}[Production écrite au choix — 150 à 180 mots]
Choisissez \textbf{un} des trois sujets ci-dessous. Respectez le format demandé (lettre formelle, dialogue théâtral, essai argumenté). Utilisez au moins \textbf{10 mots} du lexique de la leçon (surlignez-les dans votre copie).

\textbf{Sujet A — Lettre formelle (Theme/Version) :}
Vous êtes maître Nkodo, notaire à Yaoundé. Écrivez une lettre formelle aux quatre héritiers pour les convoquer à la signature de l'acte de partage définitif. Mentionnez : date, heure, lieu, documents à fournir (pièce d'identité, certificat de décès, livret de famille), conséquence en cas d'absence (procédure judiciaire). Formules d'appel et de politesse obligatoires.

\textbf{Sujet B — Dialogue théâtral (Expression orale/écrite) :}
Réécrivez la scène de réconciliation sous le manguier (dernier paragraphe du texte « El último adiós ») en développant les didascalies (ton, gestes, regards) et en ajoutant les pensées intérieures de Pierre (remords) et Marie (soulagement). Intégrez : \esp{hacer las paces, echar de menos, repartir, heredero, viudez}.

\textbf{Sujet C — Essai argumenté (Commentaire/Expression) :}
« L'héritage : source de discorde ou ciment de la famille ? » À partir du texte étudié et de votre expérience (contexte camerounais : coutumes, droit moderne, rôle du notaire, place de la mère), développez une réflexion structurée (introduction, thèse/antithèse, synthèse, conclusion). Mots-outils : \esp{por un lado, por otro, en conclusión, sin embargo, por tanto}.
\end{exercice}

\begin{methode}[Réussir sa production écrite au Bac (Espagnol LV2)]
\begin{enumerate}
\item \textbf{Analyser le sujet} : Format (lettre, dialogue, essai) ? Registre (formel, familier) ? Contraintes (mots obligatoires, nombre de mots).
\item \textbf{Planifier (brouillon)} : Idées principales $\to$ structure (paragraphes) $\to$ sélection du vocabulaire ciblé (lexique leçon + connecteurs logiques).
\item \textbf{Rédiger} : Phrases courtes et correctes. \textbf{Vérifier les accords} (genre/nombre), \textbf{conjugaisons} (subjonctif après \emph{quiero que, propongo que, es necesario que}), \textbf{ser/estar}, \textbf{por/para}.
\item \textbf{Relire (check-list)} : Accents ? ¿ ¡ ? Ponctuation ? Mots surlignés ? Nombre de mots respecté ? Lisibilité ?
\end{enumerate}
\end{methode}

\begin{aretenir}[Bilan de la leçon E05 — Compétences validées]
\begin{itemize}
\item \textbf{Compréhension écrite} : Identifier les champs lexicaux (deuil, héritage, conflit, émotions) dans un texte long authentique ; inférer le sens de mots inconnus par le contexte (\emph{tañe, delatan, esboza}).
\item \textbf{Lexique} : Maîtriser 40+ unités lexicales (noms, verbes, locutions) avec registre (formel/courant), genre, construction syntaxique (\emph{negarse a, echar de menos, dejar en herencia}).
\item \textbf{Grammaire en contexte} : Subjonctif de valeur volitive (\emph{propongo que vendamos, hagan las paces}) ; pronominaux (\emph{pelearse, reconciliarse, negarse}) ; diphtongues (\emph{muero, siento, reparte}) ; \emph{ser/estar} (\emph{está de luto, es viuda, está destrozada}).
\item \textbf{Expression orale} : Jouer un rôle codifié (notaire, héritiers) en respectant les tours de parole, le registre et la prononciation (h muet, \emph{d} intervocalique, accent tonique).
\item \textbf{Expression écrite} : Produire un texte cohérent, structuré, formaté (lettre/dialogue/essai) en mobilisant le lexique et la grammaire de la leçon.
\item \textbf{Culture} : Connaître les rites funéraires (veuvage, deuil), le vocabulaire notarial (testament, succession, notaire) et la valeur de la réconciliation familiale dans le monde hispanophone et camerounais.
\end{itemize}
\end{aretenir}

\coursec{Méthode du commentaire de texte}

Dans la section précédente, nous avons constitué ensemble le lexique du deuil et de l'héritage : \esp{viudez}, \esp{herencia}, \esp{testamento}, \esp{luto}, \esp{huérfano}... Ces mots sont nos outils. Mais à l'examen comme dans la vie, posséder des outils ne suffit pas : il faut savoir s'en servir. Nous allons maintenant apprendre à les utiliser dans l'exercice roi de la Terminale A : le \cle{commentaire de texte} (\esp{comentario de texto}), qui consiste à analyser de manière organisée un texte comme \emph{La herencia de don Ramón}.

\courssub{Qu'est-ce qu'un commentaire de texte ?}

\begin{definition}[Le comentario de texto]
Le \cle{commentaire de texte} (\esp{comentario de texto}) est une \textbf{analyse organisée} qui dégage le \textbf{sens}, les \textbf{thèmes} et l'\textbf{intérêt} d'un texte. Il ne s'agit ni de traduire, ni de résumer ligne par ligne, mais de répondre à trois questions : de quoi parle le texte ? Comment l'auteur en parle-t-il ? Quelle opinion argumentée puis-je en donner ?
\end{definition}

Observez d'abord la différence entre ces deux productions d'élèves, face au même texte :

\esp{Élève A : « Don Ramón muere. Carmen llora. Los hijos discuten por el dinero. Después hablan con un tío y se reconcilian. »} → simple résumé, phrase par phrase : \textbf{insuffisant}.

\esp{Élève B : « Este texto muestra cómo la muerte de un padre puede destruir o fortalecer una familia. La herencia, símbolo del amor del padre, se convierte aquí en fuente de conflicto. »} → analyse du sens et des enjeux : \textbf{c'est un commentaire}.

\begin{attention}[Le piège n° 1 des francophones]
Ne \textbf{traduisez pas} le texte dans votre commentaire et ne le \textbf{résumez pas} ligne par ligne. L'examinateur connaît le texte : il attend votre \textbf{analyse} et votre \textbf{avis argumenté}. Répéter \esp{« En la línea 3, el texto dice que... »} en permanence est une erreur. Citez seulement de courts extraits pour \emph{appuyer} une idée, puis commentez-les.
\end{attention}

\courssub{Étape 1 : la lecture active}

Avant d'écrire la moindre phrase, il faut \textbf{lire le texte deux fois}.

\begin{methode}[Les deux lectures]
\begin{enumerate}
\item \textbf{Première lecture} : lire pour comprendre l'histoire globale, sans dictionnaire. Qui sont les personnages ? Que se passe-t-il ?
\item \textbf{Deuxième lecture} : souligner les \cle{mots-clés} (\esp{viudez, herencia, conflicto, reconciliación}...), repérer le \textbf{thème principal} et les \textbf{thèmes secondaires}.
\item \textbf{Formuler le thème} en une phrase : \esp{« El tema principal es el conflicto familiar provocado por una herencia. »} « Le thème principal est le conflit familial provoqué par un héritage. »
\end{enumerate}
\end{methode}

Pour \emph{La herencia de don Ramón}, le thème principal est la famille en crise ; les thèmes secondaires sont le deuil (\esp{el luto}), le partage des biens (\esp{el reparto de los bienes}) et la réconciliation (\esp{la reconciliación}).

\courssub{Étape 2 : rédiger l'introduction}

\begin{methode}[L'introduction en 3 phrases]
Une bonne introduction tient en \textbf{trois phrases} :
\begin{enumerate}
\item \textbf{Nature du texte} : \esp{« Este texto es un relato / un artículo / un diálogo... »}
\item \textbf{Résumé du thème} : \esp{« Este texto narra los problemas de una familia tras la muerte del padre. »} « Ce texte raconte les problèmes d'une famille après la mort du père. »
\item \textbf{Annonce du plan} : \esp{« Estudiaremos primero..., luego..., y finalmente... »} « Nous étudierons d'abord..., ensuite..., et enfin... »
\end{enumerate}
\end{methode}

\begin{exemplebox}[Phrases d'introduction traduites]
\esp{Este texto narra los problemas de una familia tras la muerte del padre.} « Ce texte raconte les problèmes d'une famille après la mort du père. »

\esp{El tema principal es la viudez de Carmen y el reparto de la herencia.} « Le thème principal est le veuvage de Carmen et le partage de l'héritage. »

\esp{Sin embargo, la herencia provoca un conflicto entre los hermanos.} « Cependant, l'héritage provoque un conflit entre les frères et sœurs. »
\end{exemplebox}

\courssub{Étape 3 : construire le développement}

Le développement (\esp{el desarrollo}) se construit en \textbf{deux ou trois parties}, chacune consacrée à un aspect du texte. Pour notre texte, un plan en trois parties s'impose naturellement :

\begin{propriete}[Plan type du développement]
\textbf{I.} \esp{El luto y la viudez} — le deuil et le veuvage : la douleur de Carmen, la solitude, la situation des enfants devenus \esp{huérfanos}.

\textbf{II.} \esp{El conflicto por la herencia} — le conflit pour l'héritage : les disputes entre frères, le rôle du \esp{testamento}, les reproches.

\textbf{III.} \esp{La reconciliación posible} — la réconciliation possible : le dialogue, le rôle de la famille élargie, la leçon finale.
\end{propriete}

Chaque partie doit contenir trois éléments, dans cet ordre : une \textbf{idée} (phrase d'introduction de partie), une \textbf{preuve} tirée du texte (courte citation entre guillemets) et un \textbf{commentaire personnel} (ton analyse : que montre cette citation ? pourquoi l'auteur l'écrit-il ?).

\courssub{Étape 4 : conclure}

La conclusion (\esp{la conclusión}) comporte deux temps : le \textbf{bilan} (on résume l'essentiel de l'analyse en une ou deux phrases) et l'\textbf{ouverture} (on élargit le sujet). L'ouverture est très appréciée au baccalauréat : comparez la situation du texte avec la réalité \textbf{camerounaise} (les successions dans nos familles, le rôle des chefs traditionnels, les réunions de famille) ou donnez votre \textbf{opinion personnelle}.

\esp{En conclusión, este texto nos enseña que el diálogo es la única herencia que no se puede dividir. En Camerún, como en el texto, las familias se reúnen para resolver los conflictos de sucesión.} « En conclusion, ce texte nous enseigne que le dialogue est le seul héritage qu'on ne peut pas diviser. Au Cameroun, comme dans le texte, les familles se réunissent pour résoudre les conflits de succession. »

\courssub{Les connecteurs : la charpente du commentaire}

Un commentaire sans connecteurs est une maison sans poutres. Voici les outils indispensables :

\begin{center}
\begin{tabularx}{\linewidth}{|X|X|X|}
\hline
\rowcolor{popblueL} \textbf{Español} & \textbf{Français} & \textbf{Emploi} \\
\hline
\esp{en primer lugar} & en premier lieu & ouvrir la première partie \\
\hline
\esp{además / también} & de plus / aussi & ajouter une idée \\
\hline
\esp{sin embargo} & cependant & opposer deux idées \\
\hline
\esp{por un lado / por otro lado} & d'un côté / de l'autre & présenter deux aspects \\
\hline
\esp{por eso / por lo tanto} & c'est pourquoi / donc & conséquence \\
\hline
\esp{en mi opinión / pienso que} & à mon avis / je pense que & avis personnel \\
\hline
\esp{en conclusión / para terminar} & en conclusion / pour finir & ouvrir la conclusion \\
\hline
\end{tabularx}
\end{center}

\begin{attention}[Faux amis et calques à éviter]
\begin{itemize}
\item \esp{actualmente} signifie « actuellement », jamais « actuel » au sens de « réel ».
\item Ne dites pas \esp{« en el otro lado »} : la forme correcte est \esp{por un lado / por otro lado}.
\item « D'abord » se traduit \esp{primero} ou \esp{en primer lugar} ; « ensuite » se dit \esp{luego} ou \esp{después}, jamais \esp{ensuite}.
\item Attention à l'accent : \esp{opinión}, \esp{conclusión}, \esp{también}. Un accent oublié change parfois le temps du verbe : \esp{hablo} « je parle » (présent) / \esp{habló} « il a parlé » (passé simple).
\item « Devenir veuve » se dit \esp{quedarse viuda} : \esp{Carmen se queda viuda con tres hijos.} Ne calquez pas \esp{« se hace viuda »}.
\end{itemize}
\end{attention}

\begin{popfigure}
\begin{center}
\begin{tikzpicture}[
  node distance=0.55cm,
  etape/.style={rectangle, rounded corners, draw=popblue, fill=popblueL, thick, align=center, text width=5.6cm, minimum height=1.05cm, font=\small\bfseries},
  conn/.style={rectangle, rounded corners, draw=poporange, fill=poporangeL, align=center, text width=4.6cm, font=\footnotesize},
  fleche/.style={-{Stealth[length=3mm]}, thick, popdark}
]
\node[etape, align=center] (intro){INTRODUCTION\\ \footnotesize\mdseries 3 frases : naturaleza, tema, plan};
\node[conn, below=of intro] (c1) {\esp{Este texto narra... / Estudiaremos primero...}};
\node[etape, below=of c1, align=center] (des){DESARROLLO\\ \footnotesize\mdseries I. El luto y la viudez\\ II. El conflicto por la herencia\\ III. La reconciliación posible};
\node[conn, below=of des, align=center] (c2){\esp{en primer lugar, además, sin embargo,}\\\esp{por un lado / por otro lado}};
\node[etape, below=of c2, align=center] (conc){CONCLUSIÓN\\ \footnotesize\mdseries balance + apertura (Camerún, opinión)};
\node[conn, below=of conc] (c3) {\esp{en conclusión, para terminar, en mi opinión}};
\draw[fleche] (intro) -- (c1);
\draw[fleche] (c1) -- (des);
\draw[fleche] (des) -- (c2);
\draw[fleche] (c2) -- (conc);
\draw[fleche] (conc) -- (c3);
\end{tikzpicture}
\end{center}
\legende{Organigramme du commentaire de texte : chaque étape s'appuie sur ses connecteurs.}
\end{popfigure}

\courssub{Commentaire modèle : début rédigé et commenté}

Voici le début d'un commentaire modèle sur \emph{La herencia de don Ramón}, tel qu'on le construit ensemble en classe :

\begin{textebox}[Comentario modelo : introducción y primera parte]
El texto trata de la familia y sus problemas: la viudez de Carmen y el reparto de la herencia de don Ramón. Tras la muerte del padre, los tres hijos descubren el testamento y empieza un conflicto que amenaza con destruir la familia. Estudiaremos primero el dolor de la familia, luego el conflicto entre los hermanos y finalmente la posibilidad de reconciliación.

En primer lugar, el texto presenta el luto de Carmen, que se queda viuda con tres hijos. La frase «la casa quedó en silencio» muestra muy bien la soledad de la madre. Además, los hijos son ahora huérfanos y deben aprender a vivir sin su padre.
\end{textebox}

\textbf{Glossaire :} \esp{trata de} = il s'agit de ; \esp{el reparto} = le partage ; \esp{tras} = après ; \esp{amenaza con} = menace de ; \esp{se queda viuda} = elle devient veuve ; \esp{la soledad} = la solitude.

\textbf{Analyse de la méthode :} remarquez la structure en 3 phrases de l'introduction (nature du texte → thème → plan annoncé avec \esp{primero... luego... finalmente}). Dans la première partie, l'élève cite une \emph{courte} phrase du texte (« la casa quedó en silencio ») puis la \emph{commente} : idée + preuve + analyse, c'est exactement ce qu'attend l'examinateur.

\begin{savaistu}[Complément Bac — la problématique]
\emph{(Hors strict programme, mais très valorisé au baccalauréat.)} Dans l'introduction, on peut formuler une \textbf{problématique} sous forme de question en espagnol : \esp{¿Cómo afecta la herencia a las relaciones familiares?} « Comment l'héritage affecte-t-il les relations familiales ? » Tout le développement répondra alors à cette question. Autres exemples : \esp{¿Por qué el dinero divide a las familias?} ; \esp{¿Es posible reconciliarse después de un conflicto?}
\end{savaistu}

\begin{exoresolu}[Rédiger une introduction complète]
\textbf{Énoncé.} À partir du texte \emph{La herencia de don Ramón}, rédigez une introduction de commentaire en trois phrases (nature du texte, thème, annonce du plan), avec une problématique.
\tcblower
\textbf{Solution détaillée.}

\esp{Este texto es un relato sobre la vida familiar. Narra la historia de Carmen, una mujer que se queda viuda, y de sus tres hijos, que se pelean por la herencia de don Ramón. ¿Cómo afecta la herencia a las relaciones familiares? Para responder a esta pregunta, estudiaremos primero el luto y la viudez, luego el conflicto entre los hermanos y finalmente la reconciliación posible.}

Traduction : « Ce texte est un récit sur la vie familiale. Il raconte l'histoire de Carmen, une femme qui devient veuve, et de ses trois enfants, qui se disputent pour l'héritage de don Ramón. Comment l'héritage affecte-t-il les relations familiales ? Pour répondre à cette question, nous étudierons d'abord le deuil et le veuvage, ensuite le conflit entre les frères et sœurs, et enfin la réconciliation possible. »

\textbf{Points de méthode :} phrase 1 = nature du texte (\esp{es un relato}) ; phrase 2 = résumé du thème en une seule phrase ; phrase 3 = problématique (\esp{¿Cómo afecta...?}) ; phrase 4 = annonce du plan avec les connecteurs \esp{primero / luego / finalmente}. Notez l'accord : \esp{una mujer que se queda viuda} (féminin singulier).
\end{exoresolu}

\begin{exercice}[À vous de jouer]
\begin{enumerate}
\item Remettez dans l'ordre les étapes du commentaire : conclusión / lectura activa / desarrollo / introducción.
\item Traduisez en espagnol : « Cependant, l'héritage provoque un conflit entre les frères. »
\item Rédigez la phrase d'ouverture de la partie II (\esp{El conflicto por la herencia}) en utilisant \esp{sin embargo}.
\item Formulez une problématique en espagnol sur le thème du veuvage.
\end{enumerate}
\end{exercice}

\begin{corrige}[Exercice « À vous de jouer »]
\begin{enumerate}
\item Ordre correct : lectura activa → introducción → desarrollo → conclusión.
\item \esp{Sin embargo, la herencia provoca un conflicto entre los hermanos.}
\item Exemple : \esp{Sin embargo, el testamento de don Ramón provoca una fuerte discusión entre los hijos.} « Cependant, le testament de don Ramón provoque une vive dispute entre les enfants. »
\item Exemple : \esp{¿Cómo cambia la vida de una familia después de la muerte del padre?} « Comment la vie d'une famille change-t-elle après la mort du père ? »
\end{enumerate}
\end{corrige}

\begin{aretenir}[L'essentiel de la méthode]
\begin{itemize}
\item Le commentaire \textbf{analyse}, il ne traduit ni ne résume : sens, thèmes, intérêt du texte.
\item Structure en 4 temps : \textbf{lecture active} (2 lectures, mots-clés) → \textbf{introduction} (3 phrases : nature, thème, plan) → \textbf{développement} (2 ou 3 parties : idée + preuve + analyse) → \textbf{conclusion} (bilan + ouverture sur le Cameroun ou avis personnel).
\item Les connecteurs sont obligatoires : \esp{en primer lugar, además, sin embargo, por un lado / por otro lado, en conclusión}.
\item Au Bac, ajoutez une problématique : \esp{¿Cómo afecta la herencia a las relaciones familiares?}
\end{itemize}
\end{aretenir}

Maintenant que nous savons analyser un texte à l'écrit, passons à l'oral : dans la section suivante, vous incarnerez vous-mêmes Carmen et ses enfants dans un \esp{juego de rol} autour du partage de l'héritage.

\coursec{Jeu de rôle : la conversation familiale}

Après avoir appris à lire et à commenter le texte \esp{La herencia de don Ramón}, il est temps de passer de la compréhension à la production orale. Le commentaire de texte nous a permis d'analyser le conflit familial de l'extérieur ; le jeu de rôle va nous plonger à l'intérieur de la scène : chaque élève devient un personnage et doit défendre son point de vue en espagnol. C'est l'aboutissement naturel de la leçon : réutiliser le lexique du deuil et de l'héritage (\esp{viudez, herencia, testamento, huérfano, luto}) dans une situation de communication authentique.

\courssub{Objectif et mise en situation}

\begin{definition}[Qu'est-ce qu'un jeu de rôle ?]
Un \cle{jeu de rôle} (\esp{juego de rol}) est une activité orale dans laquelle chaque participant incarne un personnage dans une situation donnée. Contrairement à la récitation, le jeu de rôle exige d'\cle{improviser} à partir d'un canevas : on connaît son personnage, ses désirs et ses émotions, mais on doit réagir en direct aux répliques des autres. En classe d'espagnol, c'est l'outil idéal pour mobiliser le lexique et les structures grammaticales dans un contexte réaliste.
\end{definition}

\textbf{La situation.} Don Ramón vient de mourir. Il a laissé une maison familiale, un peu d'argent et un champ de cacao. Sa veuve Carmen réunit ses trois enfants, Luis, Ana et Pablo, pour lire le testament et décider ensemble du partage de l'héritage. Mais chacun a ses propres intérêts : Luis veut sa part d'argent tout de suite, Ana refuse de vendre la maison, Pablo cherche un compromis. La tension monte... jusqu'à la réconciliation.

\begin{textebox}[Amorce de dialogue : la réunion commence]
\textbf{Carmen :} Hijos, escuchadme. Vuestro padre quería que la familia estuviera unida.\\
\textbf{Luis :} Pero yo necesito mi parte de la herencia. Tengo deudas y no puedo esperar.\\
\textbf{Ana :} ¡No podemos vender la casa! Aquí crecimos nosotros, aquí vive mamá.\\
\textbf{Pablo :} ¿Por qué no llegamos a un acuerdo? Propongo que hablemos con calma.\\
\textbf{Carmen :} Pablo tiene razón. Escuchémonos, por favor. Vuestro padre odiaba las peleas.\\
\textbf{Luis :} Está bien, mamá. Perdóname, Ana, no quería gritar.
\end{textebox}

\textbf{Glossaire :} \esp{escuchadme} = écoutez-moi (impératif + pronom) ; \esp{estuviera unida} = soit unie (subjonctif imparfait) ; \esp{mi parte de la herencia} = ma part de l'héritage ; \esp{tengo deudas} = j'ai des dettes ; \esp{crecimos} = nous avons grandi ; \esp{llegamos a un acuerdo} = nous parvenons à un accord ; \esp{las peleas} = les disputes.

\courssub{Répartition des rôles}

La classe se divise en groupes de quatre ou cinq élèves. Chaque groupe distribue les rôles suivants :

\begin{center}
\begin{tabularx}{\linewidth}{|l|X|X|}
\hline
\rowcolor{popblueL} \textbf{Personaje} & \textbf{Carácter y deseos} & \textbf{Réplique type} \\
\hline
\esp{La viuda, Carmen} & Veut préserver l'unité de la famille ; porte encore le deuil (\esp{luto}) & \esp{Vuestro padre quería que estuviéramos unidos.} \\
\hline
\esp{El hijo mayor, Luis} & Pressé, endetté, réclame sa part & \esp{Exijo mi parte de la herencia.} \\
\hline
\esp{La hija, Ana} & Sentimentale, attachée à la maison & \esp{Me niego a vender la casa familiar.} \\
\hline
\esp{El hijo menor, Pablo} & Médiateur, cherche le compromis & \esp{Propongo que repartamos el dinero en tres partes.} \\
\hline
\esp{El notario} (optionnel) & Lit le testament, arbitre neutre & \esp{Según el testamento, la casa pertenece a los tres herederos.} \\
\hline
\end{tabularx}
\end{center}

\begin{popfigure}
\begin{tikzpicture}
\node[draw=popdark, fill=popgoldL, rounded corners, align=center, font=\bfseries, minimum width=3.2cm, minimum height=1.1cm] (mesa) at (0,0) {La mesa\\ familiar};
\node[draw=popblue, fill=popblueL, rounded corners, align=center, minimum width=3.4cm] (carmen) at (0,3.2) {\textbf{Carmen (la viuda)}\\ \footnotesize \esp{Escuchémonos, por favor.}};
\node[draw=poporange, fill=poporangeL, rounded corners, align=center, minimum width=3.4cm] (luis) at (4.6,0) {\textbf{Luis (hijo mayor)}\\ \footnotesize \esp{Necesito mi parte.}};
\node[draw=poppurple, fill=poppurpleL, rounded corners, align=center, minimum width=3.4cm] (pablo) at (0,-3.2) {\textbf{Pablo (hijo menor)}\\ \footnotesize \esp{¿Por qué no llegamos a un acuerdo?}};
\node[draw=popgreen, fill=popgreenL, rounded corners, align=center, minimum width=3.4cm] (ana) at (-4.6,0) {\textbf{Ana (la hija)}\\ \footnotesize \esp{¡No vendamos la casa!}};
\draw[-{Stealth}, popline, thick] (carmen) -- (mesa);
\draw[-{Stealth}, popline, thick] (luis) -- (mesa);
\draw[-{Stealth}, popline, thick] (pablo) -- (mesa);
\draw[-{Stealth}, popline, thick] (ana) -- (mesa);
\end{tikzpicture}
\legende{Les quatre personnages autour de la table familiale, avec une réplique type chacun.}
\end{popfigure}

\courssub{Les fonctions langagières à mobiliser}

Pour jouer la scène, il faut maîtriser cinq \cle{fonctions langagières} précises. Observons d'abord les exemples, puis dégageons les règles.

\begin{exemplebox}[Observer avant de généraliser]
\esp{Quiero que estemos unidos.} « Je veux que nous soyons unis. »\\
\esp{Me niego a vender la casa.} « Je refuse de vendre la maison. »\\
\esp{Propongo que repartamos el dinero en tres partes.} « Je propose que nous partagions l'argent en trois parts. »\\
\esp{¿Por qué no hablamos con calma?} « Pourquoi ne parlons-nous pas calmement ? »\\
\esp{Cálmate, hermano. Escuchémonos.} « Calme-toi, frère. Écoutons-nous. »\\
\esp{Perdóname, hermano, no quería pelear.} « Pardonne-moi, frère, je ne voulais pas me disputer. »\\
\esp{Hagamos las paces.} « Faisons la paix. »
\end{exemplebox}

\begin{propriete}[Exprimer un souhait : \esp{querer que} + subjonctif]
Quand le sujet de la volonté est \textbf{différent} du sujet de l'action souhaitée, l'espagnol impose la structure : \esp{querer que} + \textbf{subjonctif}. Forme : \esp{Quiero que} + verbe au subjonctif présent (\esp{estemos, hablemos, vendas, respetes}). Emploi : exprimer un désir, une exigence, un ordre indirect. Exception : si le sujet est \textbf{identique}, on utilise l'infinitif : \esp{Quiero hablar} « Je veux parler » (et non \esp{quiero que hable}).
\end{propriete}

\begin{propriete}[Refuser, proposer, apaiser, réconcilier]
\begin{itemize}
\item \textbf{Refuser} : \esp{negarse a} + infinitif (verbe pronominal : \esp{me niego a, te niegas a...}). Exemple : \esp{Me niego a firmar.} « Je refuse de signer. »
\item \textbf{Proposer} : \esp{proponer que} + subjonctif (\esp{Propongo que repartamos...}) ; ou la tournure plus légère \esp{¿Por qué no} + indicatif présent (\esp{¿Por qué no vendemos el campo?}).
\item \textbf{Apaiser} : impératif avec pronom soudé : \esp{cálmate} « calme-toi », \esp{escuchémonos} « écoutons-nous » (subjonctif de première personne du pluriel = impératif de \esp{nosotros}).
\item \textbf{Réconcilier} : \esp{perdóname} « pardonne-moi », \esp{lo siento} « je suis désolé », \esp{hagamos las paces} « faisons la paix ».
\end{itemize}
\end{propriete}

\begin{attention}[Le piège du subjonctif après \esp{quiero que} et \esp{propongo que}]
L'erreur la plus fréquente des francophones est d'utiliser l'indicatif après ces verbes de volonté, par calque du français. On dit \esp{Quiero que \textbf{estemos} unidos} et jamais \esp{quiero que estamos unidos}. De même : \esp{Propongo que \textbf{repartamos} el dinero} (et non \esp{repartimos}). Rappel des formes du subjonctif présent utiles ici : \esp{que yo venda, que tú respetes, que él reciba, que estemos, que repartamos, que hagamos}. Autre piège : ne pas oublier le \esp{que} : en espagnol, il est obligatoire, contrairement au français où « que » peut parfois être sous-entendu à l'oral.
\end{attention}

\begin{center}
\begin{tabularx}{\linewidth}{|X|X|X|}
\hline
\rowcolor{popblueL} \textbf{Fonction} & \textbf{Español} & \textbf{Français} \\
\hline
Souhaiter & \esp{Quiero que estemos unidos.} & Je veux que nous soyons unis. \\
\hline
Refuser & \esp{Me niego a vender la casa.} & Je refuse de vendre la maison. \\
\hline
Proposer & \esp{Propongo que repartamos el dinero.} & Je propose que nous partagions l'argent. \\
\hline
Proposer (léger) & \esp{¿Por qué no hablamos?} & Pourquoi ne parlons-nous pas ? \\
\hline
Apaiser & \esp{Cálmate, escuchémonos.} & Calme-toi, écoutons-nous. \\
\hline
Demander pardon & \esp{Perdóname, no quería pelear.} & Pardonne-moi, je ne voulais pas me disputer. \\
\hline
Réconcilier & \esp{Hagamos las paces.} & Faisons la paix. \\
\hline
\end{tabularx}
\end{center}

\courssub{Préparation en groupes : rédiger le script}

\begin{methode}[Comment réussir son jeu de rôle en quatre étapes]
\begin{enumerate}
\item \textbf{Distribuer les rôles} et réfléchir au caractère de son personnage : que veut-il ? Quelle émotion domine (colère, tristesse, calme) ?
\item \textbf{Rédiger un script de 8 à 12 répliques} : chaque personnage prépare au moins \textbf{3 répliques-clés} utilisant les structures de la leçon (\esp{quiero que} + subjonctif, \esp{me niego a}, \esp{propongo que}...).
\item \textbf{Intégrer au moins 5 mots du lexique} de la leçon : \esp{herencia, testamento, viuda, huérfano, luto, heredero, conflicto familiar}.
\item \textbf{Répéter à voix haute} en soignant l'intonation : la colère de Luis monte, la voix de Carmen est douce mais ferme, Pablo parle calmement. La gestuelle compte autant que les mots.
\end{enumerate}
\end{methode}

\begin{exoresolu}[Construire une réplique correcte]
Énoncé : Luis veut dire « Je veux que tu respectes la volonté de père ». Un élève écrit : \esp{Quiero que respetas la voluntad de padre.} Corrige la phrase et traduis-la.
\tcblower
Solution : après \esp{quiero que}, le verbe doit être au \textbf{subjonctif présent}. Le verbe \esp{respetar} (régulier en \esp{-ar}) donne \esp{respetes} à la deuxième personne du singulier. Il faut aussi l'article devant \esp{padre} : \esp{de papá} ou \esp{del padre}. Phrase corrigée : \esp{Quiero que \textbf{respetes} la voluntad de papá.} Traduction : « Je veux que tu respectes la volonté de papa. » L'erreur initiale (\esp{respetas}, indicatif) est le calque typique du français.
\end{exoresolu}

\courssub{Variante : la scène de réconciliation}

Une fois la scène du conflit jouée, chaque groupe peut rejouer la même réunion, mais cette fois un membre de la famille \textbf{demande pardon} et propose une solution. Par exemple, Luis reconnaît son égoïsme :

\begin{exemplebox}[Modèle de réconciliation]
\esp{Luis : Perdóname, hermana. He sido egoísta. No quería pelear por el dinero.} « Pardonne-moi, ma sœur. J'ai été égoïste. Je ne voulais pas me disputer pour l'argent. »\\
\esp{Ana : Te perdono, Luis. Hagamos las paces.} « Je te pardonne, Luis. Faisons la paix. »\\
\esp{Carmen : Gracias, hijos. Vuestro padre estaría orgulloso de vosotros.} « Merci, mes enfants. Votre père serait fier de vous. »
\end{exemplebox}

Remarque : \esp{estaría orgulloso} est au \textbf{conditionnel} (« serait fier ») : c'est l'espagnol qui exprime ici l'irréel, exactement comme le français.

\courssub{Évaluation et débriefing collectif}

\begin{propriete}[Critères d'évaluation du jeu de rôle]
\begin{itemize}
\item \textbf{Exactitude linguistique} (subjonctif correct, accords, prononciation) : 5 points.
\item \textbf{Richesse du lexique} (au moins 5 mots de la leçon employés à propos) : 5 points.
\item \textbf{Fluidité} (débit naturel, peu d'hésitations, réactions aux répliques des partenaires) : 5 points.
\item \textbf{Jeu de scène} (intonation des émotions, gestuelle, regard, occupation de l'espace) : 5 points.
\end{itemize}
\end{propriete}

\begin{experience}[Débriefing collectif]
Après les représentations, toute la classe discute, d'abord en espagnol puis en français si nécessaire : \esp{¿Qué soluciones ha encontrado cada grupo al conflicto?} « Quelles solutions chaque groupe a-t-il trouvées au conflit ? » Certains groupes auront vendu la maison et partagé l'argent ; d'autres auront gardé la maison pour la mère et partagé le champ ; d'autres encore auront inventé un notaire qui impose la volonté du testament. Le professeur relève au tableau les meilleures répliques entendues et corrige collectivement deux ou trois erreurs récurrentes (subjonctif, faux amis, prononciation de \esp{herencia}).
\end{experience}

\begin{savaistu}[Héritage et coutumes : Cameroun et Guinée équatoriale]
Au Cameroun, comme dans de nombreuses sociétés africaines, la question de l'héritage (\esp{herencia}) suit souvent le droit coutumier : transmission par lignée maternelle ou paternelle, rôle du chef de famille, place de la veuve (\esp{viuda}) et des orphelins (\esp{huérfanos}). En Guinée équatoriale, pays hispanophone voisin, le Code civil espagnol coexiste avec les traditions fang ou bubi. Le \esp{testamento} (testament) écrit gagne du terrain en ville (Yaoundé, Douala, Malabo), mais les conflits familiaux (\esp{conflicto familiar}) surgissent encore quand l'écrit contredit l'oral. Le \esp{luto} (deuil) s'exprime par des rites précis : port du noir, veillées, partage rituel des biens. Connaître ces réalités aide à jouer la scène avec justesse culturelle.
\end{savaistu}

\begin{exercice}[À vous de jouer]
En groupes de quatre, rédigez un script de 10 répliques intitulé \esp{El reparto de la herencia}. Contraintes : chaque personnage parle au moins deux fois ; le script contient au moins une fois \esp{quiero que} + subjonctif, une fois \esp{me niego a}, une fois \esp{¿por qué no...?} et se termine par une réconciliation (\esp{perdóname} / \esp{hagamos las paces}). Jouez ensuite la scène devant la classe.
\end{exercice}

\begin{corrige}[Exercice « El reparto de la herencia »]
Éléments de correction (modèle possible) : \esp{Carmen : Hijos, quiero que respetemos el testamento de vuestro padre. Luis : Me niego a esperar más, necesito mi parte. Ana : ¿Por qué no guardamos la casa para mamá? Pablo : Propongo que vendamos el campo de cacao y repartamos el dinero. Carmen : Me parece bien. Luis : Perdóname, Ana, he sido injusto. Ana : Hagamos las paces, hermano.} Vérifier : présence du subjonctif (\esp{respetemos, repartamos}), du pronom réfléchi dans \esp{me niego}, et d'au moins cinq mots du lexique (\esp{testamento, herencia, viuda, luto, herederos}).
\end{corrige}

\begin{aretenir}[L'essentiel de la section]
Le jeu de rôle transforme le vocabulaire et la grammaire de la leçon en parole vivante. Trois structures à retenir absolument : \esp{quiero que} / \esp{propongo que} + \textbf{subjonctif} ; \esp{negarse a} + infinitif ; les impératifs d'apaisement et de pardon (\esp{cálmate, perdóname, hagamos las paces}). Un bon jeu de rôle repose sur la préparation (3 répliques-clés, 5 mots de lexique) et sur l'expression des émotions par la voix et le geste.
\end{aretenir}

\coursec{Production guidée et bilan}

Après avoir joué la scène de la conversation familiale — le partage de l'héritage de don Ramón, les reproches d'Ana et de Luis, puis la réconciliation —, il est temps de fixer par écrit tout ce que nous avons appris. Cette dernière étape du module transforme l'oral en production écrite : c'est exactement la compétence exigée au baccalauréat, où le commentaire de texte et l'expression écrite personnelle pèsent lourd dans la note. Nous allons donc rédiger, nous auto-évaluer, corriger ensemble, puis vérifier notre lexique avant d'ouvrir le débat sur les coutumes camerounaises et espagnoles.

\begin{popfigure}
\begin{center}
\begin{tikzpicture}[node distance=0.35cm, every node/.style={font=\small}]
\node[draw=popblue, fill=popblueL, rounded corners, minimum width=2.6cm, minimum height=0.9cm, align=center] (a) {\textbf{texto}};
\node[draw=popgreen, fill=popgreenL, rounded corners, minimum width=2.6cm, minimum height=0.9cm, align=center, right=of a] (b) {\textbf{comprensión}};
\node[draw=poporange, fill=poporangeL, rounded corners, minimum width=2.6cm, minimum height=0.9cm, align=center, right=of b] (c) {\textbf{léxico}};
\node[draw=poppurple, fill=poppurpleL, rounded corners, minimum width=2.6cm, minimum height=0.9cm, align=center, below=0.5cm of c] (d) {\textbf{comentario}};
\node[draw=poppink, fill=poppinkL, rounded corners, minimum width=2.6cm, minimum height=0.9cm, align=center, left=of d] (e) {\textbf{juego de rol}};
\node[draw=popgold, fill=popgoldL, rounded corners, minimum width=2.6cm, minimum height=0.9cm, align=center, left=of e] (f) {\textbf{producción}};
\draw[-{Stealth}, thick, popblue] (a) -- (b);
\draw[-{Stealth}, thick, popgreen] (b) -- (c);
\draw[-{Stealth}, thick, poporange] (c) -- (d);
\draw[-{Stealth}, thick, poppurple] (d) -- (e);
\draw[-{Stealth}, thick, poppink] (e) -- (f);
\end{tikzpicture}
\end{center}
\legende{Frise de synthèse du parcours de la leçon : du texte à la production écrite.}
\end{popfigure}

\courssub{Production écrite guidée : le commentaire du texte}

\begin{methode}[Rédiger le commentaire en 15 à 20 lignes]
Le commentaire de texte au baccalauréat suit une architecture en trois temps. Pour \emph{La herencia de don Ramón}, procédez ainsi :
\begin{enumerate}
\item \textbf{Introduction} (3 à 4 lignes) : présentez le texte (type, thème, personnages) et formulez une \cle{problématique}. Exemple : \esp{¿Por qué la herencia, que debería unir a la familia, provoca aquí un conflicto?}
\item \textbf{Développement} (8 à 12 lignes) : suivez le plan détaillé construit en classe. Premier axe : le deuil et ses émotions (\esp{el luto, la tristeza, la soledad de la viuda}). Deuxième axe : le conflit de l'héritage (\esp{el testamento, los herederos, la disputa entre Ana y Luis}). Troisième axe : la réconciliation et sa leçon (\esp{el diálogo, el perdón, la unidad familiar}).
\item \textbf{Conclusion} (3 à 4 lignes) : répondez à la problématique et ouvrez sur une réflexion personnelle ou culturelle (par exemple, les successions au Cameroun).
\end{enumerate}
Consignes de langue : employez au moins \cle{cinq mots du lexique} de la leçon, des \cle{connecteurs} variés et, si possible, une phrase au \cle{subjonctif}.
\end{methode}

\begin{exemplebox}[Connecteurs utiles pour le commentaire]
\begin{center}
\begin{tabularx}{\linewidth}{|X|X|X|}
\hline
\rowcolor{popblueL} Fonction & Español & Français \\
\hline
Introduction du sujet & \esp{En primer lugar} & En premier lieu \\
\hline
Ajout & \esp{Además, Asimismo} & De plus, De même \\
\hline
Opposition & \esp{Sin embargo, No obstante} & Cependant, Néanmoins \\
\hline
Cause & \esp{Por eso, Debido a} & C'est pourquoi, À cause de \\
\hline
Conséquence & \esp{Por lo tanto, Así que} & Par conséquent, Donc \\
\hline
Conclusion & \esp{En conclusión, Para terminar} & En conclusion, Pour terminer \\
\hline
\end{tabularx}
\end{center}
\end{exemplebox}

\begin{propriete}[Rappel de grammaire : le subjonctif après les verbes de volonté et de proposition]
Après les verbes qui expriment un souhait, une proposition ou un conseil suivis de \esp{que}, la subordonnée se met au \cle{subjonctif} quand le sujet change : \esp{esperar que, querer que, proponer que, pedir que, desear que}.
\begin{itemize}
\item \esp{Te propongo que vendamos la casa.} = « Je te propose que nous vendions la maison. »
\item \esp{Espero que aceptes mi propuesta.} = « J'espère que tu accepteras ma proposition. »
\item \esp{Quiero que hablemos con calma.} = « Je veux que nous parlions calmement. »
\end{itemize}
Exception : si le sujet ne change pas, on garde l'\cle{infinitif} : \esp{Quiero hablar con calma.} = « Je veux parler calmement. » (c'est moi qui parle). Rappel des formes : \esp{vender} $\to$ \esp{venda, vendas, venda, vendamos, vendáis, vendan} ; \esp{aceptar} $\to$ \esp{acepte, aceptes, acepte, aceptemos, aceptéis, acepten}.
\end{propriete}

\begin{exoresolu}[Amorce de commentaire]
Énoncé : rédigez l'introduction et la première phrase du développement du commentaire de \emph{La herencia de don Ramón}.
\tcblower
Solution détaillée :
\esp{El texto «La herencia de don Ramón» presenta un conflicto familiar muy frecuente: tras la muerte del padre, sus dos hijos, Ana y Luis, discuten por el testamento y la casa familiar. La viuda, doña Carmen, sufre en silencio esta disputa. Podemos preguntarnos: ¿por qué la herencia divide a las familias en lugar de unirlas? En primer lugar, el texto muestra el dolor del luto: doña Carmen lleva tres meses de luto y se siente sola. Además, la codicia de Luis contrasta con la generosidad de Ana, que propone compartir la casa.}
Traduction : « Le texte "L'héritage de don Ramón" présente un conflit familial très fréquent : après la mort du père, ses deux enfants, Ana et Luis, se disputent à propos du testament et de la maison familiale. La veuve, doña Carmen, souffre en silence de cette dispute. Nous pouvons nous demander : pourquoi l'héritage divise-t-il les familles au lieu de les unir ? En premier lieu, le texte montre la douleur du deuil : doña Carmen porte le deuil depuis trois mois et se sent seule. De plus, la cupidité de Luis contraste avec la générosité d'Ana, qui propose de partager la maison. »
Remarquez : la problématique est introduite par \esp{¿por qué...?}, les connecteurs \esp{en primer lugar} et \esp{además} structurent le propos, et quatre mots du lexique sont déjà employés (\esp{herencia, testamento, viuda, luto}).
\end{exoresolu}

\courssub{Variante : la lettre de réconciliation d'Ana à Luis}

Le jeu de rôle s'est terminé par une réconciliation orale. Écrivons maintenant sa trace écrite : une lettre familière, registre intime, riche en lexique des émotions. Observons d'abord un modèle.

\begin{textebox}[Modelo : la carta de Ana a Luis]
Querido Luis:

Han pasado dos semanas desde nuestra pelea en casa de mamá y no puedo seguir así. Siento mucho nuestra disputa por la herencia de papá; sé que dije palabras duras y te pido perdón. Te echo de menos, hermano, y mamá también sufre porque no nos hablamos.

He pensado mucho en el testamento. Te propongo que vendamos la casa de Douala y que repartamos el dinero en dos partes iguales, como quería papá. Espero que aceptes mi propuesta y que vengas a cenar el domingo. Me gustaría que habláramos con calma, sin gritos, delante de mamá.

La familia es lo más importante que nos queda. No permitamos que el dinero destruya lo que papá construyó con tanto amor.

Un abrazo fuerte,

Ana
\end{textebox}

\textbf{Glossaire} : \esp{la pelea} = la dispute ; \esp{te pido perdón} = je te demande pardon ; \esp{repartir} = partager, répartir ; \esp{en dos partes iguales} = en deux parts égales ; \esp{los gritos} = les cris ; \esp{no permitamos que...} = ne permettons pas que... (subjonctif) ; \esp{un abrazo fuerte} = une forte étreinte, je t'embrasse fort.

\textbf{Questions de compréhension} :
\begin{enumerate}
\item \esp{¿Cuánto tiempo hace que Ana y Luis no se hablan?} (Depuis combien de temps Ana et Luis ne se parlent-ils plus ?)
\item \esp{¿Qué propone Ana para resolver el conflicto?} (Que propose Ana pour résoudre le conflit ?)
\item \esp{¿Qué quiere Ana que hagan el domingo?} (Que veut Ana qu'ils fassent dimanche ?)
\end{enumerate}

\begin{methode}[Écrire une lettre de réconciliation en espagnol]
\begin{enumerate}
\item \textbf{Salutation} : \esp{Querido Luis:} (avec deux-points en espagnol, jamais de virgule seule). Pour une sœur : \esp{Querida Ana:}.
\item \textbf{Expression du regret} : \esp{Siento mucho nuestra pelea.} / \esp{Te pido perdón por mis palabras.}
\item \textbf{Expression du manque} : \esp{Te echo de menos.}
\item \textbf{Proposition concrète} : \esp{Te propongo que vendamos la casa...} — attention, \esp{proponer que} + \cle{subjonctif} : \esp{que vendamos, que hablemos, que vengas}.
\item \textbf{Formule finale} : \esp{Un abrazo, Ana} ou \esp{Un abrazo fuerte, tu hermana Ana}.
\end{enumerate}
\end{methode}

\begin{exemplebox}[Phrases modèles traduites]
\begin{itemize}
\item \esp{Siento mucho nuestra pelea.} = « Je regrette beaucoup notre dispute. »
\item \esp{Te echo de menos.} = « Tu me manques. »
\item \esp{Te propongo que hablemos con calma.} = « Je te propose que nous parlions calmement. »
\item \esp{Espero que aceptes mi propuesta.} = « J'espère que tu accepteras ma proposition. » (subjonctif après \esp{esperar que})
\item \esp{Mamá sufre porque no nos hablamos.} = « Maman souffre parce que nous ne nous parlons pas. »
\end{itemize}
\end{exemplebox}

\begin{attention}[Sentir ou sentarse : le piège classique]
\esp{Sentir} (e $\to$ ie) signifie « regretter, ressentir » : \esp{Siento mucho} = « je suis désolé, je regrette beaucoup ». Ne le confondez pas avec \esp{sentarse} (s'asseoir) : \esp{Me siento} peut vouloir dire « je m'assieds » (de \esp{sentarse}) ou « je me sens » (de \esp{sentirse}) selon le contexte ! Autre piège : pour dire « tu me manques », le français part de celui qui manque, l'espagnol de celui qui ressent : \esp{te echo de menos} se construit avec la personne qui manque comme complément (\esp{te}), jamais \esp{me echas de menos} pour « tu me manques ». En Amérique latine, on dit aussi \esp{te extraño} (même construction).
\end{attention}

\courssub{Grille d'auto-évaluation et correction collective}

Avant de rendre votre copie, évaluez-vous honnêtement avec cette grille.

\begin{center}
\begin{tabularx}{\linewidth}{|X|l|l|}
\hline
\rowcolor{popgreenL} Critère & Sí & No \\
\hline
J'ai utilisé au moins 5 mots du lexique (\esp{viudez, herencia, heredero, testamento, luto, huérfano...}) & & \\
\hline
J'ai employé au moins 3 connecteurs différents & & \\
\hline
Mon commentaire contient une problématique claire & & \\
\hline
Ma lettre respecte la structure (salutation, regret, proposition, formule finale) & & \\
\hline
J'ai utilisé au moins un subjonctif correct & & \\
\hline
J'ai vérifié les accents (\esp{está, papá, atención}) et la ponctuation (¿ ? ¡ !) & & \\
\hline
\end{tabularx}
\end{center}

\begin{experience}[Correction collective au tableau]
Un élève volontaire écrit son introduction au tableau. La classe corrige en quatre passes : d'abord le contenu (la problématique est-elle présente ?), puis le lexique (surlignez les mots de la leçon), ensuite la grammaire (accords, subjonctif, accents), enfin les connecteurs (entourez-les). Chaque élève propose une amélioration : remplacer un \esp{pero} par \esp{sin embargo}, ajouter \esp{además}, transformer une phrase en subjonctif. La copie collective finale sert de modèle de référence pour tous.
\end{experience}

\courssub{Bilan lexical : quiz rapide}

\begin{exercice}[Quiz : de l'espagnol au français]
Traduisez en français, sans dictionnaire, les dix mots suivants :
\begin{enumerate}
\item \esp{la viudez}
\item \esp{la herencia}
\item \esp{el heredero / la heredera}
\item \esp{el testamento}
\item \esp{el conflicto familiar}
\item \esp{el luto}
\item \esp{el huérfano}
\item \esp{la pelea}
\item \esp{te echo de menos}
\item \esp{repartir}
\end{enumerate}
\end{exercice}

\begin{corrige}[Quiz lexical]
1. le veuvage ; 2. l'héritage ; 3. l'héritier / l'héritière ; 4. le testament ; 5. le conflit familial ; 6. le deuil ; 7. l'orphelin ; 8. la dispute ; 9. tu me manques ; 10. partager, répartir. Barème indicatif : 8 à 10 = lexique acquis ; 5 à 7 = revoir les mots manqués ; moins de 5 = refaire la fiche de vocabulaire.
\end{corrige}

\courssub{Ouverture culturelle : héritage et deuil au Cameroun et en Espagne}

\begin{savaistu}[Le Bac valorise connecteurs et subjonctif]
Au baccalauréat d'espagnol, les correcteurs d'expression écrite attribuent des points précis à la \cle{cohésion textuelle} (connecteurs variés) et à la \cle{richesse grammaticale} (subjonctif bien employé). Cette leçon vous a donné les deux : \esp{sin embargo, por lo tanto, en primer lugar} et des structures comme \esp{te propongo que hablemos} ou \esp{espero que aceptes}. Réutilisez-les dans toutes vos copies !
\end{savaistu}

\begin{experience}[Débat oral comparatif]
En groupes, comparez les coutumes autour de trois questions, en espagnol :
\begin{enumerate}
\item \esp{¿Cómo se reparte la herencia en tu región del Camerún? ¿Y en España, donde el testamento reparte los bienes entre los hijos?} (Comment partage-t-on l'héritage dans ta région du Cameroun ? Et en Espagne, où le testament répartit les biens entre les enfants ?)
\item \esp{¿Cuánto dura el luto en tu familia? ¿Qué ropa se lleva?} (Combien de temps dure le deuil dans ta famille ? Quels vêtements porte-t-on ?)
\item \esp{¿Quién hereda la casa familiar: el hijo mayor, todos los hijos, la viuda?} (Qui hérite de la maison familiale : le fils aîné, tous les enfants, la veuve ?)
\end{enumerate}
Chaque groupe présente ensuite une conclusion à la classe, en commençant par \esp{En el Camerún..., sin embargo, en España...}.
\end{experience}

\begin{aretenir}[Bilan de la leçon]
Vous savez désormais : lire et commenter un texte sur les problèmes familiaux avec une problématique et un plan en trois axes ; employer le lexique du deuil et de l'héritage (\esp{viudez, herencia, heredero, testamento, luto, huérfano}) ; écrire une lettre de réconciliation avec \esp{Siento mucho...}, \esp{Te echo de menos}, \esp{Te propongo que...} + subjonctif ; et comparer oralement les coutumes camerounaises et espagnoles. Le parcours est complet : \esp{texto, comprensión, léxico, comentario, juego de rol, producción}.
\end{aretenir}

\newpage

\coursec{Activité d'intégration}

\courssub{Objectifs de l'activité}

Cette activité d'intégration clôt la leçon \emph{La familia y sus problemas : viudez y herencia}. Elle invite l'élève à \cle{mobiliser} dans une situation de communication orale authentique tout ce qui a été étudié : le lexique du deuil et de l'héritage (\esp{viudez, herencia, heredero, testamento, luto, huérfano, conflicto familiar}), les structures d'opinion (\esp{creo que..., me parece que...}), de proposition avec le subjonctif (\esp{propongo que..., es mejor que...}) et les compétences de compréhension et de commentaire du texte \emph{La herencia de don Ramón}.

À la fin de l'activité, l'élève sera capable de :
\begin{itemize}
\item comprendre un document juridique simple en espagnol (extrait de testament) ;
\item rédiger en groupe un script de dialogue de 10 à 12 répliques ;
\item jouer une scène de conversation familiale avec une prononciation et une intonation correctes ;
\item exprimer un jugement et justifier un vote à l'oral en espagnol.
\end{itemize}

\courssub{Situation}

\textbf{Contexte.} Don Ramón Abeso Mba, commerçant de cacao très respecté du quartier Mvolyé à Yaoundé, est décédé il y a trois mois. Après la période de deuil (\esp{luto}), sa famille se réunit chez Maître Ngo Bell, notaire (\esp{notario}), pour la lecture du testament et le partage de l'héritage : une maison à Yaoundé, une plantation de cacao à Obala et un compte en banque. Sont présents : la veuve (\esp{la viuda}), doña Cristina ; le fils aîné, Andrés, qui veut vendre la plantation ; la fille, María, qui veut la garder ; le plus jeune fils, Pablo, étudiant, orphelin de père (\esp{huérfano de padre}), qui réclame de l'argent pour ses études. Le notaire dirige la réunion.

\begin{textebox}[Document : extracto del testamento de don Ramón]
« Yo, Ramón Abeso Mba, en pleno uso de mis facultades mentales, declaro que mis bienes se repartirán de la siguiente manera : la casa de Mvolyé será para mi esposa, Cristina Esono, que vivirá allí hasta su muerte. La plantación de cacao de Obala y el dinero de mi cuenta bancaria se dividirán entre mis tres hijos, Andrés, María y Pablo, en partes iguales. Pido a mis hijos que resuelvan sus diferencias con diálogo y que mantengan unida a la familia. Que nadie venda la plantación sin el acuerdo de los tres herederos. »
\end{textebox}

\textbf{Glossaire du document :} \esp{en pleno uso de mis facultades mentales} = en pleine possession de mes facultés mentales ; \esp{los bienes} = les biens ; \esp{se repartirán} = seront répartis ; \esp{en partes iguales} = en parts égales ; \esp{que nadie venda} = que personne ne vende (subjonctif) ; \esp{el acuerdo} = l'accord.

\textbf{Problème à résoudre :} Andrés veut vendre la plantation, mais le testament exige l'accord des trois héritiers. María s'y oppose. Pablo a besoin d'argent maintenant. La famille doit trouver une solution juste et se réconcilier.

\courssub{Consignes}

\begin{enumerate}
\item \textbf{Compréhension du document.} Relisez l'extrait du testament et répondez par écrit, en espagnol, à trois questions : ¿Qué bienes deja don Ramón ? ¿Quiénes son los herederos ? ¿Qué condición impone el testamento sobre la venta de la plantación ?
\item \textbf{Répartition des rôles.} Par groupes de 4 ou 5, choisissez les personnages : la viuda, el notario, Andrés, María, Pablo (un élève peut jouer deux rôles ou être le narrateur qui présente la scène).
\item \textbf{Rédaction du script.} Rédigez un dialogue de 10 à 12 répliques qui contient obligatoirement : au moins \cle{6 mots} du lexique de la leçon (\esp{viudez, herencia, heredero, testamento, luto, huérfano, conflicto familiar, notario, bienes, repartir}) ; au moins \cle{une proposition au subjonctif} (\esp{propongo que..., sugiero que..., es justo que...}) ; une \cle{scène finale de réconciliation} où la famille trouve un accord.
\item \textbf{Répétition.} Répétez la scène en soignant la prononciation, l'intonation (respect, colère, émotion, apaisement) et les gestes.
\item \textbf{Représentation.} Chaque groupe joue sa scène devant la classe en 3 à 4 minutes.
\item \textbf{Vote et justification.} La classe vote pour la solution de partage la plus juste. Chaque élève justifie son choix en une ou deux phrases en espagnol : \esp{Voto por el grupo... porque su solución respeta el testamento y ayuda a Pablo.}
\end{enumerate}

\courssub{Ressources à mobiliser}

\begin{aretenir}[Lexique utile pour la scène]
\esp{la viudez} (le veuvage), \esp{el luto} (le deuil), \esp{la herencia} (l'héritage), \esp{el/la heredero/a} (l'héritier/héritière), \esp{el testamento} (le testament), \esp{los bienes} (les biens), \esp{repartir / dividir} (partager/diviser), \esp{el conflicto familiar} (le conflit familial), \esp{reconciliarse} (se réconcilier), \esp{estar de acuerdo} (être d'accord), \esp{respetar la voluntad del difunto} (respecter la volonté du défunt).
\end{aretenir}

\begin{propriete}[Structures d'opinion et de proposition]
\begin{itemize}
\item Opinion + indicatif : \esp{Creo que la solución es justa.} « Je pense que la solution est juste. »
\item Opinion négative + subjonctif : \esp{No creo que vender sea buena idea.} « Je ne pense pas que vendre soit une bonne idée. »
\item Proposition + subjonctif : \esp{Propongo que Pablo reciba el dinero de la cuenta.} « Je propose que Pablo reçoive l'argent du compte. » ; \esp{Es mejor que no vendamos la plantación.} « Il vaut mieux que nous ne vendions pas la plantation. »
\item Accord et désaccord : \esp{Estoy de acuerdo. / No estoy de acuerdo. / Tienes razón. / Perdóname, hermano.}
\end{itemize}
\end{propriete}

\begin{attention}[Pièges à éviter]
Ne dites pas \esp{*propongo que Pablo recibe} : après \esp{propongo que}, le subjonctif est obligatoire (\esp{reciba}). Attention au faux ami : \esp{el luto} signifie « le deuil », pas « la lutte » (qui se dit \esp{la lucha}). Enfin, \esp{huérfano} prend un \emph{h} muet et un accent sur le \emph{é}.
\end{attention}

\courssub{Proposition de production et corrigé commenté}

\begin{exoresolu}[Script modèle : « La reunión familiar »]
Énoncé : rédiger et jouer un dialogue de 10 à 12 répliques intégrant 6 mots du lexique, une proposition au subjonctif et une réconciliation finale.
\tcblower
\textbf{Script modèle (12 répliques) :}

\esp{Notario : Buenos días a todos. Han pasado tres meses de luto y hoy vamos a leer el testamento de don Ramón y a repartir la herencia.}

\esp{Viuda : Gracias, señor notario. Mi viudez es muy dolorosa, pero quiero que la familia permanezca unida.}

\esp{Andrés : Yo soy el hijo mayor y el principal heredero. Creo que debemos vender la plantación de Obala y dividir el dinero.}

\esp{María : ¡No estoy de acuerdo ! El testamento dice que nadie venda la plantación sin el acuerdo de los tres herederos.}

\esp{Pablo : Yo ahora soy huérfano de padre y necesito dinero para mis estudios en la universidad.}

\esp{Notario : Escuchen, por favor. El conflicto familiar no resuelve nada. ¿Quién tiene una propuesta ?}

\esp{María : Propongo que Pablo reciba el dinero de la cuenta bancaria para sus estudios.}

\esp{Viuda : Es justo que el más joven estudie. Y sugiero que guardemos la plantación : es el trabajo de toda la vida de su padre.}

\esp{Andrés : Pero la plantación necesita inversiones...}

\esp{María : Es mejor que trabajemos juntos la plantación y que repartamos las ganancias cada año.}

\esp{Andrés : Está bien, hermana. Tienes razón. Perdóname : la tristeza y el luto me hicieron pensar solo en el dinero.}

\esp{Viuda : Hijos míos, su padre estaría orgulloso. La familia es la verdadera herencia. ¡Vengan a abrazarme !}

\textbf{Commentaire des choix de langue.} Le script contient \cle{8 mots du lexique} de la leçon : \esp{luto, testamento, herencia, viudez, heredero, huérfano, conflicto familiar, repartir}. Il respecte la règle du \cle{subjonctif} après les verbes de proposition et de souhait : \esp{quiero que la familia permanezca unida}, \esp{propongo que Pablo reciba}, \esp{sugiero que guardemos}, \esp{es justo que el más joven estudie}, \esp{es mejor que trabajemos}. La réplique de María cite le testament avec le subjonctif \esp{que nadie venda}, ce qui montre la compréhension du document. La scène se termine par une \cle{réconciliation} explicite (\esp{Perdóname, hermana... Vengan a abrazarme}) et par une phrase de clôture mémorable qui reprend le thème de la leçon : \esp{La familia es la verdadera herencia.} « La famille est le véritable héritage. »
\end{exoresolu}

\courssub{Bilan de l'activité}

Cette activité a permis de vérifier l'acquisition des trois savoir-faire de la leçon. Sur le plan de la \cle{compréhension}, les élèves ont su lire un document juridique simple et en dégager les informations essentielles (les biens, les héritiers, les conditions). Sur le plan du \cle{lexique}, les mots du deuil et de l'héritage ont été réemployés dans un contexte oral significatif, ce qui favorise leur mémorisation durable. Sur le plan de la \cle{grammaire}, le subjonctif après les verbes de proposition et de souhait a été pratiqué de manière fonctionnelle, et non mécanique. Enfin, le jeu de rôle a développé l'\cle{expression orale} : prise de parole en public, intonation, gestuelle, écoute des partenaires, et le vote final a entraîné chaque élève à justifier un jugement en espagnol. Le critère de réussite retenu : un script complet, au moins 6 mots du lexique correctement employés, un subjonctif correct et une réconciliation crédible.

\newpage

\coursec{Méthodes et exercices résolus}

\courssub{Savoir-faire 1 : Lire et comprendre un texte sur les difficultés familiales}

\begin{textebox}[La herencia de don Ramón (texte support)]
Cuando don Ramón falleció, su viuda, doña Carmen, y sus tres hijos se reunieron en la casa familiar de Yaoundé. El notario leyó el testamento: la casa era para la viuda, pero el dinero debía repartirse entre los hijos. El hijo mayor, que había cuidado a su padre durante su enfermedad, protestó: «Yo me quedé con papá mientras vosotros vivíais lejos». Sus hermanos se enfadaron y la reunión terminó en un gran conflicto familiar. Después de varios días de luto y de silencio, la madre reunió de nuevo a sus hijos y les pidió que dialogaran con calma. Al final, los hermanos se reconciliaron y aceptaron el reparto decidido por su padre.
\end{textebox}

\textbf{Traduction :} « Quand don Ramón est décédé, sa veuve, doña Carmen, et ses trois enfants se sont réunis dans la maison familiale de Yaoundé. Le notaire a lu le testament : la maison revenait à la veuve, mais l'argent devait être partagé entre les enfants. Le fils aîné, qui avait soigné son père pendant sa maladie, a protesté : "Moi, je suis resté avec papa pendant que vous viviez loin." Ses frères et sœurs se sont mis en colère et la réunion s'est terminée par un grand conflit familial. Après plusieurs jours de deuil et de silence, la mère a réuni de nouveau ses enfants et leur a demandé de dialoguer calmement. Finalement, les frères et sœurs se sont réconciliés et ont accepté le partage décidé par leur père. »

\textbf{Glossaire :} \esp{fallecer} = décéder ; \esp{el reparto} = le partage ; \esp{enfadarse} = se mettre en colère ; \esp{reunir} = réunir ; \esp{reconciliarse} = se réconcilier ; \esp{al final} = finalement.

\begin{methode}[Comprendre un texte narratif sur la famille]
Pour réussir les questions de compréhension au Baccalauréat, suis cette démarche en cinq étapes :
\begin{enumerate}
\item \textbf{Lecture globale silencieuse} : lis le texte une première fois sans dictionnaire. Repère le thème général (ici : la muerte de un padre, el reparto de la herencia, los conflictos entre hermanos).
\item \textbf{Identifier la situation} : qui parle ? à qui ? où ? quand ? Note les personnages (\esp{la viuda, los hijos, el notario}) et leurs relations.
\item \textbf{Repérer les mots-clés du champ lexical} : souligne les termes du deuil et de l'héritage (\esp{viudez, luto, testamento, herederos, reparto}). Ils guident vers le sens même si tu ne comprends pas tout.
\item \textbf{Localiser la réponse} : pour chaque question, retourne au texte et encadre la phrase qui contient la réponse. Ne réponds jamais « de mémoire ».
\item \textbf{Rédiger en phrases complètes} : réponds en espagnol, avec des phrases entières, en reformulant (ne recopie pas le texte mot à mot, sauf si la consigne demande une citation entre comillas).
\end{enumerate}
\textbf{Réflexes à retenir :} une question avec \esp{¿Por qué...?} exige une réponse avec \esp{porque...} ; une question avec \esp{¿Qué...?} exige un complément d'objet ; attention à accorder les temps (si la question est au passé, réponds au passé).
\end{methode}

\begin{exoresolu}[Compréhension de texte : « La herencia de don Ramón »]
\textbf{Questions :} 1. \esp{¿Quiénes se reunieron después de la muerte de don Ramón?} 2. \esp{¿Qué recibió la viuda según el testamento?} 3. \esp{¿Por qué protestó el hijo mayor?}
\tcblower
\textbf{Raisonnement :} la question 1 porte sur les personnages : on la localise dans la première phrase. La question 2 porte sur le contenu du testamento (deuxième phrase). La question 3 commence par \esp{¿Por qué?}, donc la réponse doit commencer par \esp{porque} et justifier la protestation (troisième phrase).

\textbf{Réponses rédigées :}
\begin{enumerate}
\item \esp{Después de la muerte de don Ramón, se reunieron su viuda, doña Carmen, y sus tres hijos, en presencia del notario.}
\item \esp{Según el testamento, la viuda recibió la casa familiar, mientras que el dinero debía repartirse entre los tres hijos.}
\item \esp{El hijo mayor protestó porque él había cuidado a su padre durante su enfermedad, mientras sus hermanos vivían lejos.}
\end{enumerate}
\textbf{Justifications grammaticales :} on emploie le plus-que-parfait \esp{había cuidado} car l'action de soigner est antérieure à la protestation (\esp{protestó}, passé simple) ; \esp{mientras} introduit ici une simultanéité au passé, donc imparfait \esp{vivían}. Note l'accord : \esp{su viuda}, \esp{sus tres hijos}.
\end{exoresolu}

\courssub{Savoir-faire 2 : Utiliser le lexique du deuil et de l'héritage}

\begin{definition}[Lexique thématique : el duelo y la herencia]
\begin{center}
\begin{tabularx}{\linewidth}{|l|X|X|}\hline
\rowcolor{popblueL} \textbf{Español} & \textbf{Français} & \textbf{Exemple} \\ \hline
\esp{la viudez} & le veuvage & \esp{La viudez cambió su vida.} \\ \hline
\esp{el viudo / la viuda} & le veuf / la veuve & \esp{La viuda recibió la casa.} \\ \hline
\esp{el luto} & le deuil & \esp{La familia guardó luto un año.} \\ \hline
\esp{el huérfano / la huérfana} & l'orphelin / l'orpheline & \esp{Quedó huérfana a los cinco años.} \\ \hline
\esp{la herencia} & l'héritage & \esp{Repartieron la herencia.} \\ \hline
\esp{el heredero / la heredera} & l'héritier / l'héritière & \esp{Los hijos son los herederos.} \\ \hline
\esp{el testamento} & le testament & \esp{El notario leyó el testamento.} \\ \hline
\esp{el conflicto familiar} & le conflit familial & \esp{El reparto causó un conflicto familiar.} \\ \hline
\esp{heredar (de)} & hériter (de) & \esp{Heredó la casa de su padre.} \\ \hline
\esp{fallecer} & décéder & \esp{Don Ramón falleció en mayo.} \\ \hline
\esp{repartir entre} & partager entre & \esp{Repartió el dinero entre los hijos.} \\ \hline
\esp{reconciliarse} & se réconcilier & \esp{Los hermanos se reconciliaron.} \\ \hline
\end{tabularx}
\end{center}
\end{definition}

\begin{methode}[Employer correctement le vocabulaire thématique]
\begin{enumerate}
\item \textbf{Apprendre les mots par familles} : \esp{heredar} (verbe) $\rightarrow$ \esp{herencia} (la chose héritée) $\rightarrow$ \esp{heredero / heredera} (la personne). De même : \esp{enviudar} $\rightarrow$ \esp{viudez} $\rightarrow$ \esp{viudo / viuda}.
\item \textbf{Mémoriser les constructions avec préposition} : \esp{morir de} (mourir de), \esp{estar de luto por} (être en deuil de), \esp{repartir algo entre} (partager entre), \esp{dejar algo a alguien en herencia} (léguer).
\item \textbf{Distinguer les genres} : \esp{el testamento, el luto, el conflicto, el huérfano} ; \esp{la viudez, la herencia, la viuda}.
\item \textbf{Utiliser les mots en contexte} : pour chaque mot, fabrique une phrase personnelle liée au texte. C'est le meilleur moyen de les réutiliser à l'écrit.
\item \textbf{Vérifier les accents} : \esp{huérfano} (accent sur la voyelle ouverte), \esp{falleció}, \esp{reunión}, \esp{atención}. Un accent oublié change parfois le temps du verbe (\esp{repartio} n'existe pas : \esp{repartió}).
\end{enumerate}
\end{methode}

\begin{exoresolu}[Lexique : compléter et réutiliser]
\textbf{Énoncé :} Complète les phrases avec le mot qui convient : \esp{viudez, herederos, testamento, luto, huérfana, herencia}.
\begin{enumerate}
\item \esp{Después de la muerte de su marido, doña Carmen llevó el \ldots\ durante un año.}
\item \esp{El notario abrió el \ldots\ delante de toda la familia.}
\item \esp{Los tres hijos son los únicos \ldots\ de don Ramón.}
\item \esp{La pequeña Lucía quedó \ldots\ a los cinco años.}
\end{enumerate}
\tcblower
\textbf{Raisonnement :} on identifie d'abord le sens attendu dans chaque phrase, puis on vérifie le genre et le nombre du nom choisi.

\textbf{Solutions :}
\begin{enumerate}
\item \esp{luto} : \esp{llevar el luto} = porter le deuil. On ne dit pas \esp{llevar la viudez} : la \esp{viudez} est l'état de veuvage, pas la tenue de deuil.
\item \esp{testamento} : c'est le document que le notaire lit après un décès. Masculin singulier : \esp{el testamento}.
\item \esp{herederos} : masculin pluriel car il s'agit des trois fils ; \esp{heredero} = celui qui reçoit l'héritage.
\item \esp{huérfana} : féminin singulier (Lucía) ; \esp{quedar huérfano/a} = devenir orphelin. Attention à l'accent : \esp{huérfana}.
\end{enumerate}
\textbf{Phrase personnelle de réutilisation :} \esp{Tras la muerte de su padre, la familia guardó luto y el notario repartió la herencia entre los herederos.} « Après la mort de son père, la famille observa le deuil et le notaire partagea l'héritage entre les héritiers. »
\end{exoresolu}

\begin{savaistu}[La famille en Guinée équatoriale]
La Guinée équatoriale, seul pays hispanophone d'Afrique subsaharienne, partage avec le Cameroun de fortes traditions familiales : la famille élargie (\esp{la familia extensa}) joue un rôle central lors des deuils et des partages d'héritage. En espagnol, on dit \esp{guardar luto} ou \esp{estar de luto} pour « observer le deuil », et les réunions familiales après un décès (\esp{la reunión familiar tras el fallecimiento}) sont un moment important de la vie sociale, comme dans le texte de don Ramón.
\end{savaistu}

\courssub{Savoir-faire 3 : Commenter un texte (commentario de texto)}

\begin{methode}[Rédiger un commentaire de texte en quatre étapes]
\begin{enumerate}
\item \textbf{Presentación} : présente le texte (type, thème, personnages, situation). Formule utile : \esp{El texto trata de... / Se sitúa en... / Los personajes son...}
\item \textbf{Resumen} : résume le contenu en trois ou quatre phrases, sans recopier. Utilise des connecteurs : \esp{primero, luego, después, finalmente}.
\item \textbf{Análisis} : analyse les idées principales et le lexique thématique. Relève le champ lexical du deuil et de l'héritage et explique son rôle : \esp{El autor utiliza palabras como... para expresar...}
\item \textbf{Opinión personal} : donne ton avis argumenté : \esp{En mi opinión... / Me parece que... / Estoy de acuerdo con... porque...} Termine par une conclusion brève.
\end{enumerate}
\textbf{Connecteurs à retenir :} \esp{sin embargo} (cependant), \esp{además} (de plus), \esp{por eso} (c'est pourquoi), \esp{en conclusión} (en conclusion).
\end{methode}

\begin{exoresolu}[Commentaire guidé]
\textbf{Énoncé :} Rédige la présentation et l'opinion personnelle d'un commentaire du texte « La herencia de don Ramón » (4 à 6 phrases au total).
\tcblower
\textbf{Raisonnement :} la présentation doit contenir le thème (les problèmes familiaux après un décès), le lieu (Yaoundé) et les personnages. L'opinion doit être argumentée avec \esp{porque}.

\textbf{Production modèle :}

\esp{El texto trata de los problemas familiares que aparecen después de la muerte de un padre. La escena se sitúa en Yaoundé, en la casa familiar, donde la viuda y los tres hijos escuchan el testamento leído por el notario. El conflicto nace del reparto de la herencia: el hijo mayor cree merecer más porque cuidó a su padre enfermo.}

\esp{En mi opinión, el dinero no debe destruir a una familia. Me parece que los hermanos deberían dialogar con calma en lugar de enfadarse, porque la unión familiar vale más que cualquier herencia.}

« Le texte traite des problèmes familiaux qui apparaissent après la mort d'un père. La scène se déroule à Yaoundé... À mon avis, l'argent ne doit pas détruire une famille... »

\textbf{Points grammaticaux :} \esp{deberían dialogar} : conditionnel pour exprimer un conseil ; \esp{en lugar de + infinitivo} (au lieu de) ; \esp{cuidó} : passé simple avec accent obligatoire.
\end{exoresolu}

\courssub{Savoir-faire 4 : Traduire (thème et version) sur le thème de la famille}

\begin{methode}[Traduire sans calquer le français]
\begin{enumerate}
\item \textbf{Analyser la phrase française} : repère le sujet, le verbe, le temps logique, les compléments.
\item \textbf{Choisir le bon temps espagnol} : action passée ponctuelle $\rightarrow$ \esp{pretérito indefinido} ; habitude ou description au passé $\rightarrow$ \esp{imperfecto} ; antériorité $\rightarrow$ \esp{pluscuamperfecto}.
\item \textbf{Attention aux constructions différentes} : « manquer à quelqu'un » = \esp{echar de menos a alguien} ; « hériter de » = \esp{heredar (de)} ; « se disputer » = \esp{pelearse / discutir}.
\item \textbf{Vérifier les accords} : sujet-verbe, nom-adjectif (\esp{la casa familiar}, \esp{los hijos mayores}).
\item \textbf{Relire à voix haute} : chaque accent, chaque article, chaque préposition.
\end{enumerate}
\end{methode}

\begin{exoresolu}[Thème : trois phrases]
\textbf{Énoncé :} Traduis en espagnol :
\begin{enumerate}
\item Après la mort de son mari, doña Carmen est restée seule avec ses enfants.
\item Le notaire a lu le testament devant toute la famille.
\item Les frères se disputaient souvent à cause de l'héritage.
\end{enumerate}
\tcblower
\textbf{Raisonnement et solutions :}
\begin{enumerate}
\item Action passée ponctuelle et achevée $\rightarrow$ passé simple : \esp{Después de la muerte de su marido, doña Carmen se quedó sola con sus hijos.} « Rester » au sens de « demeurer, se retrouver » = \esp{quedarse}, pas \esp{restar} (faux ami : \esp{restar} = soustraire).
\item \esp{El notario leyó el testamento delante de toda la familia.} \esp{leyó} : passé simple de \esp{leer}, avec accent (3\textsuperscript{e} personne du singulier) ; \esp{delante de} = devant (ne pas confondre avec \esp{antes de} = avant, dans le temps).
\item Habitude répétée au passé $\rightarrow$ imparfait : \esp{Los hermanos se peleaban a menudo por la herencia.} \esp{a menudo} = souvent ; \esp{por} exprime ici la cause du conflit ; verbe pronominal \esp{pelearse} avec le pronom \esp{se}.
\end{enumerate}
\textbf{Conclusion :} le choix entre passé simple et imparfait est la clé de ces trois phrases : actions uniques (1, 2) contre habitude (3).
\end{exoresolu}

\courssub{Savoir-faire 5 : Jouer une scène de conversation familiale}

\begin{methode}[Préparer et jouer un jeu de rôle]
\begin{enumerate}
\item \textbf{Répartir les rôles} : la viuda, el hijo mayor, la hija, el hijo menor, el notario. Chacun définit son point de vue (pour ou contre le partage).
\item \textbf{Préparer des formules de conversation} :
\begin{itemize}
\item Donner son avis : \esp{Yo creo que... / A mi parecer...}
\item Être d'accord : \esp{Tienes razón. / Estoy de acuerdo contigo.}
\item Ne pas être d'accord : \esp{No estoy de acuerdo. / No me parece justo.}
\item Proposer une solution : \esp{¿Por qué no repartimos el dinero en partes iguales?}
\item Réconcilier : \esp{Perdóname, hermano. Lo importante es nuestra familia.}
\end{itemize}
\item \textbf{Respecter les registres} : tutoiement entre frères et sœurs, ton respectueux avec la mère (\esp{mamá, por favor, escúchame}).
\item \textbf{Jouer la scène} : parler clairement, regarder son partenaire, utiliser des gestes, terminer par une réconciliation.
\end{enumerate}
\end{methode}

\begin{exoresolu}[Jeu de rôle : compléter un dialogue]
\textbf{Énoncé :} Complète ce dialogue avec des formules adaptées.

\esp{— Hijo mayor: Yo cuidé a papá durante dos años. No me parece \ldots\ que el dinero se reparta en partes iguales.}

\esp{— Hija: Entiendo tu dolor, hermano, pero no estoy \ldots\ contigo: el testamento es claro.}

\esp{— Madre: Hijos, por favor. ¿Por qué no \ldots\ con calma? La familia es lo más importante.}

\esp{— Hijo mayor: Tienes \ldots, mamá. Perdonadme, hermanos.}
\tcblower
\textbf{Solutions :}
\begin{enumerate}
\item \esp{justo} : \esp{No me parece justo} = « cela ne me semble pas juste ». Formule de désaccord poli. Remarque : après \esp{no me parece justo}, le subjonctif \esp{se reparta} est correct car on exprime un jugement négatif.
\item \esp{de acuerdo} : \esp{estar de acuerdo con alguien} ; invariable, on ne dit jamais \esp{de acuerda}.
\item \esp{hablamos / dialogamos} : après \esp{¿Por qué no...?} on utilise le présent de l'indicatif pour faire une proposition : « Pourquoi ne parlons-nous pas calmement ? »
\item \esp{razón} : \esp{tener razón} = avoir raison ; expression figée, féminin singulier.
\end{enumerate}
\textbf{Conclusion :} ce dialogue montre la progression attendue d'un jeu de rôle réussi : conflit $\rightarrow$ discussion $\rightarrow$ proposition de la mère $\rightarrow$ réconciliation (\esp{Perdonadme, hermanos}).
\end{exoresolu}

\begin{experience}[Jeu de rôle en classe : el reparto de la herencia]
Par groupes de quatre ou cinq, jouez la scène du partage de l'héritage de don Ramón :
\begin{enumerate}
\item \textbf{Distribución de papeles} : la viuda, el hijo mayor, la hija, el hijo menor, el notario.
\item \textbf{Preparación} (5 minutes) : chaque élève écrit trois répliques avec les formules du cours (avis, désaccord, proposition).
\item \textbf{Representación} : jouez la scène devant la classe ; elle doit commencer par le conflit et se terminer par une réconciliation.
\item \textbf{Evaluación} : la classe note la prononciation, l'exactitude des formules et la logique de la scène.
\end{enumerate}
Variante : imaginez une fin différente (par exemple, le notario propose une médiation).
\end{experience}

\begin{attention}[Les 5 erreurs qui coûtent des points au Bac]
\begin{enumerate}
\item \textbf{Oublier les accents} : \esp{reunion} au lieu de \esp{reunión}, \esp{leo} au lieu de \esp{leyó}, \esp{huerfano} au lieu de \esp{huérfano}. Un accent oublié peut changer le temps ou le sens du mot.
\item \textbf{Calquer le français} : « être d'accord » n'est pas \esp{ser de acuerdo} mais \esp{estar de acuerdo} ; « avoir raison » = \esp{tener razón} (jamais \esp{haber razón}) ; « manquer à » = \esp{echar de menos}, pas \esp{manchar}.
\item \textbf{Faux amis} : \esp{restar} ne signifie pas « rester » (c'est \esp{quedarse / permanecer}) ; \esp{asistir a} signifie « assister à », pas « aider » (\esp{ayudar}) ; \esp{una herida} est une blessure, pas un héritage (\esp{una herencia}).
\item \textbf{Fautes d'accord} : \esp{la viuda y los hijos está tristes} est doublement faux : \esp{la viuda y los hijos están tristes} (sujet pluriel, verbe pluriel, adjectif pluriel).
\item \textbf{Mauvais choix de temps au passé} : mélanger passé simple et imparfait sans logique. Règle d'or : action unique et achevée $\rightarrow$ \esp{falleció, leyó, protestó} ; description, habitude, arrière-plan $\rightarrow$ \esp{vivían, cuidaba, se peleaban}.
\end{enumerate}
\end{attention}

\begin{exercice}[Pour t'entraîner]
\begin{enumerate}
\item Réponds en phrases complètes : \esp{¿Qué debía repartirse entre los hijos según el testamento? ¿Cómo termina la historia?}
\item Traduis en espagnol : « La veuve a gardé le deuil pendant un an. » ; « Les héritiers se sont réconciliés après la réunion. »
\item Rédige un court paragraphe d'opinion (3 phrases) : \esp{¿Es importante el diálogo en la familia? ¿Por qué?}
\end{enumerate}
\end{exercice}

\begin{corrige}[Exercice : Pour t'entraîner]
\begin{enumerate}
\item \esp{Según el testamento, el dinero debía repartirse entre los hijos, mientras que la casa era para la viuda. La historia termina bien: los hermanos se reconcilian y aceptan el reparto decidido por su padre.}
\item \esp{La viuda guardó luto durante un año.} (action achevée $\rightarrow$ passé simple \esp{guardó}) ; \esp{Los herederos se reconciliaron después de la reunión.} (verbe pronominal au passé simple, accent sur \esp{reunión}).
\item Modèle : \esp{En mi opinión, el diálogo es muy importante en la familia porque evita los conflictos. Me parece que los hermanos deben hablar con calma antes de enfadarse. Además, la unión familiar vale más que el dinero.}
\end{enumerate}
\end{corrige}

\newpage

\coursec{Exercices}

\courssub{Je vérifie mes connaissances}

\begin{exercice}[QCM : le lexique du deuil et de l'héritage]
Pour chaque question, entoure la bonne réponse.

\begin{enumerate}
\item La \esp{viudez} désigne :
\begin{enumerate}
\item la situation d'une personne qui a perdu son conjoint
\item la cérémonie d'enterrement
\item le partage des biens entre frères
\end{enumerate}
\item Un \esp{testamento} est :
\begin{enumerate}
\item un témoin du mariage
\item un document qui exprime les dernières volontés d'une personne
\item une maison familiale
\end{enumerate}
\item \esp{Huérfano} signifie :
\begin{enumerate}
\item qui a perdu son ou ses parents
\item qui n'a pas d'enfants
\item qui est le fils aîné
\end{enumerate}
\item Le \esp{luto} correspond à :
\begin{enumerate}
\item la fête organisée après les funérailles
\item la période de deuil et ses signes extérieurs
\item le notaire chargé de la succession
\end{enumerate}
\item Un \esp{conflicto familiar} est :
\begin{enumerate}
\item une réunion de famille
\item une dispute au sein de la famille
\item un repas de fête
\end{enumerate}
\end{enumerate}
\end{exercice}

\begin{exercice}[Vrai ou faux, justifié]
Réponds par \esp{Verdadero} ou \esp{Falso} et justifie ta réponse par une phrase complète en espagnol.

\begin{enumerate}
\item \esp{La herencia} désigne uniquement l'argent liquide laissé par un défunt.
\item Une \esp{viuda} est une femme dont le mari est mort.
\item Le \esp{reparto} est le partage des biens entre les héritiers.
\item Un \esp{heredero} est toujours le fils aîné de la famille.
\end{enumerate}
\end{exercice}

\begin{exercice}[Vocabulaire : trouve le mot]
Complète chaque définition par le mot espagnol exact, choisi dans la liste : \esp{heredero, testamento, luto, viudez, huérfano, reconciliación}.

\begin{enumerate}
\item Document écrit dans lequel une personne indique comment partager ses biens après sa mort : \ldots
\item Personne qui reçoit des biens d'un défunt : \ldots
\item État d'une personne dont le conjoint est décédé : \ldots
\item Enfant qui a perdu son père ou sa mère : \ldots
\item Action de faire la paix après une dispute : \ldots
\item Vêtements sombres et période de tristesse après un décès : \ldots
\end{enumerate}
\end{exercice}

\begin{exercice}[Conjugaison : le passé simple du récit]
Complète les phrases avec le verbe entre parenthèses conjugué au \esp{pretérito indefinido}.

\begin{enumerate}
\item Don Ramón \ldots\ (morir) el mes pasado, después de una larga enfermedad.
\item Sus hijos \ldots\ (leer) el testamento con el notario.
\item La viuda \ldots\ (quedarse) sola en la casa familiar.
\item Los hermanos \ldots\ (disputarse) la tierra del cacao.
\item Finalmente, la familia \ldots\ (aceptar) el reparto propuesto.
\item Yo \ldots\ (proponer) que vendiéramos el coche de mi padre.
\end{enumerate}
\end{exercice}

\courssub{Je m'entraîne}

\begin{exercice}[Thème : traduis en espagnol]
Traduis les phrases suivantes en espagnol correct, en faisant attention aux temps et au vocabulaire de la leçon.

\begin{enumerate}
\item Après la mort de son mari, Carmen est devenue veuve.
\item Les enfants se sont disputé l'héritage de leur père.
\item Je propose que nous gardions la maison familiale. (subjonctif)
\item Le notaire a lu le testament devant toute la famille.
\item Il est important que les frères se réconcilient. (subjonctif)
\end{enumerate}
\end{exercice}

\begin{exercice}[Texte à trous : un récit de veuvage]
Complète le texte suivant avec les mots de la liste : \esp{viudez, luto, huérfanos, testamento, herederos, reparto}.

\begin{textebox}[Una historia de familia]
Cuando falleció don Mateo, toda la familia guardó \ldots\ durante varias semanas. Su esposa, doña Ángela, comenzó entonces una larga \ldots\ que duró quince años. Los tres hijos, todavía menores, quedaron \ldots\ de padre. El notario abrió el \ldots\ ante los \ldots\ reunidos en el salón de la casa. El \ldots\ de las tierras y de la casa provocó muchas discusiones, pero al final todos firmaron el acuerdo.
\end{textebox}
\end{exercice}

\begin{exercice}[Appariements : mots et définitions]
Relie chaque mot espagnol (colonne de gauche) à sa définition (colonne de droite). Attention, il y a deux définitions en trop.

\begin{center}
\begin{tabularx}{\linewidth}{|X|X|}
\hline
\rowcolor{popblueL} \textbf{Palabra} & \textbf{Definición} \\
\hline
1. la herencia & a. persona que recibe los bienes de un difunto \\
\hline
2. el heredero & b. conjunto de bienes que se transmiten tras una muerte \\
\hline
3. la viuda & c. reunión de amigos en un café \\
\hline
4. el luto & d. mujer cuyo esposo ha fallecido \\
\hline
5. la reconciliación & e. tristeza y signos exteriores tras una muerte \\
\hline
 & f. fin de un conflicto y vuelta a la armonía \\
\hline
 & g. fiesta de cumpleaños familiar \\
\hline
\end{tabularx}
\end{center}
\end{exercice}

\begin{exercice}[Transformations : du style direct au style indirect]
Mets les phrases suivantes au discours indirect, en commençant par \esp{La madre dijo que\ldots} ou \esp{El notario pidió que\ldots}. Fais attention à la concordance des temps et aux changements de pronoms.

\begin{enumerate}
\item « El testamento está en el despacho del notario. »
\item « Mis hijos se reconciliarán antes de fin de año. »
\item « Firmad el documento, por favor. » (ordre)
\item « Yo he vendido la casa de Yaoundé. »
\item « No os peleéis por el dinero. » (ordre négatif)
\end{enumerate}
\end{exercice}

\courssub{Je me prépare au Bac}

\begin{exercice}[Compréhension écrite type Bac]
Lis le texte suivant, puis réponds aux questions \textbf{en espagnol}, avec des phrases complètes.

\begin{textebox}[El reparto de la finca]
Cuando murió don Elías Nguema, sus cuatro hijos se reunieron en la casa familiar de Malabo para escuchar la lectura del testamento. El mayor, Francisco, quería vender la finca de cacao y repartir el dinero en cuatro partes iguales. Su hermana Lucía se opuso enseguida: para ella, la finca era la memoria de su padre y venderla sería traicionar su historia. Los dos hermanos menores, gemelos de veinte años, no sabían qué pensar y guardaban silencio.

La madre, doña Rosa, escuchaba las discusiones con el corazón roto. « Vuestro padre trabajó cuarenta años en esa tierra —dijo al fin—. No quiero que el dinero destruya lo que él construyó con sus manos. » Después de una larga noche de conversación, los hermanos llegaron a un acuerdo: Lucía cuidaría la finca y daría cada año a sus hermanos una parte de la cosecha. Francisco aceptó, aunque no del todo contento, y la familia cenó junta por primera vez desde el funeral. La reconciliación, pensó doña Rosa, era la mejor herencia que su marido podía dejarles.
\end{textebox}

\textbf{Glossaire :} \esp{finca} : exploitation agricole ; \esp{cosecha} : récolte ; \esp{gemelos} : jumeaux ; \esp{corazón roto} : cœur brisé.

\begin{enumerate}
\item \textbf{Comprensión.} Réponds par des phrases complètes.
\begin{enumerate}
\item ¿Dónde se reunieron los hijos de don Elías y por qué?
\item ¿Qué quería hacer Francisco con la finca?
\item ¿Por qué se opuso Lucía a la venta?
\item ¿Qué acuerdo encontraron finalmente los hermanos?
\end{enumerate}
\item \textbf{Verdadero o falso.} Justifica tu respuesta con una cita del texto.
\begin{enumerate}
\item Los gemelos participaron activamente en la discusión.
\item Doña Rosa estaba contenta con la pelea de sus hijos.
\item La familia terminó cenando unida.
\end{enumerate}
\item \textbf{Lengua.}
\begin{enumerate}
\item Busca en el texto dos palabras del campo léxico del duelo y dos del campo léxico de la herencia.
\item Conjuga los verbos \esp{oponerse} y \esp{construir} al pretérito indefinido, tercera persona del singular.
\item Explica con tus palabras el sentido de la última frase del texto.
\end{enumerate}
\end{enumerate}
\end{exercice}

\begin{exercice}[Thème et version type Bac]
\textbf{A. Thème.} Traduis en espagnol les phrases suivantes.

\begin{enumerate}
\item Mon grand-père est mort l'année dernière et toute la famille a porté le deuil.
\item Ses héritiers ont découvert qu'il n'avait pas laissé de testament.
\item Ma tante voudrait que nous vendions les terres du village. (subjonctif)
\item Depuis qu'elle est veuve, elle vit seule avec ses deux enfants orphelins.
\item Il faut que vous vous réconciliiez avant la lecture du testament. (subjonctif)
\end{enumerate}

\textbf{B. Version.} Traduis en français le passage suivant.

\begin{textebox}[Para traducir]
La viudez cambió la vida de doña Carmen. Después de la muerte de su esposo, tuvo que ocuparse sola de la casa y de los tres niños. Los parientes discutían sin cesar sobre la herencia, pero ella prefería la paz de su familia a todo el dinero del mundo. Un día, reunió a todos en el patio y les dijo: « El amor de una familia vale más que una casa. »
\end{textebox}
\end{exercice}

\begin{exercice}[Expression écrite type Bac]
\textbf{Sujet :} \esp{¿El dinero puede destruir una familia? Da tu opinión argumentada.}

Rédige un texte en espagnol de \textbf{15 à 20 lignes} (environ 150 à 200 mots). Tu devras :

\begin{enumerate}
\item présenter le sujet et annoncer ta position dans une courte introduction ;
\item développer au moins \textbf{deux arguments} illustrés d'exemples (tu peux t'inspirer de situations familiales autour de l'héritage, au Cameroun ou ailleurs) ;
\item utiliser le lexique de la leçon (\esp{herencia, testamento, conflicto, reconciliación, luto\ldots}) et au moins \textbf{deux verbes au subjonctif} après une expression d'opinion ou de sentiment ;
\item conclure en donnant ta réponse personnelle à la question.
\end{enumerate}

\textbf{Variante de production orale :} à deux, jouez la scène de la réconciliation entre un frère et une sœur qui se sont disputés à propos d'un héritage. Chaque élève prononce au moins cinq répliques ; terminez la scène par un accord.
\end{exercice}

\newpage

\coursec{Corrigés des exercices}

\courssub{Je vérifie mes connaissances}

\begin{corrige}[Exercice 1 — QCM : le lexique du deuil et de l'héritage]
\begin{enumerate}
\item \textbf{1.} La bonne réponse est \textbf{a} : \esp{la viudez} désigne la situation d'une personne qui a perdu son conjoint (veuvage).
\item \textbf{2.} La bonne réponse est \textbf{b} : un \esp{testamento} est un document juridique dans lequel une personne exprime ses dernières volontés quant à la destination de ses biens après sa mort.
\item \textbf{3.} La bonne réponse est \textbf{a} : \esp{huérfano} signifie « orphelin », qui a perdu son père, sa mère ou les deux.
\item \textbf{4.} La bonne réponse est \textbf{b} : le \esp{luto} correspond à la période de deuil et à ses manifestations extérieures (vêtements noirs, abstention de fêtes).
\item \textbf{5.} La bonne réponse est \textbf{b} : un \esp{conflicto familiar} est une dispute, un désaccord grave au sein de la famille.
\end{enumerate}
\end{corrige}

\begin{corrige}[Exercice 2 — Vrai ou faux, justifié]
\begin{enumerate}
\item \esp{Falso.} \esp{La herencia} comprende el conjunto de todos los bienes, derechos y obligaciones que una persona deja a su muerte, no solo el dinero en efectivo (pueden ser inmuebles, joyas, deudas, etc.).
\item \esp{Verdadero.} Una \esp{viuda} es la mujer cuyo esposo ha fallecido. (El masculino es \esp{viudo}).
\item \esp{Verdadero.} El \esp{reparto} (o \esp{partición}) es la operación por la cual se dividen y adjudican los bienes hereditarios entre los herederos.
\item \esp{Falso.} Un \esp{heredero} es la persona llamada a suceder al difunto en sus bienes y derechos, ya sea por testamento (heredero testamentario) o por ley (heredero legítimo/ab intestato); no tiene por qué ser el hijo mayor.
\end{enumerate}
\end{corrige}

\begin{corrige}[Exercice 3 — Vocabulaire : trouve le mot]
\begin{enumerate}
\item \esp{testamento}
\item \esp{heredero}
\item \esp{viudez}
\item \esp{huérfano}
\item \esp{reconciliación}
\item \esp{luto}
\end{enumerate}
\end{corrige}

\begin{corrige}[Exercice 4 — Conjugaison : le passé simple du récit (\esp{pretérito indefinido})]
\begin{enumerate}
\item Don Ramón \textbf{murió} (morir, 3ª pers. sg., radical \emph{u}) el mes pasado, después de una larga enfermedad.
\item Sus hijos \textbf{leyeron} (leer, 3ª pers. pl., \emph{i} → \emph{y} en la terminação) el testamento con el notario.
\item La viuda \textbf{se quedó} (quedarse, reflexivo, regular -ar) sola en la casa familiar.
\item Los hermanos \textbf{se disputaron} (disputarse, reflexivo, regular -ar) la tierra del cacao.
\item Finalmente, la familia \textbf{aceptó} (aceptar, regular -ar) el reparto propuesto.
\item Yo \textbf{propuse} (proponer, 1ª pers. sg., radical \emph{u}) que vendiéramos el coche de mi padre.
\end{enumerate}
\end{corrige}

\begin{attention}[Verbes irréguliers au pretérito indefinido]
\emph{Morir} et \emph{proponer} (composé de \emph{poner}) ont un radical en \emph{u} : \esp{murió, propuse}. \emph{Leer} change le \emph{i} en \emph{y} aux 3èmes personnes : \esp{leyó, leyeron}. \emph{Quedarse} et \emph{disputarse} sont réguliers en \emph{-ar} mais s'utilisent au pronominal.
\end{attention}


\courssub{Je m'entraîne}

\begin{corrige}[Exercice 5 — Thème : traduis en espagnol]
\begin{enumerate}
\item \esp{Después de la muerte de su marido, Carmen se quedó viuda.} (ou \esp{quedó viuda} / \esp{se hizo viuda}).
\item \esp{Los hijos se disputaron la herencia de su padre.}
\item \esp{Propongo que guardemos la casa familiar.} (Subjonctif présent : \esp{guardemos}, 1ª pers. pl. de \emph{guardar}).
\item \esp{El notario leyó el testamento delante de toda la familia.}
\item \esp{Es importante que los hermanos se reconcilien.} (Subjonctif présent : \esp{reconcilien}, 3ª pers. pl. de \emph{reconciliar}).
\end{enumerate}
\end{corrige}

\begin{corrige}[Exercice 6 — Texte à trous : un récit de veuvage]
\begin{textebox}[Una historia de familia (complété)]
Cuando falleció don Mateo, toda la familia guardó \textbf{luto} durante varias semanas. Su esposa, doña Ángela, comenzó entonces una larga \textbf{viudez} que duró quince años. Los tres hijos, todavía menores, quedaron \textbf{huérfanos} de padre. El notario abrió el \textbf{testamento} ante los \textbf{herederos} reunidos en el salón de la casa. El \textbf{reparto} de las tierras y de la casa provocó muchas discusiones, pero al final todos firmaron el acuerdo.
\end{textebox}
\end{corrige}

\begin{corrige}[Exercice 7 — Appariements : mots et définitions]
\begin{center}
\begin{tabularx}{\linewidth}{|X|X|}
\hline
\rowcolor{popblueL} \textbf{Palabra} & \textbf{Definición correcta} \\
\hline
1. la herencia & b. conjunto de bienes que se transmiten tras una muerte \\
\hline
2. el heredero & a. persona que recibe los bienes de un difunto \\
\hline
3. la viuda & d. mujer cuyo esposo ha fallecido \\
\hline
4. el luto & e. tristeza y signos exteriores tras una muerte \\
\hline
5. la reconciliación & f. fin de un conflicto y vuelta a la armonía \\
\hline
\end{tabularx}
\end{center}
\emph{Les définitions c (réunion d'amis) et g (anniversaire) sont les intrus.}
\end{corrige}

\begin{corrige}[Exercice 8 — Transformations : du style direct au style indirect]
\begin{enumerate}
\item \esp{La madre dijo que el testamento estaba en el despacho del notario.} (Présent \rightarrow Imparfait de l'indicatif).
\item \esp{La madre dijo que sus hijos se reconciliarían antes de fin de año.} (Futur \rightarrow Conditionnel).
\item \esp{El notario pidió que firmaran el documento, por favor.} (Impératif affirmatif \rightarrow Subjonctif imparfait, concordance des temps : \emph{pidió} + \emph{firmaran}).
\item \esp{El notario dijo que él había vendido la casa de Yaoundé.} (Passé composé \rightarrow Plus-que-parfait de l'indicatif, verbe \emph{decir} affirmatif = indicatif).
\item \esp{La madre pidió que no se pelearan por el dinero.} (Impératif négatif \rightarrow Subjonctif imparfait : \emph{pidió que no} + \emph{se pelearan}).
\end{enumerate}
\end{corrige}

\begin{attention}[Concordance des temps au discours indirect]
Quand le verbe introducteur est au passé (\emph{dijo, pidió}), les temps du discours direct reculent : présent $\to$ imparfait, futur $\to$ conditionnel, passé composé $\to$ plus-que-parfait. L'impératif devient subjonctif (présent si le verbe introducteur est au présent, imparfait s'il est au passé). Pour \emph{pedir que}, on utilise toujours le subjonctif.
\end{attention}


\courssub{Je me prépare au Bac}

\begin{corrige}[Exercice 9 — Compréhension écrite type Bac : \esp{El reparto de la finca}]
\textbf{Barème indicatif :} 10 points au total (Compréhension 4 pts, V/F 3 pts, Langue 3 pts).

\textbf{1. Comprensión}
\begin{enumerate}
\item \esp{Los hijos de don Elías se reunieron en la casa familiar de Malabo para escuchar la lectura del testamento.}
\item \esp{Francisco quería vender la finca de cacao y repartir el dinero en cuatro partes iguales.}
\item \esp{Lucía se opuso porque la finca era la memoria de su padre y venderla sería traicionar su historia.}
\item \esp{Llegaron al acuerdo de que Lucía cuidaría la finca y daría cada año a sus hermanos una parte de la cosecha.}
\end{enumerate}

\textbf{2. Verdadero o falso}
\begin{enumerate}
\item \esp{Falso.} \esp{« Los dos hermanos menores, gemelos de veinte años, no sabían qué pensar y guardaban silencio. »} (Ils n'ont pas participé activement).
\item \esp{Falso.} \esp{« La madre, doña Rosa, escuchaba las discusiones con el corazón roto. »} / \esp{« No quiero que el dinero destruya lo que él construyó con sus manos. »} (Elle est peinée, pas contente).
\item \esp{Verdadero.} \esp{« La familia cenó junta por primera vez desde el funeral. »}
\end{enumerate}

\textbf{3. Lengua}
\begin{enumerate}
\item \textbf{Champ lexical du deuil :} \esp{funeral, corazón roto, luto} (sous-entendu par le contexte), \esp{viuda} (statut). \textbf{Champ lexical de l'héritage :} \esp{testamento, herencia, reparto, finca, cosecha, dinero, acuerdo}.
\item \esp{Oponerse} $\to$ \textbf{se opuso} (3ª pers. sg., pret. indef., radical \emph{u} pour les verbes en \emph{-ir} à radical changeant \emph{o} $\to$ \emph{u}). \esp{Construir} $\to$ \textbf{construyó} (3ª pers. sg., pret. indef., \emph{i} $\to$ \emph{y} entre voyelles).
\item \esp{La última frase significa que la reconciliación y la unión familiar son el legado más valioso que don Elías pudo dejar a sus hijos, más importante que los bienes materiales (la finca, el dinero). La verdadera herencia es la paz y el amor entre hermanos.}
\end{enumerate}
\end{corrige}

\begin{corrige}[Exercice 10 — Thème et version type Bac]
\textbf{A. Thème (Traduction vers l'espagnol)}
\begin{enumerate}
\item \esp{Mi abuelo murió el año pasado y toda la familia guardó luto.} (ou \esp{llevó luto}).
\item \esp{Sus herederos descubrieron que no había dejado testamento.} (Plus-que-parfait de l'indicatif \esp{había dejado} car \emph{descubrir que} + fait réel ; subjonctif \esp{dejara} possible mais moins courant pour un fait avéré).
\item \esp{Mi tía quiere que vendamos las tierras del pueblo.} (Subjonctif présent \esp{vendamos}, 1ª pers. pl., après verbe de volonté \emph{querer que}).
\item \esp{Desde que es viuda / se quedó viuda, vive sola con sus dos hijos huérfanos.}
\item \esp{Es necesario que se reconcilien antes de la lectura del testamento.} (Subjonctif présent \esp{reconcilien}, 3ª pers. pl., forme \emph{ustedes} standard ; forme \emph{vosotros} : \esp{os reconciliéis}).
\end{enumerate}

\textbf{B. Version (Traduction vers le français)}
\begin{textebox}[Traduction proposée]
La veuvage a changé la vie de doña Carmen. Après la mort de son mari, elle a dû s'occuper seule de la maison et des trois enfants. Les parents se disputaient sans cesse au sujet de l'héritage, mais elle préférait la paix de sa famille à tout l'argent du monde. Un jour, elle a réuni tout le monde dans la cour et leur a dit : « L'amour d'une famille vaut plus qu'une maison. »
\end{textebox}
\end{corrige}

\begin{attention}[Points de vigilance]
\begin{itemize}
\item \esp{Viudez} = veuvage (nom abstrait), pas \emph{veuve} (adjectif/nom \esp{viuda}).
\item \esp{Ocuparse de} = s'occuper de (verbe pronominal).
\item \esp{Parientes} = parents (au sens large : famille, proches), pas \emph{pères/mères} (\esp{padres}).
\item \esp{Discutían} = imparfait (action habituelle, description passée).
\item \esp{Prefería} = imparfait (sentiment, état d'esprit passé).
\item \esp{Reunió} = pretérito indefinido (action ponctuelle, achevée).
\item \esp{Vale más que} = vaut plus que (comparatif de supériorité).
\end{itemize}
\end{attention}


\begin{corrige}[Exercice 11 — Expression écrite type Bac]
\textbf{Barème indicatif (20 points) :}
\begin{itemize}
\item Introduction (présentation + thèse) : 3 pts
\item Argument 1 + exemple + vocabulaire/grammaire : 5 pts
\item Argument 2 + exemple + vocabulaire/grammaire : 5 pts
\item Conclusion (réponse personnelle) : 3 pts
\item Correction linguistique (orthographe, accents, syntaxe) : 2 pts
\item Richesse lexicale (champ familial/héritage) + 2 subjonctifs minimum : 2 pts
\end{itemize}

\textbf{Proposition de rédaction modèle (environ 170 mots) :}

\begin{textebox}[¿El dinero puede destruir una familia?]
El dinero, sin duda, puede convertirse en un veneno para los lazos familiares, aunque no tiene por qué serlo. Cuando muere un ser querido, el \esp{luto} debería unir a los supervivientes, pero a menudo la lectura del \esp{testamento} despierta codicias dormidas.

En primer lugar, la ausencia de un testamento claro genera \esp{conflictos} terribles. He visto en Yaundé familias enteras romperse por una casa o un terreno de cacao: los hermanos se acusan de robo, los tíos exigen su parte y los \esp{herederos} legítimos acaban ante los tribunales. \textbf{Es triste que la avaricia borre años de cariño.} (Subjonctif \esp{borre} après \emph{Es triste que}).

En segundo lugar, la falta de diálogo empeora la situación. Si los hermanos no hablan, el rencor crece. Sin embargo, la \esp{reconciliación} es posible si se antepone el afecto al interés. \textbf{Mi abuelo siempre dijo: "Quiero que os queráis más que a mis tierras".} (Subjonctif \esp{queráis} après \emph{Quero que}). Cuando mi padre murió, mis tíos y mi madre vendieron la finca y repartieron el dinero, pero guardaron la casa familiar para las reuniones. Esa decisión salvó a la familia.

En conclusión, el dinero solo destruye a la familia si los corazones están vacíos. La verdadera \esp{herencia} no es el banco, sino el amor y el respeto mutuo. Por tanto, mi respuesta es no: el dinero no tiene por qué destruirnos, si sabemos priorizar lo humano.
\end{textebox}

\textbf{Variante orale (scène de réconciliation) — Guide pour l'enseignant :}
\end{corrige}

\begin{experience}[Jeu de rôle : La réconciliation]
\textbf{Personnages :} \emph{Marcos} (frère aîné, veut vendre) et \emph{Ana} (sœur cadette, veut garder la maison).
\textbf{Contexte :} Après la lecture du testament de leur père, ils se retrouvent seuls dans le salon.
\textbf{Objectif :} Minimum 5 répliques chacun, finir par un accord.
\textbf{Exemple de début :}
\esp{Marcos: Ana, por favor, escúchame. La casa es grande, cara de mantener y ninguno de los dos vive aquí.}
\esp{Ana: Pero Marcos, esta casa es nuestra historia. Aquí crecimos. Venderla es borrar a papá.}
\esp{Marcos: No es borrarlo, es ser realistas. Yo necesito mi parte para mi negocio en Douala.}
\esp{Ana: Y yo necesito un lugar para mis hijos en vacaciones. ¿Qué tal si yo me quedo la casa y te pago tu parte en plazos?}
\esp{Marcos: ¿En plazos?... Vale. Pero con un contrato ante notario.}
\esp{Ana: Por supuesto. Gracias, hermano. Papá estaría orgulloso de nosotros.}
\esp{Marcos: Sí... Brindemos por la familia.}
\end{experience}


\newpage

\coursec{Fiche bilan}

\begin{popfigure}
\begin{center}
\begin{tikzpicture}[mindmap, grow cyclic,
  every node/.style={concept, align=center, font=\small},
  root concept/.append style={concept color=popblue, font=\bfseries, minimum size=3.2cm},
  level 1 concept/.append style={level distance=4.6cm, sibling angle=60, font=\footnotesize},
  level 2 concept/.append style={level distance=2.8cm, font=\scriptsize}]
\node[root concept, align=center]{La familia y sus problemas\\ (viudez y herencia)}
  child[concept color=poporange]{ node[align=center]{Texto\\ \esp{La herencia de don Ramón}}
    child { node[concept color=poporangeL, align=center]{conflicto\\ familiar} }
    child { node[concept color=poporangeL] {reconciliación} }
  }
  child[concept color=popgreen]{ node[align=center]{Léxico\\ del duelo}
    child { node[concept color=popgreenL, align=center]{viudez, luto,\\ huérfano} }
  }
  child[concept color=popgreen!80!popblue]{ node[align=center]{Léxico\\ de la herencia}
    child { node[concept color=popgreenL, align=center]{testamento,\\ heredero} }
  }
  child[concept color=poppurple]{ node {Gramática}
    child { node[concept color=poppurpleL, align=center]{pasados :\\ indefinido /\\ imperfecto} }
    child { node[concept color=poppurpleL, align=center]{subjuntivo\\ de deseo} }
  }
  child[concept color=poppink]{ node[align=center]{Método\\ comentario}
    child { node[concept color=poppinkL, align=center]{tema, idea,\\ estructura,\\ juicio} }
  }
  child[concept color=popgold]{ node {Juego de rol}
    child { node[concept color=popgoldL, align=center]{partir la\\ herencia,\\ reconciliarse} }
  };
\end{tikzpicture}
\end{center}
\legende{Carte mentale de la leçon : la famille et ses problèmes (veuvage et héritage)}
\end{popfigure}

\begin{bilan}
\textbf{1. Lexique clé à connaître par cœur}
\begin{center}
\begin{tabularx}{\linewidth}{|X|X|X|X|}
\hline
\rowcolor{popblueL} \textbf{Español} & \textbf{Français} & \textbf{Español} & \textbf{Français} \\
\hline
la viudez & le veuvage & el/la viudo/a & le veuf / la veuve \\
\hline
la herencia & l'héritage & el/la heredero/a & l'héritier / l'héritière \\
\hline
el testamento & le testament & heredar & hériter \\
\hline
el luto / el duelo & le deuil & el/la huérfano/a & l'orphelin / l'orpheline \\
\hline
fallecer / morir & décéder / mourir & la muerte & la mort \\
\hline
el conflicto familiar & le conflit familial & reconciliarse & se réconcilier \\
\hline
repartir / partir & partager & los bienes & les biens \\
\hline
la viuda hereda la casa & la veuve hérite de la maison & dejar en herencia & laisser en héritage \\
\hline
\end{tabularx}
\end{center}

\textbf{2. Grammaire à maîtriser}
\begin{center}
\begin{tabularx}{\linewidth}{|X|X|}
\hline
\rowcolor{popgreenL} \textbf{Règle} & \textbf{Exemple} \\
\hline
Récit au passé : \cle{indefinido} pour l'action ponctuelle, \cle{imperfecto} pour le décor et les habitudes. & \esp{Don Ramón falleció ayer; era un hombre generoso.} « Don Ramón est mort hier ; c'était un homme généreux. » \\
\hline
\cle{Subjonctif} après les verbes de souhait et de sentiment : \esp{esperar que, querer que, ojalá}. & \esp{Espero que nos reconciliemos.} « J'espère que nous nous réconcilierons. » \\
\hline
\cle{Ser} + possession et origine des biens ; \esp{ser de + nom}. & \esp{La casa es de la viuda.} « La maison est à la veuve. » \\
\hline
Pronoms d'objet direct avec \esp{dejar, repartir, dar}. & \esp{Se lo dejó a sus hijos.} « Il le leur a laissé. » \\
\hline
\end{tabularx}
\end{center}

\textbf{3. Méthode express : le commentaire de texte}
\begin{enumerate}
\item Lire deux fois ; repérer le thème (\esp{viudez, herencia, conflicto}).
\item Dégager l'idée principale en une phrase.
\item Analyser la structure (situation initiale, conflit, dénouement).
\item Étudier la langue : temps des verbes, lexique du deuil, ton du texte.
\item Conclure par un jugement personnel argumenté.
\end{enumerate}

\textbf{4. Méthode express : le jeu de rôle}
\begin{enumerate}
\item Choisir les rôles : \esp{la viuda, los hijos, el notario}.
\item Préparer 4 ou 5 répliques chacun avec le lexique de l'héritage.
\item Utiliser des formules polies : \esp{me parece que..., estoy de acuerdo, perdonen, propongo que...}.
\item Viser la réconciliation : terminer par un accord (\esp{llegar a un acuerdo}).
\end{enumerate}
\end{bilan}

\begin{aretenir}[Checklist avant le Bac]
\begin{itemize}
\item[$\square$] Je connais le lexique du deuil : \esp{viudez, luto, huérfano, fallecer}.
\item[$\square$] Je connais le lexique de l'héritage : \esp{herencia, heredero, testamento, bienes}.
\item[$\square$] Je distingue \esp{morir} (mourir) de \esp{matar} (tuer) et \esp{fallecer} (décéder, registre soutenu).
\item[$\square$] Je n'oublie pas l'accent : \esp{Ramón, corazón, reconciliación}.
\item[$\square$] Je choisis correctement entre indefinido et imperfecto dans un récit au passé.
\item[$\square$] Je mets le subjonctif après \esp{espero que}, \esp{quiero que}, \esp{ojalá}.
\item[$\square$] Je sais dire « hériter de » : \esp{heredar de} et non un calque du français.
\item[$\square$] Je sais structurer un commentaire : thème, idée principale, structure, langue, jugement.
\item[$\square$] Je sais exprimer mon opinion et proposer un accord : \esp{me parece que..., propongo que...}.
\item[$\square$] Je sais jouer une scène familiale avec des formules de politesse et de réconciliation.
\item[$\square$] Je vérifie la ponctuation espagnole : ¿ ? ¡ ! dans mes dialogues.
\end{itemize}
\end{aretenir}

\findecours

\end{document}
